\input{preamble}

% OK, start here.
%
\begin{document}

\title{Obsolète}

\maketitle

\phantomsection
\label{section-phantom}

\tableofcontents




\section{Introduction}
\label{section-introduction}

\noindent
Dans ce chapitre, nous rassemblons certains lemmes devenus «~obsolètes~»
(voir \cite{Miller}).


\section{Préliminaires}
\label{section-preliminaries}

\begin{remark}
\label{remark-composition-of-adjoints-isomorphic-to-identity}
Les informations qui figuraient auparavant dans cette remarque sont désormais
contenues dans la combinaison des lemmes \ref{categories-lemma-adjoint-fully-faithful}
et \ref{categories-lemma-left-adjoint-composed-fully-faithful} du chapitre
Catégories.
\end{remark}


\section{Algèbre homologique}
\label{section-homological-algebra}

\begin{remark}
\label{remark-weak-serre-subcategory}
Les remarques suivantes sont obsolètes, car leur contenu est désormais couvert
par les lemmes \ref{homology-lemma-biregular-ss-converges} et
\ref{homology-lemma-first-quadrant-ss} du chapitre Homologie.
Soit $\mathcal{A}$ une catégorie abélienne.
Soit $\mathcal{C} \subset \mathcal{A}$
une sous-catégorie de Serre faible (voir la définition
\ref{homology-definition-serre-subcategory} du chapitre Homologie).
Supposons que $K^{\bullet, \bullet}$ soit un bicomplexe auquel
le lemme \ref{homology-lemma-first-quadrant-ss} du chapitre Homologie
s'applique et que, pour un certain $r \geq 0$, tous les objets
${}'E_r^{p, q}$ appartiennent à $\mathcal{C}$. Alors tous les objets de cohomologie
$H^n(sK^\bullet)$ appartiennent à $\mathcal{C}$. En effet, les hypothèses entraînent
que les noyaux et les images des ${}'d_r^{p, q}$ appartiennent à $\mathcal{C}$.
On en déduit que chaque ${}'E_{r + 1}^{p, q}$ appartient à $\mathcal{C}$.
Par récurrence, chaque ${}'E_\infty^{p, q}$ appartient à $\mathcal{C}$.
Par conséquent, chaque $H^n(sK^\bullet)$ possède une filtration finie dont les
sous-quotients appartiennent à $\mathcal{C}$. Comme $\mathcal{C}$ est stable par
extensions, on conclut que $H^n(sK^\bullet)$ appartient à $\mathcal{C}$, comme annoncé.
Le même résultat vaut pour la seconde suite spectrale associée
à $K^{\bullet, \bullet}$. De même, si $(K^\bullet, F)$ est un complexe filtré auquel
le lemme \ref{homology-lemma-biregular-ss-converges} du chapitre Homologie
s'applique et si, pour un certain $r \geq 0$, tous les objets $E_r^{p, q}$
appartiennent à $\mathcal{C}$, alors chaque $H^n(K^\bullet)$ est
un objet de $\mathcal{C}$.
\end{remark}




\section{Lemmes d'algèbre obsolètes}
\label{section-algebra}

\begin{lemma}
\label{lemma-finite-presentation-module-independent}
Soit $M$ un $R$-module de présentation finie.
Pour toute surjection $\alpha : R^{\oplus n} \to M$, le
noyau de $\alpha$ est un $R$-module de type fini.
\end{lemma}

\begin{proof}
C'est un cas particulier du lemme \ref{algebra-lemma-extension} du chapitre Algèbre.
\end{proof}

\begin{lemma}
\label{lemma-p-ring-map}
Soit $\varphi : R \to S$ un homomorphisme d'anneaux. Si
\begin{enumerate}
\item pour tout $x \in S$, il existe $n > 0$ tel que
$x^n$ appartienne à l'image de $\varphi$, et
\item pour tout $x \in \Ker(\varphi)$, il existe $n > 0$
tel que $x^n = 0$,
\end{enumerate}
alors $\varphi$ induit un homéomorphisme sur les spectres. Soit
$p$ un nombre premier tel que
\begin{enumerate}
\item[(a)] $S$ soit engendrée comme $R$-algèbre par des éléments $x$ pour
lesquels il existe $n > 0$ tel que $x^{p^n} \in \varphi(R)$ et
$p^nx \in \varphi(R)$, et
\item[(b)] le noyau de $\varphi$ soit engendré par des éléments nilpotents.
\end{enumerate}
Alors (1) et (2) sont satisfaites et, pour tout homomorphisme d'anneaux $R \to R'$,
l'homomorphisme $R' \to R' \otimes_R S$ satisfait aussi (a), (b), (1) et (2) ;
en particulier, il induit un homéomorphisme sur les spectres.
\end{lemma}

\begin{proof}
Cela résulte de la combinaison des lemmes
\ref{algebra-lemma-powers} et \ref{algebra-lemma-p-ring-map} du chapitre Algèbre.
\end{proof}

\noindent
Le lemme technique suivant affirme que l'on peut relever à la situation globale,
en un sens convenable, toute suite de relations donnée sur une fibre d'un
homomorphisme d'anneaux essentiellement de type fini.

\begin{lemma}
\label{lemma-lift-elements-ideal}
Soit $R \to S$ un homomorphisme d'anneaux.
Soit $\mathfrak p \subset R$ un idéal premier.
Soit $\mathfrak q \subset S$ un idéal premier au-dessus de $\mathfrak p$.
Supposons que $S_{\mathfrak q}$ soit essentiellement de type fini sur $R_\mathfrak p$.
Supposons donnés
\begin{enumerate}
\item un entier $n \geq 0$,
\item un idéal premier $\mathfrak a \subset \kappa(\mathfrak p)[x_1, \ldots, x_n]$,
\item un homomorphisme surjectif de $\kappa(\mathfrak p)$-algèbres
$$
\psi : (\kappa(\mathfrak p)[x_1, \ldots, x_n])_{\mathfrak a}
\longrightarrow
S_{\mathfrak q}/\mathfrak p S_{\mathfrak q},
$$
et
\item des éléments $\overline{f}_1, \ldots, \overline{f}_e$ de $\Ker(\psi)$.
\end{enumerate}
Il existe alors
\begin{enumerate}
\item un entier $m \geq 0$,
\item un élément $g \in S$, $g \not\in \mathfrak q$,
\item un homomorphisme
$$
\Psi :
R[x_1, \ldots, x_n, x_{n + 1}, \ldots, x_{n + m}]
\longrightarrow
S_g,
$$
et
\item des éléments $f_1, \ldots, f_e, f_{e + 1}, \ldots, f_{e + m}$
de $\Ker(\Psi)$
\end{enumerate}
tels que
\begin{enumerate}
\item le diagramme suivant soit commutatif :
$$
\xymatrix{
R[x_1, \ldots, x_{n + m}] \ar[d]_\Psi
\ar[rr]_-{x_{n + j} \mapsto 0} & &
(\kappa(\mathfrak p)[x_1, \ldots, x_n])_{\mathfrak a} \ar[d]^\psi \\
S_g \ar[rr] & &
S_{\mathfrak q}/\mathfrak p S_{\mathfrak q}
},
$$
\item pour $i \leq e$, l'élément $f_i$ ait pour image un multiple de
$\overline{f}_i$ par une unité dans l'anneau local
$$
(\kappa(\mathfrak p)[x_1, \ldots, x_{n + m}])_{
(\mathfrak a, x_{n + 1}, \ldots, x_{n + m})},
$$
\item l'élément $f_{e + j}$ ait pour image un multiple de
$x_{n + j}$ par une unité dans le même anneau local, et
\item l'homomorphisme induit $R[x_1, \ldots, x_{n + m}]_{\mathfrak b}
\to S_{\mathfrak q}$ soit surjectif, où
$\mathfrak b = \Psi^{-1}(\mathfrak qS_g)$.
\end{enumerate}
\end{lemma}

\begin{proof}
Nous affirmons qu'il suffit de démontrer le lemme lorsque $R$
et $S$ sont locaux, d'idéaux maximaux respectifs $\mathfrak p$ et $\mathfrak q$.
Supposons en effet avoir construit
$$
\Psi' : R_{\mathfrak p}[x_1, \ldots, x_{n + m}]
\longrightarrow
S_{\mathfrak q}
$$
et $f_1', \ldots, f_{e + m}' \in R_{\mathfrak p}[x_1, \ldots, x_{n + m}]$
possédant toutes les propriétés requises. Il existe alors un élément
$f \in R$, $f \not \in \mathfrak p$, tel que chacun des
$ff_k'$ provienne d'un élément $f_k \in R[x_1, \ldots, x_{n + m}]$.
De plus, pour un $g \in S$, $g \not \in \mathfrak q$, convenable,
les éléments $\Psi'(x_i)$ sont les images d'éléments
$y_i \in S_g$. Soit $\Psi$ l'homomorphisme de $R$-algèbres défini
par la règle $\Psi(x_i) = y_i$. Puisque $\Psi(f_i)$ est nul
dans le localisé $S_{\mathfrak q}$, nous pouvons, quitte à remplacer
$g$, supposer que $\Psi(f_i) = 0$. L'affirmation est démontrée.

\medskip\noindent
Nous pouvons donc supposer $R$ et $S$ locaux,
d'idéaux maximaux $\mathfrak p$ et $\mathfrak q$.
Choisissons $y_1, \ldots, y_n \in S$ tels que
$y_i \bmod \mathfrak pS = \psi(x_i)$.
Soient $y_{n + 1}, \ldots, y_{n + m} \in S$ des éléments qui engendrent
une $R$-algèbre dont $S$ est une localisation.
Ils existent puisque $S$ est essentiellement de type fini sur $R$.
Comme $\psi$ est surjectif, on peut écrire
$y_{n + j} \bmod \mathfrak pS = \psi(h_j)$ pour
un certain $h_j \in \kappa(\mathfrak p)[x_1, \ldots, x_n]_{\mathfrak a}$.
Écrivons $h_j = g_j/d$, avec $g_j \in \kappa(\mathfrak p)[x_1, \ldots, x_n]$
et un dénominateur commun $d \in \kappa(\mathfrak p)[x_1, \ldots, x_n]$,
$d \not \in \mathfrak a$. Choisissons des relèvements
$G_j, D \in R[x_1, \ldots, x_n]$ de $g_j$ et $d$. Posons
$y_{n + j}' = D(y_1, \ldots, y_n) y_{n + j} - G_j(y_1, \ldots, y_n)$.
Par construction, $y_{n + j}' \in \mathfrak p S$.
Il est clair que $y_1, \ldots, y_n, y_{n + 1}', \ldots, y_{n + m}'$
engendrent une $R$-sous-algèbre de $S$ dont $S$ est une localisation.
On définit
$$
\Psi : R[x_1, \ldots, x_{n + m}] \to S
$$
comme l'homomorphisme qui envoie $x_i$ sur $y_i$ pour $i = 1, \ldots, n$
et $x_{n + j}$ sur $y'_{n + j}$ pour $j = 1, \ldots, m$. Les propriétés
(1) et (4) sont immédiates par construction. De plus, l'idéal
$\mathfrak b$ s'envoie sur l'idéal
$(\mathfrak a, x_{n + 1}, \ldots, x_{n + m})$
de l'anneau de polynômes $\kappa(\mathfrak p)[x_1, \ldots, x_{n + m}]$.

\medskip\noindent
Notons $J = \Ker(\Psi)$. On a une suite exacte courte
$$
0 \to J_{\mathfrak b}
\to R[x_1, \ldots, x_{n + m}]_{\mathfrak b}
\to S_{\mathfrak q}
\to 0.
$$
La surjectivité résulte du choix fait ci-dessus de
$y_1, \ldots, y_n, y_{n + 1}', \ldots, y_{n + m}'$.
Il en résulte que la suite
$$
J_{\mathfrak b}/ \mathfrak pJ_{\mathfrak b}
\to \kappa(\mathfrak p)[x_1, \ldots, x_{n + m}]_{
(\mathfrak a, x_{n + 1}, \ldots, x_{n + m})}
\to S_{\mathfrak q}/\mathfrak pS_{\mathfrak q}
\to 0
$$
est exacte. Par construction, $x_i$ s'envoie sur $\psi(x_i)$ et
$x_{n + j}$ sur zéro par le dernier homomorphisme.
On choisit alors aisément les $f_i$ comme en
(2) et (3) du lemme.
\end{proof}

\begin{lemma}
\label{lemma-smooth-independent-presentation}
Soit $R \to S$ un homomorphisme d'anneaux de présentation finie.
Si, pour une présentation $\alpha$ de $S$ sur $R$, le
complexe cotangent naïf $\NL(\alpha)$ est quasi-isomorphe
à un $S$-module projectif de type fini placé en degré $0$, il en va
de même pour toute présentation.
\end{lemma}

\begin{proof}
Cela résulte immédiatement du lemme \ref{algebra-lemma-NL-homotopy} du chapitre Algèbre.
\end{proof}

\begin{remark}[Résolutions projectives]
\label{remark-projective-resolution}
Soit $R$ un anneau.
Pour tout ensemble $S$, notons $F(S)$ le $R$-module libre de base $S$.
Tout $R$-module à gauche possède alors la résolution en deux étapes suivante :
$$
F(M \times M) \oplus F(R \times M) \to F(M) \to M \to 0.
$$
Le premier homomorphisme est donné par la règle
$$
[m_1, m_2] \oplus [r, m] \mapsto [m_1 + m_2] - [m_1] - [m_2] + [rm] - r[m].
$$
\end{remark}

\begin{lemma}
\label{lemma-spec-localization-first}
Soit $S$ une partie multiplicative de $A$. Alors l'homomorphisme
$$
f: \Spec(S^{-1}A)\longrightarrow \Spec(A)
$$
induit par l'homomorphisme canonique d'anneaux
$A \to S^{-1}A$ est un homéomorphisme sur son image et
$\Im(f) = \{ \mathfrak p \in \Spec(A) : \mathfrak p\cap S = \emptyset \}$.
\end{lemma}

\begin{proof}
C'est un doublon du lemme \ref{algebra-lemma-spec-localization} du chapitre Algèbre.
\end{proof}

\begin{lemma}
\label{lemma-finite-type-flat-over-integral-algebra}
Soit $A \to B$ un homomorphisme d'anneaux plat de type fini, où $A$ est
intègre. Alors $B$ est une $A$-algèbre de présentation finie.
\end{lemma}

\begin{proof}
C'est un cas particulier de la proposition
\ref{flat-proposition-flat-finite-type-finite-presentation-domain}
du chapitre Compléments sur la platitude.
\end{proof}

\begin{lemma}
\label{lemma-helper-finite-type-flat-finite-presentation}
Soit $R$ un anneau intègre de corps des fractions $K$.
Soit $S = R[x_1, \ldots, x_n]$ un anneau de polynômes sur $R$.
Soit $M$ un $S$-module de type fini. Supposons que $M$ soit plat sur $R$.
Si, pour tout sous-anneau $R \subset R' \subset K$, $R \not = R'$,
le module $M \otimes_R R'$ est de présentation finie
sur $S \otimes_R R'$, alors $M$ est de présentation finie sur $S$.
\end{lemma}

\begin{proof}
Ce lemme est vrai parce que $M$ est de présentation finie même sans
l'hypothèse selon laquelle $M \otimes_R R'$ est de présentation finie pour tout
$R'$ comme dans l'énoncé. Cela résulte de la proposition
\ref{flat-proposition-flat-finite-type-finite-presentation-domain}
du chapitre Compléments sur la platitude.
À l'origine, ce lemme comportait une preuve erronée (nous remercions Ofer Gabber
d'en avoir repéré la lacune) et servait dans une autre preuve de la proposition
citée. Pour rétablir ce lemme, il nous faut un argument correct lorsque $R$ est
un anneau intègre local normal, n'utilisant que des résultats des chapitres
d'algèbre commutative ; si vous connaissez un tel argument, veuillez écrire à
\href{mailto:stacks.project@gmail.com}{stacks.project@gmail.com}.
\end{proof}

\begin{lemma}
\label{lemma-relative-effective-cartier-algebra}
Soit $A \to B$ un homomorphisme d'anneaux. Soit $f \in B$. Supposons que
\begin{enumerate}
\item $A \to B$ soit plat,
\item $f$ soit un non-diviseur de zéro, et
\item $A \to B/fB$ soit plat.
\end{enumerate}
Alors, pour tout idéal $I \subset A$, l'homomorphisme
$f : B/IB \to B/IB$ est injectif.
\end{lemma}

\begin{proof}
Remarquons que $IB = I \otimes_A B$ et $I(B/fB) = I \otimes_A B/fB$
par platitude de $B$ et de $B/fB$ sur $A$.
En particulier, $IB/fIB \cong I \otimes_A B/fB$ s'injecte
dans $B/fB$. Le résultat découle donc du lemme du serpent appliqué
au diagramme
$$
\xymatrix{
0 \ar[r] &
I \otimes_A B \ar[r] \ar[d]^f &
B \ar[r] \ar[d]^f &
B/IB \ar[r] \ar[d]^f &
0 \\
0 \ar[r] &
I \otimes_A B \ar[r] &
B \ar[r] &
B/IB \ar[r] &
0
}
$$
dont les lignes sont exactes.
\end{proof}

\begin{lemma}
\label{lemma-faithfully-flat-injective}
Si $R \to S$ est un homomorphisme d'anneaux fidèlement plat, alors, pour tout
$R$-module $M$, l'homomorphisme $M \to S \otimes_R M$, $x \mapsto 1 \otimes x$,
est injectif.
\end{lemma}

\begin{proof}
Ce lemme est un doublon du lemme
\ref{algebra-lemma-faithfully-flat-universally-injective} du chapitre Algèbre.
\end{proof}

\begin{remark}
\label{remark-section-colimits}
Cette référence et cette étiquette renvoyaient autrefois à une section du
chapitre Lissage des homomorphismes d'anneaux, mais son contenu a depuis été
repris dans la section \ref{algebra-section-colimits-flat} du chapitre Algèbre.
\end{remark}

\begin{lemma}
\label{lemma-not-domain}
Soit $(R, \mathfrak m)$ un anneau local noethérien réduit de dimension $1$,
et soit $x \in \mathfrak m$ un non-diviseur de zéro. Soient
$\mathfrak q_1, \ldots, \mathfrak q_r$ les idéaux premiers minimaux de $R$.
Alors
$$
\text{length}_R(R/(x)) = \sum\nolimits_i \text{ord}_{R/\mathfrak q_i}(x)
$$
\end{lemma}

\begin{proof}
C'est un cas particulier (très simple) du lemme
\ref{chow-lemma-additivity-divisors-restricted} du chapitre Homologie de Chow.
\end{proof}

\begin{lemma}
\label{lemma-bound-primes}
Soit $A$ un anneau intègre local noethérien normal de dimension $2$.
Pour $f \in \mathfrak m$ non nul, notons
$\text{div}(f) = \sum n_i (\mathfrak p_i)$
le diviseur associé à $f$ sur le spectre épointé de $A$.
Posons $|f| = \sum n_i$. Il existe des entiers $N$ et $M$
tels que $|f + g| \leq M$ pour tout $g \in \mathfrak m^N$.
\end{lemma}

\begin{proof}
Choisissons $h \in \mathfrak m$ tel que $f, h$ soit une suite régulière dans $A$
(cela résulte des lemmes \ref{algebra-lemma-criterion-normal} et
\ref{algebra-lemma-depth-drops-by-one} du chapitre Algèbre).
Nous allons démontrer le lemme avec $M = \text{length}_A(A/(f, h))$ et avec
n'importe quel entier $N$ tel que $\mathfrak m^N \subset (f, h)$. Un tel
entier $N$ existe puisque $\sqrt{(f, h)} = \mathfrak m$. Remarquons que
$M = \text{length}_A(A/(f + g, h))$ pour tout $g \in \mathfrak m^N$
car $(f, h) = (f + g, h)$. Cela entraîne en outre que $f + g, h$
est encore une suite régulière dans $A$ ; voir le lemme
\ref{algebra-lemma-reformulate-CM} du chapitre Algèbre.
Supposons maintenant que $\text{div}(f + g ) = \sum m_j (\mathfrak q_j)$.
Considérons alors l'homomorphisme
$$
c : A/(f + g) \longrightarrow \prod A/\mathfrak q_j^{(m_j)},
$$
où $\mathfrak q_j^{(m_j)}$ est la puissance symbolique ; voir la section
\ref{algebra-section-symbolic-power} du chapitre Algèbre.
Comme $A$ est normal, $A_{\mathfrak q_i}$ est un anneau de valuation discrète et
$$
A_{\mathfrak q_i}/(f + g) =
A_{\mathfrak q_i}/\mathfrak q_i^{m_i} A_{\mathfrak q_i} =
(A/\mathfrak q_i^{(m_i)})_{\mathfrak q_i}
$$
Puisque $V(f + g, h) = \{\mathfrak m\}$, il en résulte que $c$ devient
un isomorphisme après inversion de $h$ (nous omettons un détail mineur).
Comme $h$ est un non-diviseur de zéro sur $A/(f + g)$, la longueur de
$A/(f + g, h)$ est égale au quotient de Herbrand
$e_A(A/(f + g), 0, h)$ défini dans la section
\ref{chow-section-periodic-complexes} du chapitre Homologie de Chow.
De même, la longueur de $A/(h, \mathfrak q_j^{(m_j)})$ est égale à
$e_A(A/\mathfrak q_j^{(m_j)}, 0, h)$. On a alors
\begin{align*}
M & = \text{length}_A(A/(f + g, h)) \\
& =
e_A(A/(f + g), 0, h) \\
& =
\sum\nolimits_j e_A(A/\mathfrak q_j^{(m_j)}, 0, h) \\
& =
\sum\nolimits_j \sum\nolimits_{m = 0, \ldots, m_j - 1}
e_A(\mathfrak q_j^{(m)}/\mathfrak q_j^{(m + 1)}, 0, h).
\end{align*}
Les égalités résultent des lemmes
\ref{chow-lemma-additivity-periodic-length} et
\ref{chow-lemma-finite-periodic-length} du chapitre Homologie de Chow,
en utilisant notamment le fait que
le conoyau de $c$ est de longueur finie, comme expliqué ci-dessus.
Le lemme de Nakayama permet de vérifier directement que
$e_A(\mathfrak q^{(m)}/\mathfrak q^{(m + 1)}, 0, h)$
est au moins égal à $1$. Cela achève la démonstration du lemme.
\end{proof}

\begin{lemma}
\label{lemma-flat-over-gorenstein-gorenstein-fibre}
Soit $A \to B$ un homomorphisme local plat d'anneaux locaux noethériens.
Si $A$ et $B/\mathfrak m_A B$ sont de Gorenstein, alors $B$ est de Gorenstein.
\end{lemma}

\begin{proof}
Cela résulte immédiatement du lemme
\ref{dualizing-lemma-flat-under-gorenstein} du chapitre Complexes dualisants.
\end{proof}

\begin{lemma}
\label{lemma-kill-local}
Soit $(A, \mathfrak m)$ un anneau local noethérien.
Soit $I \subset A$ un idéal. Soit $M$ un $A$-module de type fini.
Soit $s$ un entier. Supposons que
\begin{enumerate}
\item $A$ possède un complexe dualisant,
\item si $\mathfrak p \not \in V(I)$ et
$V(\mathfrak p) \cap V(I) \not = \{\mathfrak m\}$, alors
$\text{depth}_{A_\mathfrak p}(M_\mathfrak p) + \dim(A/\mathfrak p) > s$.
\end{enumerate}
Il existe alors un entier $n > 0$ et un idéal $J \subset A$
tels que $V(J) \cap V(I) = \{\mathfrak m\}$ et que $JI^n$ annule
$H^i_\mathfrak m(M)$ pour $i \leq s$.
\end{lemma}

\begin{proof}
D'après le lemme \ref{local-cohomology-lemma-sitting-in-degrees}
du chapitre Cohomologie locale, il suffit de le démontrer pour le $A$-module de type fini
$E^i = \text{Ext}^{-i}_A(M, \omega_A^\bullet)$
pour $i \leq s$. Le support $Z$ de $E^0 \oplus \ldots \oplus E^s$
est fermé dans $\Spec(A)$ et ne contient aucun idéal premier comme en (2).
Il est donc contenu dans $V(JI^n)$ pour un certain $J$ comme dans
l'énoncé du lemme.
\end{proof}

\begin{lemma}
\label{lemma-algebraize-local-cohomology-bis-bis}
Soit $(A, \mathfrak m)$ un anneau local noethérien.
Soit $I \subset A$ un idéal. Soit $M$ un $A$-module de type fini.
Soient $s$ et $d$ des entiers. Supposons que
\begin{enumerate}
\item[(a)] $A$ possède un complexe dualisant,
\item[(b)] $\text{cd}(A, I) \leq d$,
\item[(c)] si $\mathfrak p \not \in V(I)$, alors
$\text{depth}_{A_\mathfrak p}(M_\mathfrak p) > s$ ou
$\text{depth}_{A_\mathfrak p}(M_\mathfrak p) + \dim(A/\mathfrak p) > d + s$.
\end{enumerate}
Alors les hypothèses du lemme
\ref{algebraization-lemma-algebraize-local-cohomology-bis}
du chapitre Géométrie algébrique et formelle sont satisfaites
pour $A, I, \mathfrak m, M$, et l'homomorphisme
$H^i_\mathfrak m(M) \to \lim H^i_\mathfrak m(M/I^nM)$
est un isomorphisme pour $i \leq s$ ; ces modules sont en outre
annulés par une puissance de $I$.
\end{lemma}

\begin{proof}
Les hypothèses du lemme
\ref{algebraization-lemma-algebraize-local-cohomology-bis}
du chapitre Géométrie algébrique et formelle découlent du lemme plus général
\ref{algebraization-lemma-bootstrap-bis-bis} du même chapitre.
La conclusion du lemme
\ref{algebraization-lemma-algebraize-local-cohomology-bis}
donne alors la seconde assertion.
\end{proof}

\begin{lemma}
\label{lemma-combine-one}
Dans la situation \ref{algebraization-situation-bootstrap} du chapitre
Géométrie algébrique et formelle, on a
$H^s_\mathfrak a(M) = \lim H^s_\mathfrak a(M/I^nM)$.
\end{lemma}

\begin{proof}
Cela résulte immédiatement du théorème
\ref{algebraization-theorem-final-bootstrap} du chapitre Géométrie algébrique et formelle.
La version originale de ce lemme, qui comportait des hypothèses supplémentaires,
a été remplacée par ce théorème.
\end{proof}

\begin{lemma}
\label{lemma-fully-faithful-simple-one}
Soit $A$ un anneau noethérien. Soit $f \in \mathfrak a$ un élément d'un
idéal de $A$. Soit $U = \Spec(A) \setminus V(\mathfrak a)$. Supposons que
\begin{enumerate}
\item $A$ possède un complexe dualisant et soit complet pour la topologie $f$-adique,
\item $A_f$ soit $(S_2)$ et que, pour tout idéal premier minimal
$\mathfrak p \subset A$, on ait $f \not \in \mathfrak p$ et, pour tout
$\mathfrak q \in V(\mathfrak p) \cap V(\mathfrak a)$,
$\dim((A/\mathfrak p)_\mathfrak q) \geq 3$.
\end{enumerate}
Alors le foncteur de complétion
$$
\textit{Coh}(\mathcal{O}_U)
\longrightarrow
\textit{Coh}(U, I\mathcal{O}_U),
\quad
\mathcal{F} \longmapsto \mathcal{F}^\wedge
$$
est pleinement fidèle sur la sous-catégorie pleine des objets localement libres
de type fini.
\end{lemma}

\begin{proof}
Ce lemme est un cas particulier du lemme
\ref{algebraization-lemma-fully-faithful-simple-two}
du chapitre Géométrie algébrique et formelle.
\end{proof}



\section{Lemmes relatifs au th\'eor\`eme principal de Zariski}
\label{section-ZMT}

\noindent
Les lemmes de cette section servaient initialement dans la d\'emonstration de
la version alg\'ebrique du th\'eor\`eme principal de Zariski,
Alg\`ebre, Th\'eor\`eme \ref{algebra-theorem-main-theorem}.

\begin{lemma}
\label{lemma-change-equation-multiply}
Soient $R$ un anneau et $\varphi : R[x] \to S$ un homomorphisme
d'anneaux. Soit $t \in S$. Si $t$ est entier sur
$R[x]$, alors il existe un entier $\ell \geq 0$ tel que,
pour tout $a \in R$, l'\'el\'ement $\varphi(a)^\ell t$
soit entier relativement \`a $\varphi_a : R[y] \to S$, d\'efini par
$y \mapsto \varphi(ax)$ et $r \mapsto \varphi(r)$
pour $r\in R$.
\end{lemma}

\begin{proof}
\'Ecrivons $t^d + \sum_{i < d} \varphi(f_i)t^i = 0$
avec $f_i \in R[x]$. Soit $\ell$ le maximum des degr\'es
en $x$ de tous les $f_i$. En multipliant l'\'equation
par $\varphi(a)^\ell$, on obtient
$\varphi(a)^\ell t^d + \sum_{i < d} \varphi(a^\ell f_i)t^i = 0$.
Observons que chaque $\varphi(a^\ell f_i)$ appartient \`a l'image de
$\varphi_a$. Le r\'esultat d\'ecoule alors de
Alg\`ebre, Lemme \ref{algebra-lemma-make-integral-trivial}.
\end{proof}

\begin{lemma}
\label{lemma-make-integral-less-trivial}
Soit $\varphi : R \to S$ un homomorphisme d'anneaux.
Supposons que $t \in S$ v\'erifie la
relation $\varphi(a_0) + \varphi(a_1)t + \ldots + \varphi(a_n) t^n = 0$.
Posons $u_n = \varphi(a_n)$, $u_{n-1} = u_n t + \varphi(a_{n-1})$,
et ainsi de suite jusqu'\`a $u_1 = u_2 t + \varphi(a_1)$.
Alors les \'el\'ements $u_n, u_{n-1}, \ldots, u_1$ ainsi que
$u_nt, u_{n-1}t, \ldots, u_1t$ sont entiers sur $R$,
et les id\'eaux $(\varphi(a_0), \ldots, \varphi(a_n))$ et
$(u_n, \ldots, u_1)$ de $S$ sont \'egaux.
\end{lemma}

\begin{proof}
Nous raisonnons par r\'ecurrence sur $n$. Comme $u_n = \varphi(a_n)$,
il r\'esulte d'Alg\`ebre, Lemme \ref{algebra-lemma-make-integral-trivial},
que $u_nt$ est entier sur $R$. Bien entendu,
$u_n = \varphi(a_n)$ est entier sur $R$. Par cons\'equent,
$u_{n - 1} = u_n t  + \varphi(a_{n - 1})$ est entier sur $R$ (voir
Alg\`ebre, Lemme \ref{algebra-lemma-integral-closure-is-ring})
et l'on a
$$
\varphi(a_0) + \varphi(a_1)t + \ldots + \varphi(a_{n - 1})t^{n - 1} +
u_{n - 1}t^{n - 1} = 0.
$$
L'hypoth\`ese de r\'ecurrence, appliqu\'ee \`a l'homomorphisme
$S' \to S$, o\`u $S'$ est la cl\^oture int\'egrale de $R$ dans $S$,
et \`a l'\'equation affich\'ee, montre donc que
$u_{n-1}, \ldots, u_1$ et $u_{n-1}t, \ldots, u_1t$
appartiennent eux aussi tous \`a $S'$. L'assertion concernant les id\'eaux
d\'ecoule imm\'ediatement de la forme des \'el\'ements et du fait que
$u_1t + \varphi(a_0) = 0$.
\end{proof}

\begin{lemma}
\label{lemma-make-integral-not-in-ideal}
Soit $\varphi : R \to S$ un homomorphisme d'anneaux.
Supposons que $t \in S$ v\'erifie la
relation $\varphi(a_0) + \varphi(a_1)t + \ldots + \varphi(a_n) t^n = 0$.
Soit $J \subset S$ un id\'eal tel qu'il existe
au moins un $i$ pour lequel $\varphi(a_i) \not \in J$.
Alors il existe un \'el\'ement $u \in S$, $u \not\in J$, tel
que $u$ et $ut$ soient tous deux entiers sur $R$.
\end{lemma}

\begin{proof}
Cela d\'ecoule imm\'ediatement du Lemme \ref{lemma-make-integral-less-trivial},
car l'un des \'el\'ements $u_i$ n'appartient pas \`a $J$.
\end{proof}

\noindent
Les deux lemmes suivants permettent de d\'ecrire les sous-sch\'emas ferm\'es de
$\mathbf{P}^1_R$ d\'efinis par une seule \'equation (non d\'eg\'en\'er\'ee).

\begin{lemma}
\label{lemma-P1}
Soit $R$ un anneau.
Soit $F(X, Y) \in R[X, Y]$ un polyn\^ome homog\`ene de degr\'e
$d$. Supposons que, pour tout id\'eal premier $\mathfrak p$ de $R$,
au moins un coefficient de $F$ n'appartienne pas \`a $\mathfrak p$.
Munissons $S = R[X, Y]/(F)$ de sa structure d'anneau gradu\'e.
Alors, pour tout $n \geq d$, le $R$-module $S_n$
est localement libre de type fini, de rang $d$.
\end{lemma}

\begin{proof}
Le $R$-module $S_n$ admet une pr\'esentation
$$
R[X, Y]_{n-d} \to R[X, Y]_n \to S_n \to 0.
$$
D'apr\`es Alg\`ebre, Lemme \ref{algebra-lemma-cokernel-flat},
il suffit donc de montrer que la multiplication
par $F$ induit une application injective
$\kappa(\mathfrak p)[X, Y]
\to \kappa(\mathfrak p)[X, Y]$
pour tout id\'eal premier $\mathfrak p$.
Cela r\'esulte aussit\^ot de l'hypoth\`ese selon laquelle
l'image de $F$ modulo $\mathfrak p$ n'est pas le polyn\^ome nul.
L'assertion concernant les rangs en d\'ecoule elle aussi.
\end{proof}

\begin{lemma}
\label{lemma-rel-prime-pols}
Soit $k$ un corps. Soient $F, G \in k[X, Y]$ des polyn\^omes homog\`enes
de degr\'es respectifs $d, e$. Supposons que $F, G$ soient premiers entre eux.
Alors la multiplication par $G$ est injective sur $S = k[X, Y]/(F)$.
\end{lemma}

\begin{proof}
C'est une mani\`ere de d\'efinir l'expression \og premiers entre eux \fg. Si
l'on adopte une autre d\'efinition, on peut montrer qu'elle lui est \'equivalente.
\end{proof}

\begin{lemma}
\label{lemma-P1-localize}
Soit $R$ un anneau. Soit $F(X, Y) \in R[X, Y]$ homog\`ene de degr\'e
$d$. Munissons $S = R[X, Y]/(F)$ de sa structure d'anneau gradu\'e.
Soit $\mathfrak p \subset R$ un id\'eal premier tel qu'un
coefficient de $F$ n'appartienne pas \`a $\mathfrak p$.
Il existe un \'el\'ement $f \in R$, $f \not\in \mathfrak p$,
un entier $e$ et un \'el\'ement $G \in R[X, Y]_e$
tels que la multiplication par $G$ induise des isomorphismes
$(S_n)_f \to (S_{n + e})_f$ pour tout $n \geq d$.
\end{lemma}

\begin{proof}
Au cours de la d\'emonstration, nous pourrons remplacer $R$ par $R_f$
pour un $f\in R$, $f\not\in \mathfrak p$ (un nombre fini de fois).
Dans un premier temps, effectuons un tel remplacement de sorte
qu'un coefficient de $F$ soit inversible dans $R$.
En particulier, les modules $S_n$ sont d\`es lors localement
libres de rang $d$ pour $n \geq d$, d'apr\`es le Lemme \ref{lemma-P1}.
Choisissons un \'el\'ement $G \in R[X, Y]_e$ tel que les images de
$G$ dans $\kappa(\mathfrak p)[X, Y]$ et de
$F(X, Y)$ soient premi\`eres entre elles (un tel choix est possible pour un certain $e$).
Appliquons Alg\`ebre, Lemme \ref{algebra-lemma-cokernel-flat}, \`a l'application
induite par la multiplication par $G$ de $S_d \to S_{d + e}$.
Le choix de $G$ et le Lemme \ref{lemma-rel-prime-pols}
montrent que
$S_d \otimes \kappa(\mathfrak p) \to S_{d + e} \otimes \kappa(\mathfrak p)$
est bijective. Ainsi, apr\`es avoir remplac\'e $R$ par $R_f$ pour un
$f$ convenable, nous pouvons supposer que $G : S_d \to S_{d + e}$
est bijective. Il s'ensuit que l'image
de $G$ dans $\kappa(\mathfrak p')[X, Y]$ et celle de
$F$ sont premi\`eres entre elles, pour tout id\'eal premier $\mathfrak p'$
de $R$. En appliquant de nouveau Alg\`ebre, Lemme
\ref{algebra-lemma-cokernel-flat}, on voit alors que toutes les applications
$G : S_n \to S_{n + e}$, $n \geq d$, sont des isomorphismes.
\end{proof}

\begin{remark}
\label{remark-algebra}
Soit $R$ un anneau. Supposons donn\'es $F \in R[X, Y]_d$
et $G \in R[X, Y]_e$ tels que, si l'on pose $S = R[X, Y]/(F)$,
on ait (1) $S_n$ localement libre de type fini, de rang $d$, pour
tout $n \geq d$, et (2) que la multiplication par $G$ d\'efinisse des
isomorphismes $S_n \to S_{n + e}$ pour tout $n \geq d$. Dans ce
cas, on peut d\'efinir une $R$-alg\`ebre finie et localement libre
$A$ de la mani\`ere suivante :
\begin{enumerate}
\item comme $R$-module, $A = S_{ed}$ ;
\item la multiplication $A \times A \to A$ est donn\'ee par
la r\`egle suivante : $H_1 H_2 = H_3$ si et seulement si $G^d H_3 = H_1 H_2$
dans $S_{2ed}$.
\end{enumerate}
Cette d\'efinition est licite, car la multiplication par $G^d$
induit une application bijective $S_{de} \to S_{2de}$.
On v\'erifie ais\'ement qu'elle d\'efinit une structure d'anneau.
Notons le fait d\'econcertant que l'\'el\'ement $G^d$
d\'efinit l'\'el\'ement unit\'e de l'anneau $A$.
\end{remark}

\begin{lemma}
\label{lemma-finite-after-localization}
Soient $R$ un anneau et $f \in R$.
Supposons donn\'es $S$, $S'$ et les fl\`eches pleines
formant le diagramme commutatif d'anneaux suivant :
$$
\xymatrix{
& S'' \ar@{-->}[rd] \ar@{-->}[dd] &
\\
R \ar[rr] \ar@{-->}[ru] \ar[d] &  & S \ar[d]
\\
R_f \ar[r] & S' \ar[r] & S_f
}
$$
Supposons que $R_f \to S'$ soit fini. Alors on peut trouver
un homomorphisme d'anneaux fini $R \to S''$ et des fl\`eches pointill\'ees comme
dans le diagramme, de sorte que $S' = (S'')_f$.
\end{lemma}

\begin{proof}
En effet, supposons que $S'$ soit engendr\'ee par
les $x_i$ sur $R_f$, $i = 1, \ldots, w$. Soit $P_i(t) \in R_f[t]$
un polyn\^ome unitaire tel que $P_i(x_i) = 0$.
Le polyn\^ome $P_i$ est de degr\'e $d_i > 0$. \'Ecrivons
$P_i(t) = t^{d_i} + \sum_{j < d_i} (a_{ij}/f^n) t^j$
pour un m\^eme entier $n$. \'Ecrivons aussi
l'image de $x_i$ dans $S_f$ sous la forme $g_i / f^n$
avec un \'el\'ement convenable $g_i \in S$. On sait alors
que l'\'el\'ement
$\xi_i = f^{nd_i} g_i^{d_i} + \sum_{j < d_i} f^{n(d_i - j)} a_{ij} g_i^j$
de $S$ est annul\'e par une puissance de $f$.
En rempla\c cant $n$ par un entier $n'$ plus grand, ce qui revient \`a remplacer
$g_i$ par $f^{n' - n}g_i$, on peut donc supposer que $\xi_i = 0$.
L'alg\`ebre $S'$ est alors engendr\'ee par les \'el\'ements
$f^n x_i$, chacun \'etant un z\'ero du
polyn\^ome unitaire $Q_i(t) = t^{d_i} +
\sum_{j < d_i} f^{n(d_i - j)} a_{ij} t^j$
\`a coefficients dans $R$. Par construction, on a aussi
$Q_i(f^ng_i) = 0$ dans $S$. On obtient ainsi une $R$-alg\`ebre finie
$S'' = R[z_1, \ldots, z_w]/(Q_1(z_1), \ldots, Q_w(z_w))$
qui s'ins\`ere dans un diagramme commutatif comme ci-dessus.
L'homomorphisme $\alpha : S'' \to S$ envoie $z_i$ sur $f^ng_i$, et
l'homomorphisme $\beta : S'' \to S'$ envoie $z_i$ sur $f^nx_i$.
Il se peut que $\beta$ n'induise pas encore un
isomorphisme $(S'')_f \cong S'$.
Pour l'instant, on sait seulement que cette application
est surjective. Le probl\`eme est qu'il pourrait exister des
\'el\'ements $h/f^n \in (S'')_f$ qui s'envoient sur z\'ero
dans $S'$ sans \^etre eux-m\^emes nuls. Dans ce cas, $\alpha(h)$
est un \'el\'ement de $S$ tel que $f^N \alpha(h) = 0$
pour un certain $N$. Ainsi, $f^N h$ appartient \`a l'id\'eal
$J = \{h \in S'' \mid \alpha(h) = 0 \text{ et }
\beta(h) = 0\}$ de $S''$. Il est alors facile de voir que
$S''/J$ convient.
\end{proof}
\section{Homomorphismes formellement lisses}
\label{section-formally-smooth}

\begin{lemma}
\label{lemma-formally-smooth-smooth}
Soit $R$ un anneau. Soit $S$ une $R$-alg\`ebre.
Si $S$ est de pr\'esentation finie et formellement lisse sur $R$,
alors $S$ est lisse sur $R$.
\end{lemma}

\begin{proof}
Voir Alg\`ebre, Proposition \ref{algebra-proposition-smooth-formally-smooth}.
\end{proof}

\begin{remark}
\label{remark-equation-derivatives}
Cette \`etiquette renvoyait \`a une \`equation de la d\'emonstration de
la Proposition \ref{restricted-proposition-approximate}
d'Alg\'ebrisation des espaces formels.
Apr\`es r\'eorganisation du contenu, cette \`etiquette n'est plus utilis\'ee.
\end{remark}

\begin{remark}
\label{remark-equation-ci}
Cette \`etiquette renvoyait \`a une \`equation de la d\'emonstration de
la Proposition \ref{restricted-proposition-approximate}
d'Alg\'ebrisation des espaces formels.
Apr\`es r\'eorganisation du contenu, cette \`etiquette n'est plus utilis\'ee.
\end{remark}

\begin{remark}
\label{remark-equation-in-ideal}
Cette \`etiquette renvoyait \`a une \`equation de la d\'emonstration de
la Proposition \ref{restricted-proposition-approximate}
d'Alg\'ebrisation des espaces formels.
Apr\`es r\'eorganisation du contenu, cette \`etiquette n'est plus utilis\'ee.
\end{remark}

\begin{remark}
\label{remark-equation-derivatives-analogue}
Cette \`etiquette renvoyait \`a une \`equation de la d\'emonstration de
la Proposition \ref{restricted-proposition-approximate}
d'Alg\'ebrisation des espaces formels.
Apr\`es r\'eorganisation du contenu, cette \`etiquette n'est plus utilis\'ee.
\end{remark}

\begin{remark}
\label{remark-equation-go-down}
Cette \`etiquette renvoyait \`a une \`equation de la d\'emonstration du
Lemme \ref{restricted-lemma-lift-approximation}
d'Alg\'ebrisation des espaces formels.
Apr\`es r\'eorganisation du contenu, cette \`etiquette n'est plus utilis\'ee.
\end{remark}

\begin{lemma}
\label{lemma-get-morphism-general}
Soit $A$ un anneau noeth\'erien. Soit $I \subset A$ un id\'eal.
Soit $t$ le nombre minimum de g\'en\'erateurs de $I$.
Soit $C$ une $A$-alg\`ebre noeth\'erienne $I$-adiquement compl\`ete.
Il existe un entier $d \geq 0$ ne d\'ependant que de
$I \subset A \to C$ et poss\'edant la propri\'et\'e suivante : \`etant donn\'es
\begin{enumerate}
\item un entier $c \geq 0$ et l'alg\`ebre $B$ de l'\`Equation
(\ref{restricted-equation-C-prime}) d'Alg\'ebrisation des espaces formels,
tels que, pour tout $a \in I^c$,
la multiplication par $a$ sur $\NL_{B/A}^\wedge$ soit nulle dans $D(B)$,
\item un entier $n > 2t\max(c, d)$,
\item un homomorphisme de $A/I^n$-alg\`ebres
$\psi_n : B/I^nB \to C/I^nC$,
\end{enumerate}
il existe un homomorphisme de $A$-alg\`ebres $\varphi : B \to C$ tel
que $\psi_n \bmod I^{m - c} = \varphi \bmod I^{m - c}$,
o\`u $m = \lfloor \frac{n}{t} \rfloor$.
\end{lemma}

\begin{proof}
Ce lemme est devenu obsol\`ete en raison du r\'esultat plus fort
d'Alg\'ebrisation des espaces formels,
Lemme \ref{restricted-lemma-get-morphism-general-better}.
Nous allons en fait l'en d\'eduire.

\medskip\noindent
Soit $I \subset A \to C$ comme dans l'\'enonc\'e ci-dessus.
Notons $d(\text{Gr}_I(C))$ et $q(\text{Gr}_I(C))$ les entiers obtenus dans
Cohomologie locale, Section \ref{local-cohomology-section-uniform}.
Remarquons que $t$ majore le nombre minimum de g\'en\'erateurs de $IC$;
on a donc $d(\text{Gr}_I(C)) + 1 \leq t$, voir la discussion de
Cohomologie locale, Section \ref{local-cohomology-section-uniform}.
Nous pouvons supposer, et nous supposerons, que $t \geq 1$, car sinon le lemme
n'affirme rien. Nous affirmons que le lemme est vrai en prenant
$$
d = q(\text{Gr}_I(C))
$$
En effet, supposons que $c$, $B$, $n$ et $\psi_n$ soient comme dans l'\'enonc\'e
ci-dessus. On a alors
$$
n > 2t\max(c, d) \Rightarrow n \geq 2tc + 1 \Rightarrow
n \geq 2(d(\text{Gr}_I(C)) + 1)c + 1
$$
D'autre part,
$$
n > 2t\max(c, d) \Rightarrow n > t(c + d) \Rightarrow
n \geq q(C) + tc \geq q(\text{Gr}_I(C)) + (d(\text{Gr}_I(C)) + 1)c
$$
Les hypoth\`eses du Lemme \ref{restricted-lemma-get-morphism-general-better}
d'Alg\'ebrisation des espaces formels sont donc satisfaites,
et nous obtenons un homomorphisme de $A$-alg\`ebres
$\varphi : B \to C$ congru \`a $\psi_n$
modulo $I^{n - (d(\text{Gr}_I(C)) + 1)c}C$.
Comme
\begin{align*}
n - (d(\text{Gr}_I(C)) + 1)c
& = \frac{n}{t} + \frac{(t - 1)n}{t} - (d(\text{Gr}_I(C)) + 1)c \\
& \geq \frac{n}{t} + \frac{d(\text{Gr}_I(C))n}{t} - (d(\text{Gr}_I(C)) + 1)c \\
& > \frac{n}{t} + \frac{d(\text{Gr}_I(C))2tc}{t} - (d(\text{Gr}_I(C)) + 1)c \\
& = \frac{n}{t} + 2d(\text{Gr}_I(C))c - (d(\text{Gr}_I(C)) + 1)c \\
& = \frac{n}{t} + d(\text{Gr}_I(C))c - c \\
& \geq m - c
\end{align*}
on voit que $\varphi$ et $\psi_n$ sont congrus modulo $I^{m - c}C$, comme voulu.
\end{proof}
\section{Sites et faisceaux}
\label{section-sites}

\begin{remark}[Absence de morphisme du prolongement par zéro vers l'image directe]
\label{remark-from-shriek-to-star}
Soit $U$ un objet d'un site $\mathcal{C}$. Pour tout faisceau abélien
$\mathcal{G}$ sur $\mathcal{C}/U$, on peut se demander s'il existe
un morphisme canonique
$$
c : j_{U!}\mathcal{G} \longrightarrow j_{U*}\mathcal{G}
$$
Construire un tel morphisme revient à construire un morphisme
$j_U^{-1}j_{U!}\mathcal{G} \to \mathcal{G}$.
Remarquons que la restriction commute au passage au faisceau associé.
Nous pouvons donc utiliser le préfaisceau du lemme
Modules sur les sites, lemme \ref{sites-modules-lemma-extension-by-zero}.
Il suffit ainsi de définir, pour $V/U$, un morphisme
$$
\bigoplus\nolimits_{\varphi \in \Mor_\mathcal{C}(V, U)}
\mathcal{G}(V \xrightarrow{\varphi} U)
\longrightarrow
\mathcal{G}(V/U)
$$
compatible aux restrictions. Il semble que l'on puisse prendre le morphisme
nul sur tous les facteurs directs, sauf sur celui pour lequel $\varphi$
est le morphisme structural $\varphi_0 : V \to U$, où l'on prend $1$.
Cependant, ce morphisme n'est pas compatible aux applications de restriction.
En effet, si $\alpha : V' \to V$ est un morphisme de $\mathcal{C}$, notons
$V'/U$ l'objet de $\mathcal{C}/U$ dont le morphisme structural est
$\varphi'_0 = \varphi_0 \circ \alpha$.
Il nous faut vérifier que le diagramme
$$
\xymatrix{
\bigoplus\nolimits_{\varphi \in \Mor_\mathcal{C}(V, U)}
\mathcal{G}(V \xrightarrow{\varphi} U)
\ar[d] \ar[r] &
\mathcal{G}(V/U) \ar[d] \\
\bigoplus\nolimits_{\varphi' \in \Mor_\mathcal{C}(V', U)}
\mathcal{G}(V' \xrightarrow{\varphi'} U)
\ar[r] &
\mathcal{G}(V'/U)
}
$$
est commutatif. Le problème est qu'il peut exister un morphisme
$\varphi : V \to U$, distinct de $\varphi_0$, tel que
$\varphi \circ \alpha = \varphi'_0$.
La flèche verticale de gauche envoie alors le facteur direct correspondant
à $\varphi$ dans celui sur lequel la flèche horizontale inférieure est
égale à $1$, et il est presque certain que le diagramme ne commute pas.
\end{remark}




\section{Cohomologie}
\label{section-cohomology}

\noindent
Lemmes devenus obsolètes concernant la cohomologie.

\begin{lemma}
\label{lemma-ML-general}
Soit $I$ un idéal d'un anneau $A$. Soit $X$ un schéma sur $\Spec(A)$. Soit
$$
\ldots \to \mathcal{F}_3 \to \mathcal{F}_2 \to \mathcal{F}_1
$$
un système projectif de $\mathcal{O}_X$-modules tel que
$\mathcal{F}_n = \mathcal{F}_{n + 1}/I^n\mathcal{F}_{n + 1}$.
Supposons que
$$
\bigoplus\nolimits_{n \geq 0} H^1(X, I^n\mathcal{F}_{n + 1})
$$
satisfasse la condition de chaîne croissante en tant que
$\bigoplus_{n \geq 0} I^n/I^{n + 1}$-module gradué.
Alors le système projectif $M_n = \Gamma(X, \mathcal{F}_n)$ satisfait
la condition de Mittag-Leffler.
\end{lemma}

\begin{proof}
C'est un cas particulier du lemme plus général
Cohomologie, lemme \ref{cohomology-lemma-ML-general}.
\end{proof}

\begin{lemma}
\label{lemma-ML-general-better}
Soit $I$ un idéal d'un anneau $A$. Soit $X$ un schéma sur $\Spec(A)$. Soit
$$
\ldots \to \mathcal{F}_3 \to \mathcal{F}_2 \to \mathcal{F}_1
$$
un système projectif de $\mathcal{O}_X$-modules tel que
$\mathcal{F}_n = \mathcal{F}_{n + 1}/I^n\mathcal{F}_{n + 1}$.
Pour tout $n$, posons
$$
H^1_n =
\bigcap\nolimits_{m \geq n}
\Im\left(
H^1(X, I^n\mathcal{F}_{m + 1}) \to H^1(X, I^n\mathcal{F}_{n + 1})
\right)
$$
Si $\bigoplus H^1_n$ satisfait la condition de chaîne croissante en tant que
$\bigoplus_{n \geq 0} I^n/I^{n + 1}$-module gradué, alors le système
projectif $M_n = \Gamma(X, \mathcal{F}_n)$ satisfait la condition
de Mittag-Leffler.
\end{lemma}

\begin{proof}
C'est un cas particulier du lemme plus général
Cohomologie, lemme \ref{cohomology-lemma-ML-general-better}.
\end{proof}

\begin{lemma}
\label{lemma-topology-I-adic-general}
Soit $I$ un idéal de type fini d'un anneau $A$.
Soit $X$ un schéma sur $\Spec(A)$. Soit
$$
\ldots \to \mathcal{F}_3 \to \mathcal{F}_2 \to \mathcal{F}_1
$$
un système projectif de $\mathcal{O}_X$-modules tel que
$\mathcal{F}_n = \mathcal{F}_{n + 1}/I^n\mathcal{F}_{n + 1}$.
Supposons que
$$
\bigoplus\nolimits_{n \geq 0} H^0(X, I^n\mathcal{F}_{n + 1})
$$
satisfasse la condition de chaîne croissante en tant que
$\bigoplus_{n \geq 0} I^n/I^{n + 1}$-module gradué.
Alors la topologie limite sur $M = \lim \Gamma(X, \mathcal{F}_n)$
est la topologie $I$-adique.
\end{lemma}

\begin{proof}
C'est un cas particulier du lemme plus général
Cohomologie, lemme \ref{cohomology-lemma-topology-I-adic-general}.
\end{proof}

\begin{lemma}
\label{lemma-pullback-K-flat}
Soit $(\Sh(\mathcal{C}), \mathcal{O}_\mathcal{C})$ un topos annelé.
Pour tout complexe de $\mathcal{O}_\mathcal{C}$-modules
$\mathcal{G}^\bullet$, il existe un quasi-isomorphisme
$\mathcal{K}^\bullet \to \mathcal{G}^\bullet$ tel que
$f^*\mathcal{K}^\bullet$ soit un complexe K-plat de
$\mathcal{O}_\mathcal{D}$-modules pour tout morphisme de topos annelés
$f : (\Sh(\mathcal{D}), \mathcal{O}_\mathcal{D}) \to
(\Sh(\mathcal{C}), \mathcal{O}_\mathcal{C})$.
\end{lemma}

\begin{proof}
Cela résulte de Cohomologie sur les sites, lemmes
\ref{sites-cohomology-lemma-K-flat-resolution} et
\ref{sites-cohomology-lemma-pullback-K-flat}.
\end{proof}

\begin{remark}
\label{remark-pullback-K-flat}
Cette remarque examinait autrefois ce que l'on savait du caractère K-plat,
ou non, des images réciproques de complexes K-plats; elle est désormais
rendue obsolète par Cohomologie sur les sites, lemme
\ref{sites-cohomology-lemma-pullback-K-flat}.
\end{remark}

\noindent
Le lemme suivant calcule, dans un cas particulier, les faisceaux de
cohomologie de la limite dérivée.

\begin{lemma}
\label{lemma-Rlim-of-system}
Soit $(\mathcal{C}, \mathcal{O})$ un site annelé. Soit $(K_n)$
un système projectif d'objets de $D(\mathcal{O})$.
Soit $\mathcal{B} \subset \Ob(\mathcal{C})$ un sous-ensemble.
Soit $d \in \mathbf{N}$. Supposons que
\begin{enumerate}
\item $K_n$ soit un objet de $D^+(\mathcal{O})$ pour tout $n$,
\item pour tout $q \in \mathbf{Z}$, il existe
$n(q)$ tel que $H^q(K_{n + 1}) \to H^q(K_n)$ soit un isomorphisme pour
$n \geq n(q)$,
\item tout objet de $\mathcal{C}$ admette un recouvrement dont les membres
appartiennent à $\mathcal{B}$,
\item pour tout $U \in \mathcal{B}$, on ait $H^p(U, H^q(K_n)) = 0$
pour $p > d$ et tout $q$.
\end{enumerate}
Alors $H^m(R\lim K_n) = \lim H^m(K_n)$ pour tout $m \in \mathbf{Z}$.
\end{lemma}

\begin{proof}
Posons $K = R\lim K_n$. Soit $U \in \mathcal{B}$. Pour chaque $n$, il existe
une suite spectrale
$$
H^p(U, H^q(K_n)) \Rightarrow H^{p + q}(U, K_n)
$$
qui converge puisque $K_n$ est borné inférieurement; voir
Catégories dérivées, lemme \ref{derived-lemma-two-ss-complex-functor}.
Fixons $m \in \mathbf{Z}$. L'hypothèse (4) montre que seuls les groupes
$H^p(U, H^q(K_n))$ contribuent à $H^m(U, K_n)$ pour
$0 \leq p \leq d$ et $m - d \leq q \leq m$.
L'hypothèse (2) implique donc que
$H^m(U, K_{n + 1}) \to H^m(U, K_n)$ est un isomorphisme dès que
$n \geq \max{n(m), \ldots, n(m - d)}$. Le foncteur $R\Gamma(U, -)$
commute aux limites dérivées d'après
Injectifs, lemme \ref{injectives-lemma-RF-commutes-with-Rlim}.
Ainsi,
$$
H^m(U, K) = H^m(R\lim R\Gamma(U, K_n))
$$
D'autre part, nous venons de voir que les groupes de cohomologie des
complexes $R\Gamma(U, K_n)$ sont stationnaires à partir d'un certain rang.
Il résulte donc de Compléments d'algèbre, remarque
\ref{more-algebra-remark-compare-derived-limit}, que $H^m(U, K)$ est égal
à $H^m(U, K_n)$ pour tout $n \gg 0$, avec une borne indépendante de
$U \in \mathcal{B}$. Fixons un tel $n$. Enfin, rappelons que $H^m(K)$
est le faisceau associé au préfaisceau $U \mapsto H^m(U, K)$, tandis que
$H^m(K_n)$ est le faisceau associé au préfaisceau
$U \mapsto H^m(U, K_n)$. Sur les éléments de $\mathcal{B}$, ces
préfaisceaux prennent les mêmes valeurs. L'hypothèse (3) garantit donc
que les faisceaux associés sont eux aussi identiques. Le lemme en résulte.
\end{proof}

\begin{lemma}
\label{lemma-trivialities-cohomological-descent-abelian}
Dans la situation de Espaces simpliciaux,
\ref{spaces-simplicial-situation-simplicial-site}, soit $a_0$ une
augmentation vers un site $\mathcal{D}$, comme dans Espaces simpliciaux,
remarque \ref{spaces-simplicial-remark-augmentation-site}.
Supposons données des sous-catégories de Serre faibles strictement pleines
$$
\mathcal{A} \subset \textit{Ab}(\mathcal{D}),\quad
\mathcal{A}_n \subset \textit{Ab}(\mathcal{C}_n)
$$
Alors
\begin{enumerate}
\item[(1)]
la collection des faisceaux abéliens $\mathcal{F}$ sur
$\mathcal{C}_{total}$ dont la restriction à $\mathcal{C}_n$ appartient à
$\mathcal{A}_n$ pour tout $n$ est une sous-catégorie de Serre faible
strictement pleine
$\mathcal{A}_{total} \subset \textit{Ab}(\mathcal{C}_{total})$.
\end{enumerate}
Si $a_n^{-1}$ envoie $\mathcal{A}$ dans $\mathcal{A}_n$ pour tout $n$, alors
\begin{enumerate}
\item[(2)] $a^{-1}$ envoie $\mathcal{A}$ dans $\mathcal{A}_{total}$, et
\item[(3)] $a^{-1}$ envoie $D_\mathcal{A}(\mathcal{D})$ dans
$D_{\mathcal{A}_{total}}(\mathcal{C}_{total})$.
\end{enumerate}
Si $R^qa_{n, *}$ envoie $\mathcal{A}_n$ dans $\mathcal{A}$
pour tous $n, q$, alors
\begin{enumerate}
\item[(4)] $R^qa_*$ envoie $\mathcal{A}_{total}$ dans $\mathcal{A}$ pour tout
$q$, et
\item[(5)] $Ra_*$ envoie
$D_{\mathcal{A}_{total}}^+(\mathcal{C}_{total})$
dans $D_\mathcal{A}^+(\mathcal{D})$.
\end{enumerate}
\end{lemma}

\begin{proof}
Les seules assertions qui demandent une explication sont (4) et (5).
L'assertion (4) résulte de la suite spectrale de Espaces simpliciaux,
lemme \ref{spaces-simplicial-lemma-augmentation-spectral-sequence},
et de Homologie, lemme \ref{homology-lemma-biregular-ss-converges}.
L'assertion (5) s'en déduit en considérant la suite spectrale associée à
la filtration canonique par les troncations d'un objet
$K$ de $D_{\mathcal{A}_{total}}^+(\mathcal{C}_{total})$.
Nous omettons les détails.
\end{proof}

\begin{remark}
\label{remark-cohomology-topics}
Ce tag renvoyait autrefois à une section du chapitre consacré à la
cohomologie qui énumérait les sujets restant à traiter.
\end{remark}

\begin{remark}
\label{remark-sites-cohomology-topics}
Ce tag renvoyait autrefois à une section du chapitre consacré à la
cohomologie qui énumérait les sujets restant à traiter.
\end{remark}

\begin{remark}
\label{remark-V-implies-C}
Ce tag renvoyait autrefois au cas particulier de
Cohomologie sur les sites, lemme
\ref{sites-cohomology-lemma-V-implies-C-general}, dans la situation décrite
par Cohomologie sur les sites, lemme
\ref{sites-cohomology-lemma-compare-qc-zar}.
\end{remark}

\begin{remark}
\label{remark-V-implies-cohomology}
Ce tag renvoyait autrefois au cas particulier de
Cohomologie sur les sites, lemme
\ref{sites-cohomology-lemma-V-implies-cohomology-general}, dans la situation
décrite par Cohomologie sur les sites, lemme
\ref{sites-cohomology-lemma-compare-qc-zar}.
\end{remark}

\begin{remark}
\label{remark-induction-step-V-C}
Ce tag renvoyait autrefois au cas particulier de
Cohomologie sur les sites, lemme
\ref{sites-cohomology-lemma-induction-step-V-C-general}, dans la situation
décrite par Cohomologie sur les sites, lemme
\ref{sites-cohomology-lemma-compare-qc-zar}.
\end{remark}

\begin{remark}
\label{remark-V-implies-C-etale-fppf}
Ce tag renvoyait autrefois au cas particulier de
Cohomologie sur les sites, lemme
\ref{sites-cohomology-lemma-V-implies-C-general}, dans la situation décrite
par Cohomologie étale, lemme
\ref{etale-cohomology-lemma-compare-fppf-etale}.
\end{remark}

\begin{remark}
\label{remark-V-implies-cohomology-etale-fppf}
Ce tag renvoyait autrefois au cas particulier de
Cohomologie sur les sites, lemme
\ref{sites-cohomology-lemma-V-implies-cohomology-general}, dans la situation
décrite par Cohomologie étale, lemme
\ref{etale-cohomology-lemma-compare-fppf-etale}.
\end{remark}

\begin{remark}
\label{remark-induction-step-V-C-etale-fppf}
Ce tag renvoyait autrefois au cas particulier de
Cohomologie sur les sites, lemme
\ref{sites-cohomology-lemma-induction-step-V-C-general}, dans la situation
décrite par Cohomologie étale, lemme
\ref{etale-cohomology-lemma-compare-fppf-etale}.
\end{remark}

\begin{remark}
\label{remark-V-implies-C-etale-ph}
Ce tag renvoyait autrefois au cas particulier de
Cohomologie sur les sites, lemme
\ref{sites-cohomology-lemma-V-implies-C-general}, dans la situation décrite
par Cohomologie étale, lemme
\ref{etale-cohomology-lemma-compare-ph-etale}.
\end{remark}

\begin{remark}
\label{remark-V-implies-cohomology-etale-ph}
Ce tag renvoyait autrefois au cas particulier de
Cohomologie sur les sites, lemme
\ref{sites-cohomology-lemma-V-implies-cohomology-general}, dans la situation
décrite par Cohomologie étale, lemme
\ref{etale-cohomology-lemma-compare-ph-etale}.
\end{remark}

\begin{remark}
\label{remark-V-implies-cohomology-etale-ph-extra}
Ce tag renvoyait autrefois au cas particulier de
Cohomologie sur les sites, lemme
\ref{sites-cohomology-lemma-V-implies-cohomology-extra-general}, dans la
situation décrite par Cohomologie étale, lemme
\ref{etale-cohomology-lemma-compare-ph-etale}.
\end{remark}

\begin{remark}
\label{remark-make-class-zero}
Ce tag renvoyait autrefois au cas particulier de
Cohomologie sur les sites, lemme
\ref{sites-cohomology-lemma-make-class-zero-general}, dans la situation
décrite par Cohomologie étale, lemme
\ref{etale-cohomology-lemma-compare-ph-etale}.
\end{remark}

\begin{remark}
\label{remark-induction-step-V-C-etale-ph}
Ce tag renvoyait autrefois au cas particulier de
Cohomologie sur les sites, lemme
\ref{sites-cohomology-lemma-induction-step-V-C-general}, dans la situation
décrite par Cohomologie étale, lemme
\ref{etale-cohomology-lemma-compare-ph-etale}.
\end{remark}

\begin{remark}
\label{remark-V-implies-C-etale-h}
Ce tag renvoyait autrefois au cas particulier de
Cohomologie sur les sites, lemme
\ref{sites-cohomology-lemma-V-implies-C-general}, dans la situation décrite
par Cohomologie étale, lemme
\ref{etale-cohomology-lemma-compare-h-etale}.
\end{remark}

\begin{remark}
\label{remark-V-implies-cohomology-etale-h}
Ce tag renvoyait autrefois au cas particulier de
Cohomologie sur les sites, lemme
\ref{sites-cohomology-lemma-V-implies-cohomology-general}, dans la situation
décrite par Cohomologie étale, lemme
\ref{etale-cohomology-lemma-compare-h-etale}.
\end{remark}

\begin{remark}
\label{remark-V-implies-cohomology-etale-h-extra}
Ce tag renvoyait autrefois au cas particulier de
Cohomologie sur les sites, lemme
\ref{sites-cohomology-lemma-V-implies-cohomology-extra-general}, dans la
situation décrite par Cohomologie étale, lemme
\ref{etale-cohomology-lemma-compare-h-etale}.
\end{remark}

\begin{remark}
\label{remark-induction-step-V-C-etale-h}
Ce tag renvoyait autrefois au cas particulier de
Cohomologie sur les sites, lemme
\ref{sites-cohomology-lemma-induction-step-V-C-general}, dans la situation
décrite par Cohomologie étale, lemme
\ref{etale-cohomology-lemma-compare-h-etale}.
\end{remark}

\begin{remark}
\label{remark-how-used}
Ce tag se trouvait autrefois dans le chapitre consacré à la cohomologie
étale, mais une réorganisation fait qu'il n'y a désormais plus sa place.
Son contenu était le suivant :
Cohomologie étale, lemme \ref{etale-cohomology-lemma-when-ctf},
permet de prouver que, si $f : X \to Y$ est un morphisme de schémas séparé
et de type fini et si $Y$ est noethérien, alors $Rf_!$ induit un foncteur
$D_{ctf}(X_\etale, \Lambda) \to D_{ctf}(Y_\etale, \Lambda)$.
Un exemple de cet argument, lorsque $Y$ est le spectre d'un corps et que
$X$ est une courbe, est donné dans La formule des traces,
proposition \ref{trace-proposition-projective-curve-constructible-cohomology}.
\end{remark}

\begin{lemma}
\label{lemma-glue-f-upper-shriek}
Soit $f : X \to Y$ un morphisme de schémas localement quasi-fini.
Il existe un unique foncteur
$f^! : \textit{Ab}(Y_\etale) \to \textit{Ab}(X_\etale)$ tel que
\begin{enumerate}
\item pour toute immersion ouverte $j : U \to X$ telle que $f \circ j$
soit séparé, il existe un isomorphisme canonique
$j^! \circ f^! = (f \circ j)^!$, et
\item ces isomorphismes, pour $U \subset U' \subset X$, soient compatibles
avec ceux de Compléments de cohomologie étale, lemme
\ref{more-etale-lemma-upper-shriek-restriction}.
\end{enumerate}
\end{lemma}

\begin{proof}
C'est une conséquence immédiate de Compléments de cohomologie étale, lemmes
\ref{more-etale-lemma-lqf-f-upper-shriek} et
\ref{more-etale-lemma-upper-shriek-restriction}.
\end{proof}

\begin{proposition}
\label{proposition-lqf-shriek}
Soit $f : X \to Y$ un morphisme localement quasi-fini. Il existe des
foncteurs adjoints $f_! : \textit{Ab}(X_\etale) \to \textit{Ab}(Y_\etale)$
et $f^! : \textit{Ab}(Y_\etale) \to \textit{Ab}(X_\etale)$
ayant les propriétés suivantes :
\begin{enumerate}
\item le foncteur $f^!$ est celui construit dans Compléments de cohomologie
étale, lemme \ref{more-etale-lemma-lqf-f-upper-shriek},
\item pour toute immersion ouverte $j : U \to X$ telle que $f \circ j$
soit séparé, il existe un isomorphisme canonique
$f_! \circ j_! = (f \circ j)_!$, et
\item ces isomorphismes, pour $U \subset U' \subset X$, sont compatibles
avec ceux de Compléments de cohomologie étale, lemme
\ref{more-etale-lemma-f-shriek-composition}.
\end{enumerate}
\end{proposition}

\begin{proof}
Voir Compléments de cohomologie étale, sections
\ref{more-etale-section-finite-support} et
\ref{more-etale-section-duality-locally-quasi-finite}.
\end{proof}

\begin{lemma}
\label{lemma-lqf-f-upper-shriek-stalk}
Soit $f : X \to Y$ un morphisme de schémas localement quasi-fini.
Pour un groupe abélien $A$ et un point géométrique
$\overline{y} : \Spec(k) \to Y$, on a
$f^!(\overline{y}_*A) = \prod\nolimits_{f(\overline{x}) = \overline{y}}
\overline{x}_*A$.
\end{lemma}

\begin{proof}
Cela résulte de l'assertion correspondante de Compléments de cohomologie
étale, lemme \ref{more-etale-lemma-lqf-f-upper-shriek}.
\end{proof}

\begin{lemma}
\label{lemma-lqf-f-shriek-composition}
Soient $f : X \to Y$ et $g : Y \to Z$ des morphismes de schémas
localement quasi-finis et composables. Alors
$g_! \circ f_! = (g \circ f)_!$ et
$f^! \circ g^! = (g \circ f)^!$.
\end{lemma}

\begin{proof}
Cela résulte de la combinaison des lemmes de Compléments de cohomologie étale
\ref{more-etale-lemma-lqf-shriek-composition} et
\ref{more-etale-lemma-upper-shriek-restriction}.
\end{proof}
\section{Algèbre différentielle graduée}
\label{section-dga}


\begin{lemma}
\label{lemma-P-not-preserved-base-change}
Soient $(A, \text{d})$ et $(B, \text{d})$ des algèbres différentielles graduées.
Soit $N$ un $(A, B)$-bimodule différentiel gradué possédant la propriété
(P). Soit $M$ un $A$-module différentiel gradué possédant la propriété (P).
Alors $Q = M \otimes_A N$ est un $B$-module différentiel gradué qui représente
$M \otimes_A^\mathbf{L} N$ dans $D(B)$ et qui admet une filtration
$$
0 = F_{-1}Q \subset F_0Q \subset F_1Q \subset \ldots \subset Q
$$
par des sous-modules différentiels gradués telle que $Q = \bigcup F_pQ$,
que les inclusions $F_iQ \to F_{i + 1}Q$ soient des monomorphismes admissibles,
et que les quotients $F_{i + 1}Q/F_iQ$ soient isomorphes, comme
$B$-modules différentiels gradués, à une somme directe de modules de la forme
$(A \otimes_R B)[k]$.
\end{lemma}

\begin{proof}
Choisissons des filtrations $F_\bullet$ sur $M$ et $N$. Considérons alors la filtration
de $Q = M \otimes_A N$ définie par
$$
F_n(Q) = \sum\nolimits_{i + j = n}
F_i(M) \otimes_A F_j(N)
$$
Il s'agit manifestement d'un $B$-sous-module différentiel gradué. On a
$$
F_n(Q)/F_{n - 1}(Q) =
\bigoplus\nolimits_{i + j = n}
F_i(M)/F_{i - 1}(M) \otimes_A F_j(N)/F_{j - 1}(N)
$$
par exemple parce que la filtration de $M$ est scindée dans la catégorie
des $A$-modules gradués. Par hypothèse, les quotients du membre de droite
sont isomorphes à des sommes directes de décalés de $A$ et de
$A \otimes_R B$; comme
$A \otimes_A (A \otimes_R B) = A \otimes_R B$,
on en conclut que le membre de gauche est une somme directe de décalés
de $A \otimes_R B$ comme $B$-module différentiel gradué.
(Attention : $Q$ n'est pas muni d'une structure de $(A, B)$-bimodule.)
Cela établit l'assertion relative à la filtration.
L'assertion relative au produit tensoriel dérivé résulte immédiatement
de la définition du foncteur dans
Algèbre différentielle graduée, lemme \ref{dga-lemma-derived-bc}.
\end{proof}













\section{Méthodes simpliciales}
\label{section-simplicial}


\begin{lemma}
\label{lemma-equiv}
Sous les hypothèses et avec les notations de
Méthodes simpliciales, lemme \ref{simplicial-lemma-section},
il existe une section $g : U \to V$ du morphisme $f$, et
le composé $g \circ f$ est homotope à l'identité de $V$.
En particulier, le morphisme $f$ est une équivalence d'homotopie.
\end{lemma}

\begin{proof}
Cela résulte immédiatement de Méthodes simpliciales, lemmes
\ref{simplicial-lemma-section} et
\ref{simplicial-lemma-trivial-kan-homotopy}.
\end{proof}

\begin{lemma}
\label{lemma-cosk-hom-deltak}
Soit $\mathcal{C}$ une catégorie admettant des sommes finies
et des limites finies. Soit $X$ un objet de $\mathcal{C}$.
Soit $k \geq 0$. Le morphisme canonique
$$
\Hom(\Delta[k], X)
\longrightarrow
\text{cosk}_1 \text{sk}_1 \Hom(\Delta[k], X)
$$
est un isomorphisme.
\end{lemma}

\begin{proof}
Pour tout objet simplicial $V$, on a
\begin{eqnarray*}
\Mor(V, \text{cosk}_1 \text{sk}_1 \Hom(\Delta[k], X))
& = &
\Mor(\text{sk}_1 V, \text{sk}_1 \Hom(\Delta[k], X)) \\
& = &
\Mor(i_{1!} \text{sk}_1 V, \Hom(\Delta[k], X)) \\
& = &
\Mor(i_{1!} \text{sk}_1 V \times \Delta[k], X)
\end{eqnarray*}
La première égalité résulte de l'adjonction entre $\text{sk}$ et $\text{cosk}$,
la deuxième de l'adjonction entre $i_{1!}$ et $\text{sk}_1$, et
la troisième de
Méthodes simpliciales, définition \ref{simplicial-definition-hom-from-simplicial-set},
où le dernier $X$ désigne l'objet simplicial constant de valeur $X$.
D'après Méthodes simpliciales, lemme \ref{simplicial-lemma-augmentation-howto},
un élément de cet ensemble ne dépend que des composantes de degrés $0$ et $1$
de $i_{1!} \text{sk}_1 V \times \Delta[k]$. Celles-ci
coïncident avec les composantes de degrés $0$ et $1$ de
$V \times \Delta[k]$; voir
Méthodes simpliciales, lemme \ref{simplicial-lemma-recovering-U-for-real}.
L'ensemble ci-dessus s'identifie donc à
$\Mor(V \times \Delta[k], X) = \Mor(V, \Hom(\Delta[k], X))$.
\end{proof}

\begin{lemma}
\label{lemma-cosk0-hom-deltak}
Soit $\mathcal{C}$ une catégorie. Soit $X$ un objet de $\mathcal{C}$
dont les puissances cartésiennes $X \times \ldots \times X$ existent.
Soit $k \geq 0$, et soit $C[k]$ l'objet défini dans
Méthodes simpliciales, exemple \ref{simplicial-example-simplex-cosimplicial-set}.
Avec les notations de
Méthodes simpliciales, lemme \ref{simplicial-lemma-morphism-into-product},
le morphisme canonique
$$
\Hom(C[k], X)_1
\longrightarrow
(\text{cosk}_0 \text{sk}_0 \Hom(C[k], X))_1
$$
s'identifie à l'application
$$
\prod\nolimits_{\alpha : [k] \to [1]} X
\longrightarrow
X \times X
$$
qui est la projection sur les facteurs correspondant aux applications
constantes $\alpha$.
\end{lemma}

\begin{proof}
C'est établi dans la démonstration de
Hyperrecouvrements, lemme \ref{hypercovering-lemma-covering}.
\end{proof}



\section{Résultats sur les schémas}
\label{section-devissage}

\noindent
Lemmes qui paraissent superflus.

\begin{lemma}
\label{lemma-stein-projective}
Soit $(R, \mathfrak m, \kappa)$ un anneau local.
Soit $X \subset \mathbf{P}^n_R$ un sous-schéma fermé.
Supposons que $R = \Gamma(X, \mathcal{O}_X)$. Alors la fibre spéciale
$X_\kappa$ est géométriquement connexe.
\end{lemma}

\begin{proof}
C'est un cas particulier de
Compléments sur les morphismes, théorème
\ref{more-morphisms-theorem-stein-factorization-general}.
\end{proof}

\begin{lemma}
\label{lemma-property-irreducible-higher-rank}
Soit $X$ un schéma noethérien.
Soit $Z_0 \subset X$ une partie fermée irréductible de point générique $\xi$.
Soit $\mathcal{P}$ une propriété des faisceaux cohérents sur $X$ telle que
\begin{enumerate}
\item Pour toute suite exacte courte de faisceaux cohérents, si deux
des trois termes possèdent la propriété $\mathcal{P}$, alors le troisième
la possède aussi.
\item Si une somme directe de deux faisceaux cohérents possède la propriété
$\mathcal{P}$, alors chacun des deux la possède.
\item Pour tout sous-schéma fermé intègre $Z \subset Z_0 \subset X$,
$Z \not = Z_0$, et tout faisceau quasi-cohérent d'idéaux
$\mathcal{I} \subset \mathcal{O}_Z$, la propriété $\mathcal{P}$ est satisfaite
par $(Z \to X)_*\mathcal{I}$.
\item Il existe un faisceau cohérent $\mathcal{G}$ sur $X$ tel que
\begin{enumerate}
\item $\text{Supp}(\mathcal{G}) = Z_0$,
\item $\mathcal{G}_\xi$ soit annulé par $\mathfrak m_\xi$, et
\item la propriété $\mathcal{P}$ soit satisfaite par $\mathcal{G}$.
\end{enumerate}
\end{enumerate}
Alors la propriété $\mathcal{P}$ est satisfaite par tout faisceau cohérent
$\mathcal{F}$ sur $X$ dont le support est contenu dans $Z_0$.
\end{lemma}

\begin{proof}
La démonstration est une variante de celle de
Cohomologie des schémas, lemme \ref{coherent-lemma-property-irreducible}.
Exactement comme dans cette démonstration, on voit que
tout faisceau cohérent dont le support est strictement contenu dans $Z_0$
possède la propriété $\mathcal{P}$.

\medskip\noindent
Considérons un faisceau cohérent $\mathcal{G}$ comme en (4).
D'après Cohomologie des schémas, lemme \ref{coherent-lemma-prepare-filter-irreducible},
il existe un faisceau d'idéaux $\mathcal{I}$ sur $Z_0$ et
une suite exacte courte
$$
0 \to
\left((Z_0 \to X)_*\mathcal{I}\right)^{\oplus r} \to
\mathcal{G} \to
\mathcal{Q} \to 0
$$
où le support de $\mathcal{Q}$ est strictement contenu dans $Z_0$.
En particulier, $r > 0$ et $\mathcal{I}$ est non nul,
car le support de $\mathcal{G}$ est égal à $Z_0$.
Puisque $\mathcal{Q}$ possède la propriété $\mathcal{P}$, on en conclut que
$\left((Z_0 \to X)_*\mathcal{I}\right)^{\oplus r}$
possède également la propriété $\mathcal{P}$.
D'après (2), on en déduit que la propriété $\mathcal{P}$ est satisfaite par
$(Z_0 \to X)_*\mathcal{I}$. En insérant ce fait au point voulu dans la démonstration de
Cohomologie des schémas, lemme \ref{coherent-lemma-property-irreducible},
on obtient le lemme.
Nous omettons quelques détails.
\end{proof}

\begin{lemma}
\label{lemma-property-higher-rank}
Soit $X$ un schéma noethérien.
Soit $\mathcal{P}$ une propriété des faisceaux cohérents sur $X$ telle que
\begin{enumerate}
\item Pour toute suite exacte courte de faisceaux cohérents, si deux
des trois termes possèdent la propriété $\mathcal{P}$, alors le troisième
la possède aussi.
\item Si une somme directe de deux faisceaux cohérents possède la propriété
$\mathcal{P}$, alors chacun des deux la possède.
\item Pour tout sous-schéma fermé intègre $Z \subset X$
de point générique $\xi$, il existe
un faisceau cohérent $\mathcal{G}$ tel que
\begin{enumerate}
\item $\text{Supp}(\mathcal{G}) = Z$,
\item $\mathcal{G}_\xi$ soit annulé par $\mathfrak m_\xi$, et
\item la propriété $\mathcal{P}$ soit satisfaite par $\mathcal{G}$.
\end{enumerate}
\end{enumerate}
Alors la propriété $\mathcal{P}$ est satisfaite par tout faisceau cohérent sur $X$.
\end{lemma}

\begin{proof}
Cela résulte du lemme \ref{lemma-property-irreducible-higher-rank}
exactement de la même manière que
Cohomologie des schémas, lemme \ref{coherent-lemma-property} résulte de
Cohomologie des schémas, lemme \ref{coherent-lemma-property-irreducible}.
\end{proof}

\begin{lemma}
\label{lemma-section-maps-back-into}
Soit $X$ un schéma.
Soit $\mathcal{L}$ un $\mathcal{O}_X$-module inversible.
Soit $s \in \Gamma(X, \mathcal{L})$ une section.
Soient $\mathcal{F}' \subset \mathcal{F}$ des
$\mathcal{O}_X$-modules quasi-cohérents. Supposons que
\begin{enumerate}
\item $X$ soit quasi-compact,
\item $\mathcal{F}$ soit de type fini, et
\item $\mathcal{F}'|_{X_s} = \mathcal{F}|_{X_s}$.
\end{enumerate}
Alors il existe un entier $n \geq 0$ tel que
la multiplication par $s^n$ sur $\mathcal{F}$ se factorise
par $\mathcal{F}'$.
\end{lemma}

\begin{proof}
Autrement dit, nous affirmons que
$s^n\mathcal{F} \subset
\mathcal{F}' \otimes_{\mathcal{O}_X} \mathcal{L}^{\otimes n}$
pour un certain $n \geq 0$. Autrement dit, nous affirmons que le morphisme
quotient $\mathcal{F} \to \mathcal{F}/\mathcal{F}'$ devient
nul après multiplication par une puissance de $s$.
Cela résulte de Propriétés, lemme
\ref{properties-lemma-section-maps-backwards}.
\end{proof}

\begin{lemma}
\label{lemma-bound-degree-in-nbhd-generic-point}
Soit $f : X \to Y$ un morphisme de schémas. Supposons que
\begin{enumerate}
\item $X$ et $Y$ soient des schémas intègres,
\item $f$ soit localement de type fini et dominant,
\item $f$ soit quasi-compact ou séparé,
\item $f$ soit génériquement fini, c'est-à-dire que l'une des conditions
(1) -- (5) de Morphismes, lemme \ref{morphisms-lemma-finite-degree}, soit satisfaite.
\end{enumerate}
Alors il existe un ouvert non vide $V \subset Y$ tel que
$f^{-1}(V) \to V$ soit fini localement libre de degré $\deg(X/Y)$.
En particulier, les degrés des fibres de $f^{-1}(V) \to V$
sont majorés par $\deg(X/Y)$.
\end{lemma}

\begin{proof}
On peut choisir $V$ de sorte que $f^{-1}(V) \to V$ soit fini.
Quitte à restreindre $V$, on peut alors supposer que $f^{-1}(V) \to V$
est plat et de présentation finie, d'après le théorème de platitude générique
(Morphismes, proposition \ref{morphisms-proposition-generic-flatness}).
Le morphisme est alors fini localement libre d'après
Morphismes, lemme \ref{morphisms-lemma-finite-flat}.
Comme $V$ est irréductible, le morphisme est de degré constant.
La dernière assertion en résulte avec
Morphismes, lemme \ref{morphisms-lemma-finite-locally-free-universally-bounded}.
\end{proof}





\section{Catégories dérivées des variétés}
\label{section-equiv}

\noindent
Un lemme qui figurait initialement dans le chapitre consacré aux catégories
dérivées des variétés, mais qui n'est désormais plus nécessaire.

\begin{lemma}
\label{lemma-preserves-Coh}
Soit $k$ un corps. Soit $X$ un schéma séparé de type fini sur $k$ qui est
régulier. Soit $F : D_{perf}(\mathcal{O}_X) \to D_{perf}(\mathcal{O}_X)$
un foncteur exact $k$-linéaire. Supposons que, pour tout
$\mathcal{O}_X$-module cohérent $\mathcal{F}$ tel que
$\dim(\text{Supp}(\mathcal{F})) = 0$, il existe un isomorphisme d'espaces
vectoriels sur $k$
$$
\Hom_X(\mathcal{F}, M) = \Hom_X(\mathcal{F}, F(M))
$$
fonctoriel en $M$ dans $D_{perf}(\mathcal{O}_X)$. Alors il existe un
automorphisme $f : X \to X$ sur $k$ qui induit l'identité sur l'espace
sous-jacent\footnote{Cela force souvent $f$ à être l'identité ;
voir Variétés, Lemme \ref{varieties-lemma-automorphism}.}
et un $\mathcal{O}_X$-module inversible $\mathcal{L}$ tels que $F$ et
$F'(M) = f^*M \otimes_{\mathcal{O}_X}^\mathbf{L} \mathcal{L}$ soient
apparentés.
\end{lemma}

\begin{proof}
Le Lemme \ref{equiv-lemma-duality-at-point} de Catégories dérivées des
variétés montre que, pour tout $\mathcal{O}_X$-module cohérent
$\mathcal{F}$ dont le support est un point fermé, il existe des
isomorphismes
$$
H^0(X, M \otimes^\mathbf{L}_{\mathcal{O}_X} \mathcal{F}) =
H^0(X, F(M) \otimes^\mathbf{L}_{\mathcal{O}_X} \mathcal{F})
$$
fonctoriels en $M$.

\medskip\noindent
Soit $x \in X$ un point fermé et appliquons ce qui précède à
$\mathcal{F} = \mathcal{O}_x$, le faisceau gratte-ciel de fibre
$\kappa(x)$ en $x$. On obtient
$$
\dim_{\kappa(x)}
\text{Tor}^{\mathcal{O}_{X, x}}_p(M_x, \kappa(x)) =
\dim_{\kappa(x)}
\text{Tor}^{\mathcal{O}_{X, x}}_p(F(M)_x, \kappa(x))
$$
pour tout $p \in \mathbf{Z}$. En particulier, si $H^i(M) = 0$ pour
$i > 0$, alors $H^i(F(M)) = 0$ pour $i > 0$, d'après le Lemme
\ref{equiv-lemma-orthogonal-point-sheaf} de Catégories dérivées des
variétés.

\medskip\noindent
Si $\mathcal{E}$ est localement libre de rang $r$, alors
$F(\mathcal{E})$ est localement libre de rang $r$. En effet, un complexe
parfait $K$ sur $\mathcal{O}_{X, x}$ tel que
$$
\dim_{\kappa(x)} \text{Tor}^{\mathcal{O}_{X, x}}_i(K, \kappa(x)) =
\left\{
\begin{matrix}
r & \text{si} & i = 0 \\
0 & \text{si} & i \not = 0
\end{matrix}
\right.
$$
est isomorphe à un module libre de rang $r$ placé en degré $0$. Voir, par
exemple, Algèbre, compléments, Lemme
\ref{more-algebra-lemma-lift-perfect-from-residue-field}.

\medskip\noindent
Si $M$ est supporté sur un sous-schéma fermé $Z \subset X$, alors $F(M)$
est également supporté sur $Z$. En effet,
$M \otimes_{\mathcal{O}_X}^\mathbf{L} \mathcal{O}_x = 0$ pour
$x \not \in Z$ ; il en va donc de même pour $F(M)$, et la conclusion
résulte du Lemme \ref{equiv-lemma-orthogonal-point-sheaf} de Catégories
dérivées des variétés.

\medskip\noindent
En particulier, $F(\mathcal{O}_x)$ est supporté sur $\{x\}$. Soit
$i \in \mathbf{Z}$ le plus petit entier tel que
$H^i(F(\mathcal{O}_x)) \not = 0$. On sait que $i \leq 0$. Si $i < 0$,
il existe alors un morphisme
$\mathcal{O}_x[-i] \to F(\mathcal{O}_x)$, ce qui contredit le fait que
tous les morphismes $\mathcal{O}_x[-i] \to \mathcal{O}_x$ sont nuls.
Ainsi $F(\mathcal{O}_x) = \mathcal{H}[0]$, où $\mathcal{H}$ est un
faisceau gratte-ciel en $x$.

\medskip\noindent
Soit $\mathcal{G}$ un $\mathcal{O}_X$-module cohérent tel que
$\dim(\text{Supp}(\mathcal{G})) = 0$. Il existe alors une filtration
$$
0 = \mathcal{G}_0 \subset \mathcal{G}_1 \subset \ldots \subset
\mathcal{G}_n = \mathcal{G}
$$
telle que, pour $n \geq i \geq 1$, le quotient
$\mathcal{G}_i/\mathcal{G}_{i - 1}$ soit isomorphe à
$\mathcal{O}_{x_i}$ pour un certain point fermé $x_i \in X$. On obtient
alors des triangles distingués
$$
F(\mathcal{G}_{i - 1}) \to F(\mathcal{G}_i) \to F(\mathcal{O}_{x_i})
$$
et une récurrence montre que $F(\mathcal{G}_i)$ est un faisceau cohérent
placé en degré $0$.

\medskip\noindent
Soit $\mathcal{G}$ un $\mathcal{O}_X$-module cohérent. Nous savons que
$H^i(F(\mathcal{G})) = 0$ pour $i > 0$. Supposons, en vue d'une
contradiction, que $H^i(F(\mathcal{G}))$ soit non nul pour un certain
$i < 0$. Choisissons $i$ minimal avec cette propriété, de sorte que nous
disposons d'un morphisme
$H^i(F(\mathcal{G}))[-i] \to F(\mathcal{G})$ dans
$D_{perf}(\mathcal{O}_X)$. Choisissons un point fermé $x \in X$ dans le
support de $H^i(F(\mathcal{G}))$. D'après Algèbre, compléments, Lemme
\ref{more-algebra-lemma-kollar-kovacs-pseudo-coherent}, il existe un
$n > 0$ tel que
$$
H^i(F(\mathcal{G}))_x \otimes_{\mathcal{O}_{X, x}}
\mathcal{O}_{X, x}/\mathfrak m_x^n
\longrightarrow
\text{Tor}^{\mathcal{O}_{X, x}}_{-i}(F(\mathcal{G})_x,
\mathcal{O}_{X, x}/\mathfrak m_x^n)
$$
soit non nul. Prenons ensuite $m \geq 1$ et considérons la suite exacte
courte
$$
0 \to \mathfrak m_x^m \mathcal{G} \to \mathcal{G} \to
\mathcal{G}/\mathfrak m_x^m\mathcal{G} \to 0
$$
D'après ce qui précède,
$F(\mathcal{G}/\mathfrak m_x^m\mathcal{G})$ est un faisceau placé en
degré $0$. Par conséquent,
$H^i(F(\mathfrak m_x^m \mathcal{G})) \to H^i(F(\mathcal{G}))$ est un
isomorphisme. Considérons le diagramme commutatif
$$
\xymatrix{
H^i(F(\mathfrak m_x^m\mathcal{G}))_x \otimes_{\mathcal{O}_{X, x}}
\mathcal{O}_{X, x}/\mathfrak m_x^n \ar[r] \ar[d] &
\text{Tor}^{\mathcal{O}_{X, x}}_{-i}(F(\mathfrak m_x^m\mathcal{G})_x,
\mathcal{O}_{X, x}/\mathfrak m_x^n) \ar[d] \\
H^i(F(\mathcal{G}))_x \otimes_{\mathcal{O}_{X, x}}
\mathcal{O}_{X, x}/\mathfrak m_x^n \ar[r] &
\text{Tor}^{\mathcal{O}_{X, x}}_{-i}(F(\mathcal{G})_x,
\mathcal{O}_{X, x}/\mathfrak m_x^n)
}
$$
Puisque la flèche verticale de gauche est un isomorphisme et que la flèche
horizontale du bas est non nulle, la flèche verticale de droite est non
nulle pour tout $m \geq 1$. D'autre part, le premier paragraphe de la
démonstration montre que cette flèche s'identifie à la flèche
$$
\text{Tor}^{\mathcal{O}_{X, x}}_{-i}(\mathfrak m_x^m\mathcal{G}_x,
\mathcal{O}_{X, x}/\mathfrak m_x^n)
\longrightarrow
\text{Tor}^{\mathcal{O}_{X, x}}_{-i}(\mathcal{G}_x,
\mathcal{O}_{X, x}/\mathfrak m_x^n)
$$
Cependant, cette flèche est nulle pour $m \gg n$ d'après Algèbre,
compléments, Lemme \ref{more-algebra-lemma-tor-annihilated}, ce qui donne
la contradiction cherchée.

\medskip\noindent
Nous savons donc que $F$ préserve les modules cohérents. D'après le Lemme
\ref{equiv-lemma-exact-functor-preserving-Coh} de Catégories dérivées des
variétés, $F$ est apparenté au foncteur de Fourier--Mukai $F'$ défini par
un $\mathcal{O}_{X \times X}$-module cohérent $\mathcal{K}$, plat sur
$X$ via $\text{pr}_1$ et fini sur $X$ via $\text{pr}_2$. Comme
$F(\mathcal{O}_X)$ est un $\mathcal{O}_X$-module inversible
$\mathcal{L}$ placé en degré $0$, on a
$$
\mathcal{L} \cong F(\mathcal{O}_X) \cong F'(\mathcal{O}_X) \cong
\text{pr}_{2, *}\mathcal{K}
$$
Le Lemme \ref{functors-lemma-pushforward-invertible-pre} de Foncteurs et
morphismes fournit donc un morphisme $s : X \to X \times X$ tel que
$\text{pr}_2 \circ s = \text{id}_X$ et $\mathcal{K} = s_*\mathcal{L}$.
Posons $f = \text{pr}_1 \circ s$. On a alors
\begin{align*}
F'(M)
& =
R\text{pr}_{2, *}(L\text{pr}_1^*M \otimes \mathcal{K}) \\
& =
R\text{pr}_{2, *}(L\text{pr}_1^*M \otimes s_*\mathcal{L}) \\
& =
R\text{pr}_{2, *}(Rs_*(Lf^*M \otimes \mathcal{L})) \\
& =
Lf^*M \otimes \mathcal{L}
\end{align*}
Dans la troisième étape, nous avons utilisé Catégories dérivées des
schémas, Lemme \ref{perfect-lemma-cohomology-base-change}.
Comme, pour tout point fermé $x \in X$, le module
$F(\mathcal{O}_x)$ est supporté en $x$, on voit que $f$ induit l'identité
sur l'espace sous-jacent de $X$. Il reste à montrer que $f$
est un isomorphisme ; c'est l'objet du paragraphe suivant.

\medskip\noindent
Soit $x \in X$ un point fermé. Pour $n \geq 1$, notons
$\mathcal{O}_{x, n}$ le faisceau gratte-ciel en $x$ de fibre
$\mathcal{O}_{X, x}/\mathfrak m_x^n$. Nous avons
$$
\Hom_X(\mathcal{O}_{x, m}, \mathcal{O}_{x, n}) \cong
\Hom_X(\mathcal{O}_{x, m}, F(\mathcal{O}_{x, n})) \cong
\Hom_X(\mathcal{O}_{x, m}, f^*\mathcal{O}_{x, n} \otimes \mathcal{L})
$$
fonctoriellement par rapport aux homomorphismes de $\mathcal{O}_X$-modules
entre les $\mathcal{O}_{x, n}$. (Le premier isomorphisme existe par
hypothèse, et le second parce que $F$ et $F'$ sont apparentés.) Pour
$m \geq n$, l'action sur $\mathcal{O}_{x, n}$ donne
$\mathcal{O}_{X, x}/\mathfrak m_x^n =
\Hom_X(\mathcal{O}_{x, m}, \mathcal{O}_{x, n})$ ; on en déduit que
$f^\sharp : \mathcal{O}_{X, x}/\mathfrak m_x^n \to
\mathcal{O}_{X, x}/\mathfrak m_x^n$ est bijectif pour tout $n$. Ainsi,
$f$ induit des isomorphismes sur les anneaux locaux complétés aux points
fermés et est donc étale (Morphismes étales, Lemme
\ref{etale-lemma-characterize-etale-completions}). En considérant les
points fermés, on voit que
$\Delta_f : X \to X \times_{f, X, f} X$ (qui est une immersion ouverte
puisque $f$ est étale) est bijective, donc est un isomorphisme. Par suite,
$f$ est un monomorphisme. Enfin, Descente, Lemme
\ref{descent-lemma-flat-surjective-quasi-compact-monomorphism-isomorphism}
montre que $f$ est un isomorphisme.
\end{proof}








\section{Représentabilité dans le cas propre et régulier}
\label{section-regular-proper}

\noindent
Cette section est obsolète : le Théorème
\ref{equiv-theorem-bondal-van-den-bergh} de Catégories dérivées des variétés
a été amélioré de sorte qu'il s'applique à tous les schémas propres sur un
corps (alors qu'auparavant nous ne l'avions démontré que pour les schémas
projectifs sur un corps).

\begin{lemma}
\label{lemma-trace-map}
\begin{reference}
La démonstration donnée ici reprend l'argument de
\cite[remarque 3.4]{MS}.
\end{reference}
Soit $f : X' \to X$ un morphisme propre et birationnel entre schémas
noethériens intègres, avec $X$ régulier. Le morphisme
$\mathcal{O}_X \to Rf_*\mathcal{O}_{X'}$ admet canoniquement une rétraction
dans $D(\mathcal{O}_X)$.
\end{lemma}

\begin{proof}
Posons $E = Rf_*\mathcal{O}_{X'}$ dans $D(\mathcal{O}_X)$.
Remarquons que $E$ appartient à
$D^b_{\textit{Coh}}(\mathcal{O}_X)$ d'après
Catégories dérivées des schémas, Lemme
\ref{perfect-lemma-direct-image-coherent}.
D'après Catégories dérivées des schémas, Lemme
\ref{perfect-lemma-perfect-on-regular}, $E$ est donc un objet parfait de
$D(\mathcal{O}_X)$. Comme $\mathcal{O}_{X'}$ est un faisceau d'algèbres,
le cup-produit relatif fournit le morphisme
$\mu : E \otimes_{\mathcal{O}_X}^\mathbf{L} E \to E$ de Cohomologie,
Remarque \ref{cohomology-remark-cup-product}.
Soit $\sigma : E \otimes E^\vee \to E^\vee \otimes E$ la contrainte de
commutativité de la catégorie monoïdale symétrique $D(\mathcal{O}_X)$
(Cohomologie, Lemme \ref{cohomology-lemma-symmetric-monoidal-derived}).
Notons $\eta : \mathcal{O}_X \to E \otimes E^\vee$ et
$\epsilon : E^\vee \otimes E \to \mathcal{O}_X$ les morphismes
construits dans Cohomologie, Exemple \ref{cohomology-example-dual-derived}.
Nous pouvons alors considérer le morphisme
$$
E \xrightarrow{\eta \otimes 1} E \otimes E^\vee \otimes E
\xrightarrow{\sigma \otimes 1} E^\vee \otimes E \otimes E
\xrightarrow{1 \otimes \mu} E^\vee \otimes E
\xrightarrow{\epsilon} \mathcal{O}_X
$$
Nous affirmons que ce morphisme est un inverse à gauche du morphisme figurant
dans l'énoncé du lemme. Pour le vérifier, il suffit de montrer que le composé
$\mathcal{O}_X \to \mathcal{O}_X$ est l'identité. Il suffit de le vérifier au
point générique de $X$ (ou sur un sous-schéma ouvert de $X$ au-dessus duquel
$f$ est un isomorphisme). Dans ce cas, $E = \mathcal{O}_X$ et $\mu$ est le
morphisme de multiplication usuel, d'où l'assertion.
\end{proof}

\begin{lemma}
\label{lemma-characterize-dbcoh-proper-regular}
Soit $X$ un schéma régulier, propre sur un corps $k$. Soit
$K \in \Ob(D_\QCoh(\mathcal{O}_X))$. Les assertions suivantes sont
équivalentes :
\begin{enumerate}
\item $K \in D^b_{\textit{Coh}}(\mathcal{O}_X) = D_{perf}(\mathcal{O}_X)$ ;
\item $\sum_{i \in \mathbf{Z}} \dim_k \Ext^i_X(E, K) < \infty$
pour tout objet parfait $E$ de $D(\mathcal{O}_X)$.
\end{enumerate}
\end{lemma}

\begin{proof}
L'égalité de (1) résulte de Catégories dérivées des schémas,
Lemme \ref{perfect-lemma-perfect-on-regular}.
L'implication (1) $\Rightarrow$ (2) résulte de
Catégories dérivées des variétés, Lemme \ref{equiv-lemma-finiteness}.
L'implication (2) $\Rightarrow$ (1) résulte de
Compléments sur les morphismes, Lemme
\ref{more-morphisms-lemma-characterize-relatively-perfect}.
\end{proof}

\begin{lemma}
\label{lemma-bondal-van-den-bergh}
Soit $X$ un schéma régulier, propre sur un corps $k$.
\begin{enumerate}
\item Soit $F : D_{perf}(\mathcal{O}_X)^{opp} \to \text{Vect}_k$
un foncteur cohomologique $k$-linéaire tel que
$$
\sum\nolimits_{n \in \mathbf{Z}} \dim_k F(E[n]) < \infty
$$
pour tout $E \in D_{perf}(\mathcal{O}_X)$. Alors $F$ est isomorphe à un
foncteur de la forme $E \mapsto \Hom_X(E, K)$, où
$K \in D_{perf}(\mathcal{O}_X)$.
\item Soit $G : D_{perf}(\mathcal{O}_X) \to \text{Vect}_k$
un foncteur homologique $k$-linéaire tel que
$$
\sum\nolimits_{n \in \mathbf{Z}} \dim_k G(E[n]) < \infty
$$
pour tout $E \in D_{perf}(\mathcal{O}_X)$. Alors $G$ est isomorphe à un
foncteur de la forme $E \mapsto \Hom_X(K, E)$, où
$K \in D_{perf}(\mathcal{O}_X)$.
\end{enumerate}
\end{lemma}

\begin{proof}
Cela résulte de Catégories dérivées des variétés, Théorème
\ref{equiv-theorem-bondal-van-den-bergh}, et du
Lemme \ref{equiv-lemma-homological-representable}.
Nous donnons également une autre démonstration ci-dessous.

\medskip\noindent
Démonstration de (1). La catégorie dérivée $D_\QCoh(\mathcal{O}_X)$ admet les
sommes directes, est compactement engendrée, et
$D_{perf}(\mathcal{O}_X)$ en est la sous-catégorie pleine des objets compacts ;
voir Catégories dérivées des schémas, Lemme
\ref{perfect-lemma-quasi-coherence-direct-sums},
Théorème \ref{perfect-theorem-bondal-van-den-Bergh}, et
Proposition \ref{perfect-proposition-compact-is-perfect}.
Le lemme de Catégories dérivées des variétés
\ref{equiv-lemma-van-den-bergh} permet de supposer que
$F(E) = \Hom_X(E, K)$ pour un certain
$K \in \Ob(D_\QCoh(\mathcal{O}_X))$.
$K$ appartient alors à $D^b_{\textit{Coh}}(\mathcal{O}_X)$ d'après le
Lemme \ref{lemma-characterize-dbcoh-proper-regular}.

\medskip\noindent
Démonstration de (2). Considérons le foncteur contravariant
$E \mapsto E^\vee$ sur $D_{perf}(\mathcal{O}_X)$ ; voir
Cohomologie, Lemme \ref{cohomology-lemma-dual-perfect-complex}.
C'est une anti-équivalence exacte de $D_{perf}(\mathcal{O}_X)$ sur elle-même.
On peut donc appliquer (1) au foncteur $F(E) = G(E^\vee)$ : il existe
$K \in D_{perf}(\mathcal{O}_X)$ tel que
$G(E^\vee) = \Hom_X(E, K)$.
Il s'ensuit que $G(E) = \Hom_X(E^\vee, K) = \Hom_X(K^\vee, E)$ ;
ainsi, $K^\vee$ convient.
\end{proof}








\section{Foncteur des quotients}
\label{section-quotients}

\begin{lemma}
\label{lemma-factors-through-quotient}
Soit $S = \Spec(R)$ un sch\'ema affine. Soit $X$ un espace alg\'ebrique sur
$S$. Soient $q_i : \mathcal{F} \to \mathcal{Q}_i$, $i = 1, 2$,
des morphismes surjectifs de $\mathcal{O}_X$-modules quasi-coh\'erents.
Supposons $\mathcal{Q}_1$ plat sur $S$. Soit $T \to S$ un morphisme quasi-compact
de sch\'emas tel qu'il existe une factorisation
$$
\xymatrix{
& \mathcal{F}_T \ar[rd]^{q_{2, T}} \ar[ld]_{q_{1, T}} \\
\mathcal{Q}_{1, T} & & \mathcal{Q}_{2, T} \ar@{..>}[ll]
}
$$
Il existe alors un sous-sch\'ema ferm\'e $Z \subset S$ tel que
(a) $T \to S$ se factorise par $Z$ et (b)
$q_{1, Z}$ se factorise par $q_{2, Z}$.
Si $\Ker(q_2)$ est un $\mathcal{O}_X$-module de type fini et si $X$
est quasi-compact, on peut choisir $Z \to S$ de pr\'esentation finie.
\end{lemma}

\begin{proof}
On applique Platitude sur les espaces, Lemme \ref{spaces-flat-lemma-F-zero-somewhat-closed},
au morphisme $\Ker(q_2) \to \mathcal{Q}_1$.
\end{proof}







\section{Espaces et recouvrements fpqc}
\label{section-fpqc}

\noindent
Le contenu de cette section a \'et\'e rendu obsol\`ete par l'argument de Gabber, qui montre que
les espaces alg\'ebriques v\'erifient la condition de faisceau pour les
recouvrements fpqc. Voir
Propri\'et\'es des espaces, Section \ref{spaces-properties-section-fpqc}.

\begin{lemma}
\label{lemma-separated-fpqc}
Soit $S$ un sch\'ema. Soit $X$ un espace alg\'ebrique sur $S$.
Soit $\{f_i : T_i \to T\}_{i \in I}$ un recouvrement fpqc de sch\'emas sur $S$.
Alors l'application
$$
\Mor_S(T, X)
\longrightarrow
\prod\nolimits_{i \in I} \Mor_S(T_i, X)
$$
est injective.
\end{lemma}

\begin{proof}
C'est une cons\'equence imm\'ediate de la Proposition
\ref{spaces-properties-proposition-sheaf-fpqc}
de Propri\'et\'es des espaces.
\end{proof}

\begin{lemma}
\label{lemma-sheaf-fpqc-open-covering}
Soit $S$ un sch\'ema. Soit $X$ un espace alg\'ebrique sur $S$.
Soit $X = \bigcup_{j \in J} X_j$ un recouvrement de Zariski; voir
Espaces, D\'efinition \ref{spaces-definition-Zariski-open-covering}.
Si chacun des $X_j$ v\'erifie la condition de faisceau pour la topologie fpqc,
alors $X$ v\'erifie la condition de faisceau pour la topologie fpqc.
\end{lemma}

\begin{proof}
Cela r\'esulte du fait que tous les espaces alg\'ebriques v\'erifient la condition
de faisceau pour la topologie fpqc; voir
Propri\'et\'es des espaces, Proposition
\ref{spaces-properties-proposition-sheaf-fpqc}.
\end{proof}

\begin{lemma}
\label{lemma-sheaf-fpqc-quasi-separated}
Soit $S$ un sch\'ema. Soit $X$ un espace alg\'ebrique sur $S$.
Si $X$ est localement quasi-s\'epar\'e sur $S$ pour la topologie de Zariski,
alors $X$ v\'erifie la condition de faisceau pour la topologie fpqc.
\end{lemma}

\begin{proof}
C'est une cons\'equence imm\'ediate de la Proposition g\'en\'erale
\ref{spaces-properties-proposition-sheaf-fpqc}
de Propri\'et\'es des espaces.
\end{proof}

\begin{remark}
\label{remark-proof-works-when}
Cette remarque examinait autrefois dans quelle mesure la d\'emonstration originelle du
Lemme \ref{lemma-sheaf-fpqc-quasi-separated} (dat\'ee du 18 d\'ecembre 2009)
se g\'en\'eralise.
\end{remark}





\section{Espaces algébriques très raisonnables}
\label{section-very-reasonable}

\noindent
Éléments quelque peu obsolètes.

\begin{lemma}
\label{lemma-reasonable-kolmogorov}
Soit $S$ un schéma.
Soit $X$ un espace algébrique raisonnable au-dessus de $S$.
Alors $|X|$ est un espace de Kolmogorov (voir
Topologie, Définition \ref{topology-definition-generic-point}).
\end{lemma}

\begin{proof}
Cela résulte des définitions et du
Lemme \ref{decent-spaces-lemma-kolmogorov} d'Espaces décents.
\end{proof}

\noindent
Dans la suite de cette section, nous formulons quelques remarques sur les
espaces algébriques très raisonnables. S'il existe un schéma $U$ et un
morphisme surjectif, étale et quasi-compact
$U \to X$, alors $X$ est très raisonnable ; voir
Lemme \ref{decent-spaces-lemma-characterize-very-reasonable} d'Espaces décents.

\begin{lemma}
\label{lemma-scheme-very-reasonable}
Tout schéma est très raisonnable.
\end{lemma}

\begin{proof}
Cela tient à ce que le morphisme identité est étale, quasi-compact
et surjectif.
\end{proof}

\begin{lemma}
\label{lemma-very-reasonable-Zariski-local}
Soit $S$ un schéma.
Soit $X$ un espace algébrique au-dessus de $S$.
S'il existe un recouvrement ouvert de Zariski $X = \bigcup X_i$ tel que
chaque $X_i$ soit très raisonnable, alors $X$ est très raisonnable.
\end{lemma}

\begin{proof}
C'est le cas $(\epsilon)$ du
Lemme \ref{decent-spaces-lemma-properties-local} d'Espaces décents.
\end{proof}

\begin{lemma}
\label{lemma-quasi-separated-very-reasonable}
Tout espace algébrique localement quasi-séparé pour la topologie de Zariski est très raisonnable.
En particulier, tout espace algébrique quasi-séparé est très raisonnable.
\end{lemma}

\begin{proof}
C'est l'une des implications du
Lemme \ref{decent-spaces-lemma-bounded-fibres} d'Espaces décents.
\end{proof}

\begin{lemma}
\label{lemma-representable-very-reasonable}
Soit $S$ un schéma.
Soient $X$ et $Y$ des espaces algébriques au-dessus de $S$.
Soit $Y \to X$ un morphisme représentable.
Si $X$ est très raisonnable, alors $Y$ l'est aussi.
\end{lemma}

\begin{proof}
C'est le cas $(\epsilon)$ du
Lemme \ref{decent-spaces-lemma-representable-properties} d'Espaces décents.
\end{proof}

\begin{remark}
\label{remark-very-reasonable-Zariski-locally-quasi-separated}
Les espaces algébriques très raisonnables forment une classe strictement plus
grande que celle des espaces algébriques localement quasi-séparés pour la
topologie de Zariski. Considérons un espace algébrique de la forme
$X = [U/G]$ (voir
Espaces, Définition \ref{spaces-definition-quotient}),
où $G$ est un groupe fini agissant sans points fixes sur un schéma
$U$ non quasi-séparé. En effet, dans ce cas,
$U \times_X U = U \times G$, et les deux projections vers $U$ sont
manifestement quasi-compactes ; ainsi, $X$ est très raisonnable. En revanche,
la diagonale $U \times_X U \to U \times U$ n'est pas quasi-compacte ; cet
espace algébrique n'est donc pas quasi-séparé. Prenons maintenant pour $U$
l'espace affine de dimension infinie sur un corps $k$ de caractéristique
$\not = 2$, dont l'origine est dédoublée ; voir
Schémas, Exemple \ref{schemes-example-not-quasi-separated}.
Notons $0_1, 0_2$ les deux origines de $U$. Soit $G = \{+1, -1\}$ et
faisons agir $-1$ par multiplication par $-1$ sur toutes les coordonnées,
tout en échangeant $0_1$ et $0_2$. Alors $[U/G]$ est très raisonnable, mais
n'est pas localement quasi-séparé pour la topologie de Zariski (détails omis).
\end{remark}

\noindent
Avertissement : le lemme suivant doit être utilisé avec prudence, car les
schémas $U_i$ qui y interviennent ne sont pas nécessairement séparés, ni même
quasi-séparés.

\begin{lemma}
\label{lemma-very-reasonable-quasi-compact-pieces}
Soit $S$ un schéma.
Soit $X$ un espace algébrique très raisonnable au-dessus de $S$.
Il existe une famille de schémas
$U_i$ et des morphismes $U_i \to X$ tels que
\begin{enumerate}
\item chaque $U_i$ soit un schéma quasi-compact,
\item chaque $U_i \to X$ soit étale,
\item les deux projections $U_i \times_X U_i \to U_i$ soient quasi-compactes, et
\item le morphisme $\coprod U_i \to X$ soit surjectif (et étale).
\end{enumerate}
\end{lemma}

\begin{proof}
La Définition \ref{decent-spaces-definition-very-reasonable} d'Espaces décents
assure l'existence de $U_i \to X$ satisfaisant (2), (3) et (4).
Fixons $i$, posons $R_i = U_i \times_X U_i$ et notons $s, t : R_i \to U_i$
les projections.
Pour tout ouvert affine $W \subset U_i$, l'ouvert $W' = t(s^{-1}(W)) \subset U_i$
est quasi-compact et invariant par $R_i$ (voir
Groupoïdes, Lemme \ref{groupoids-lemma-constructing-invariant-opens}).
Ainsi, $W'$ est un schéma quasi-compact, $W' \to X$ est étale et
$W' \times_X W' = s^{-1}(W') = t^{-1}(W')$ ; les deux projections
$W' \times_X W' \to W'$ sont donc quasi-compactes. La famille des
$W' \to X$, où $W \subset U_i$ parcourt les membres de recouvrements ouverts
affines des $U_i$, possède par conséquent les propriétés voulues.
\end{proof}




\section{Lemmes obsolètes sur les espaces algébriques}
\label{section-obsolete-on-spaces}

\noindent
Lemmes qui semblent superflus ou ne sont plus utilisés dans le texte.

\begin{lemma}
\label{lemma-vanishing-surjective}
Dans la Situation \ref{spaces-cohomology-situation-vanishing} de Cohomologie des espaces,
le morphisme $p : X \to \Spec(A)$ est surjectif.
\end{lemma}

\begin{proof}
Ce lemme intervenait à l'origine dans la démonstration de la
Proposition \ref{spaces-cohomology-proposition-vanishing-affine} de
Cohomologie des espaces, mais il en est désormais une conséquence.
\end{proof}

\begin{lemma}
\label{lemma-vanishing-universally-closed}
Dans la Situation \ref{spaces-cohomology-situation-vanishing} de Cohomologie des espaces,
le morphisme $p : X \to \Spec(A)$ est universellement fermé.
\end{lemma}

\begin{proof}
Ce lemme intervenait à l'origine dans la démonstration de la
Proposition \ref{spaces-cohomology-proposition-vanishing-affine} de
Cohomologie des espaces, mais il en est désormais une conséquence.
\end{proof}

\begin{remark}
\label{remark-equation-first}
Ce tag renvoyait autrefois à une équation de la démonstration du
Lemme \ref{formal-spaces-lemma-iff-adic-star} d'Espaces formels.
\end{remark}

\begin{remark}
\label{remark-equation-second}
Ce tag renvoyait autrefois à une équation de la démonstration du
Lemme \ref{formal-spaces-lemma-iff-adic-star} d'Espaces formels.
\end{remark}



\section{Lemmes obsolètes sur les champs algébriques}
\label{section-obsolete-on-stacks}

\noindent
Lemmes qui semblent superflus ou ne sont plus utilisés dans le texte.

\begin{lemma}
\label{lemma-infinite-sequence}
Soit $S$ un schéma localement noethérien. Soit $\mathcal{X}$ une catégorie
fibrée en groupoïdes au-dessus de $(\Sch/S)_{fppf}$ vérifiant (RS*).
Soit $x$ un objet de
$\mathcal{X}$ au-dessus d'un schéma affine $U$ de type fini sur $S$.
Soient $u_n \in U$, $n \geq 1$, des points de type fini deux à deux distincts
tels que $x$ ne soit versel en $u_n$ pour tout $n$. Après avoir remplacé
$u_n$ par une sous-suite, il existe des morphismes
$$
x \to x_1 \to x_2 \to \ldots
\quad\text{dans }\mathcal{X}\text{, au-dessus de}\quad
U \to U_1 \to U_2 \to \ldots
$$
sur $S$ tels que
\begin{enumerate}
\item pour tout $n$, le morphisme $U \to U_n$ est un épaississement
du premier ordre,
\item pour tout $n$, on a une suite exacte courte
$$
0 \to \kappa(u_n) \to \mathcal{O}_{U_n} \to \mathcal{O}_{U_{n - 1}} \to 0
$$
où $U_0 = U$ lorsque $n = 1$,
\item pour tout $n$, il n'existe {\bf aucun} couple $(W, \alpha)$
formé d'un voisinage ouvert $W \subset U_n$ de $u_n$
et d'un morphisme $\alpha : x_n|_W \to x$
tel que le composé
$$
x|_{U \cap W} \xrightarrow{\text{restriction de }x \to x_n}
x_n|_W \xrightarrow{\alpha} x
$$
soit le morphisme canonique $x|_{U \cap W} \to x$.
\end{enumerate}
\end{lemma}

\begin{proof}
À l'origine, ce lemme intervenait dans la preuve d'un critère
d'ouverture de la versalité
(Artin's Axioms, lemme \ref{artin-lemma-SGE-implies-openness-versality}),
mais il a été remplacé par Artin's Axioms,
lemme \ref{artin-lemma-infinite-sequence-pre}, dont il découle immédiatement.
En effet, quitte à remplacer $u_n$, $n \geq 1$, par une sous-suite,
nous pouvons supposer qu'il n'existe aucune relation de spécialisation
entre ces points ; voir Properties,
lemme \ref{properties-lemma-thin-infinite-sequence}.
Il suffit alors d'appliquer Artin's Axioms,
lemme \ref{artin-lemma-infinite-sequence-pre} pour achever la preuve.
\end{proof}




\section{Variantes du complexe cotangent pour les schémas}
\label{section-cotangent-schemes-variant}

\noindent
Cette section donne une autre construction du complexe cotangent
d'un morphisme de schémas. Elle figure actuellement dans le chapitre
obsolète, car, pour les applications, la version plus élémentaire étudiée dans
Complexe cotangent, Section \ref{cotangent-section-cotangent-schemes-variant}
suffit.

\medskip\noindent
Soit $f : X \to Y$ un morphisme de schémas. Soit $\mathcal{C}_{X/Y}$ la
catégorie dont les objets sont les diagrammes commutatifs
\begin{equation}
\label{equation-object}
\vcenter{
\xymatrix{
X \ar[d] & U \ar[l] \ar[d] \ar[r]_i & A \ar[ld] \\
Y & V \ar[l]
}
}
\end{equation}
de schémas tels que
\begin{enumerate}
\item $U$ soit un sous-schéma ouvert de $X$,
\item $V$ soit un sous-schéma ouvert de $Y$, et
\item il existe un isomorphisme $A = V \times \Spec(P)$ au-dessus de $V$,
où $P$ est une algèbre de polynômes sur $\mathbf{Z}$ (en un ensemble
quelconque de variables).
\end{enumerate}
Autrement dit, $A$ est un espace affine (de dimension éventuellement infinie)
sur $V$. Les morphismes sont donnés par des diagrammes commutatifs.

\medskip\noindent
{\bf Notation.} Un objet de $\mathcal{C}_{X/Y}$, c'est-à-dire un diagramme
(\ref{equation-object}), sera souvent noté $U \to A$, étant entendu que
(a) $U$ est un sous-schéma ouvert de $X$, (b) $U \to A$ est un morphisme
au-dessus de $Y$, (c) l'image du morphisme structural $A \to Y$ est un ouvert
$V \subset Y$, et (d) le morphisme $A \to V$ présente sa source comme un
espace affine sur sa base.
Nous écrirons $U \to A/V$ pour indiquer que $V \subset Y$ est l'image de
$A \to Y$.
Rappelons que $X_{Zar}$ désigne le petit site de Zariski de $X$.
Il existe des foncteurs d'oubli
$$
\mathcal{C}_{X/Y} \to X_{Zar},\ (U \to A) \mapsto U
\quad\text{et}\quad
\mathcal{C}_{X/Y} \to Y_{Zar},\ (U \to A/V) \mapsto V.
$$

\begin{lemma}
\label{lemma-category-fibred}
Soit $X \to Y$ un morphisme de schémas.
\begin{enumerate}
\item La catégorie $\mathcal{C}_{X/Y}$ est fibrée sur $X_{Zar}$.
\item La catégorie $\mathcal{C}_{X/Y}$ est fibrée sur $Y_{Zar}$.
\item La catégorie $\mathcal{C}_{X/Y}$ est fibrée sur la
catégorie des couples $(U, V)$, où $U \subset X$ et $V \subset Y$ sont
ouverts et où $f(U) \subset V$.
\end{enumerate}
\end{lemma}

\begin{proof}
Pour (1). Étant donnés un objet $U \to A$ de $\mathcal{C}_{X/Y}$ et un
morphisme $U' \to U$ de $X_{Zar}$, considérons l'objet $i' : U' \to A$ de
$\mathcal{C}_{X/Y}$, où $i'$ est le composé de $i$ et de $U' \to U$.
Le morphisme $(U' \to A) \to (U \to A)$ de $\mathcal{C}_{X/Y}$
est fortement cartésien au-dessus de $X_{Zar}$.

\medskip\noindent
Pour (2). Étant donnés un objet $U \to A/V$ et un morphisme $V' \to V$,
posons $U' = U \cap f^{-1}(V')$ et $A' = V' \times_V A$; on obtient ainsi
un morphisme fortement cartésien $(U' \to A') \to (U \to A)$ au-dessus de
$V' \to V$.

\medskip\noindent
Pour (3). Notons $(X/Y)_{Zar}$ la catégorie de (3). Étant donnés
$U \to A/V$ et un morphisme $(U', V') \to (U, V)$ dans $(X/Y)_{Zar}$,
considérons $A' = V' \times_V A$. Alors le morphisme
$(U' \to A'/V') \to (U \to A/V)$ est fortement cartésien dans
$\mathcal{C}_{X/Y}$ au-dessus de $(X/Y)_{Zar}$.
\end{proof}

\noindent
Nous obtenons une topologie $\tau_X$ sur $\mathcal{C}_{X/Y}$ en
prenant la topologie induite par celle de $X_{Zar}$ (voir
Champs, Section \ref{stacks-section-topology}). Sauf mention contraire,
c'est cette topologie sur $\mathcal{C}_{X/Y}$ que nous considérerons.
Précisément, une famille de morphismes
$\{(U_i \to A_i) \to (U \to A)\}$ est couvrante dans
$\mathcal{C}_{X/Y}$ si et seulement si
\begin{enumerate}
\item $U = \bigcup U_i$, et
\item $A_i = A$ pour tout $i$.
\end{enumerate}
On obtient la même catégorie de faisceaux si l'on autorise
$A_i \cong A$ dans (2). Le foncteur d'oubli vers $X_{Zar}$ définit un
morphisme de topos
$\pi : \Sh(\mathcal{C}_{X/Y}) \to \Sh(X_{Zar})$.

\medskip\noindent
Le site $\mathcal{C}_{X/Y}$ est muni de plusieurs faisceaux d'anneaux.
\begin{enumerate}
\item Le faisceau $\mathcal{O}$ défini par la règle
$(U \to A) \mapsto \mathcal{O}(A)$.
\item Le faisceau $\underline{\mathcal{O}}_X = \pi^{-1}\mathcal{O}_X$ défini par
la règle $(U \to A) \mapsto \mathcal{O}(U)$.
\item Le faisceau $\underline{\mathcal{O}}_Y$ défini par la règle
$(U \to A/V) \mapsto \mathcal{O}(V)$.
\end{enumerate}
Nous en déduisons des morphismes de topos annelés
\begin{equation}
\label{equation-pi-schemes}
\vcenter{
\xymatrix{
(\Sh(\mathcal{C}_{X/Y}), \underline{\mathcal{O}}_X) \ar[r]_i \ar[d]_\pi &
(\Sh(\mathcal{C}_{X/Y}), \mathcal{O}) \\
(\Sh(X_{Zar}), \mathcal{O}_X)
}
}
\end{equation}
Le morphisme $i$ est l'identité sur les topos sous-jacents, et
$i^\sharp : \mathcal{O} \to \underline{\mathcal{O}}_X$
est le morphisme évident.
Le morphisme $\pi$ est un cas particulier de Cohomologie sur les sites, Situation
\ref{sites-cohomology-situation-fibred-category}.
Les foncteurs dérivés
$
Li^* : D(\mathcal{O}) \longrightarrow D(\underline{\mathcal{O}}_X)
$,
adjoint à gauche de
$Ri_* = i_* : D(\underline{\mathcal{O}}_X) \to D(\mathcal{O})$,
et
$
L\pi_! : D(\underline{\mathcal{O}}_X) \longrightarrow D(\mathcal{O}_X)
$,
adjoint à gauche de
$\pi^* = \pi^{-1} : D(\mathcal{O}_X) \to D(\underline{\mathcal{O}}_X)$,
joueront un rôle important dans la suite.

\begin{remark}
\label{remark-different-topologies}
Nous obtenons une deuxième topologie $\tau_Y$ sur $\mathcal{C}_{X/Y}$
en prenant la topologie induite par celle de $Y_{Zar}$.
Il existe une troisième topologie $\tau_{X \to Y}$, pour laquelle une famille
de morphismes $\{(U_i \to A_i) \to (U \to A)\}$ est couvrante si et seulement
si $U = \bigcup U_i$, $V = \bigcup V_i$ et $A_i \cong V_i \times_V A$.
C'est la topologie induite par celle du site $(X/Y)_{Zar}$, dont la catégorie
sous-jacente est la catégorie des couples $(U, V)$ du
Lemme \ref{lemma-category-fibred}, partie (3). Les familles couvrantes de
$(X/Y)_{Zar}$ sont les familles $\{(U_i, V_i) \to (U, V)\}$ telles que
$U = \bigcup U_i$ et $V = \bigcup V_i$. Il existe des morphismes de topos
$$
\xymatrix{
\Sh(\mathcal{C}_{X/Y})
= \Sh(\mathcal{C}_{X/Y}, \tau_X) &
\Sh(\mathcal{C}_{X/Y}, \tau_{X \to Y}) \ar[l] \ar[r] &
\Sh(\mathcal{C}_{X/Y}, \tau_Y)
}
$$
(rappelons que $\tau_X$ est notre topologie «~par défaut~»). Les foncteurs
d'image inverse associés à ces flèches sont donnés par le passage au faisceau
associé, tandis que l'image directe est l'identité sur les préfaisceaux
sous-jacents. Le diagramme de topos
$$
\xymatrix{
\Sh(X_{Zar}) \ar[d]^f & \Sh(\mathcal{C}_{X/Y}) \ar[l]^\pi &
\Sh(\mathcal{C}_{X/Y}, \tau_{X \to Y}) \ar[l] \ar[d] \\
\Sh(Y_{Zar}) & & \Sh(\mathcal{C}_{X/Y}, \tau_Y) \ar[ll]
}
$$
n'est {\bf pas} commutatif. En effet, l'image inverse d'un faisceau abélien
non nul sur $Y$ est un faisceau abélien non nul sur
$(\mathcal{C}_{X/Y}, \tau_{X \to Y})$, mais il existe certainement des
exemples où son image inverse sur $X$ est nulle. Notons que tout préfaisceau
$\mathcal{F}$ sur $Y_{Zar}$ définit un faisceau $\underline{\mathcal{F}}$ sur
$\mathcal{C}_{X/Y}$ par la règle qui associe à $(U \to A/V)$ l'ensemble
$\mathcal{F}(V)$. Même lorsque $\mathcal{F}$ est un faisceau, on n'a pas en
général $\underline{\mathcal{F}} = \pi^{-1}f^{-1}\mathcal{F}$.
Cela tient à la non-commutativité du diagramme ci-dessus, car
$\underline{\mathcal{F}}$ est l'image directe de l'image inverse de
$\mathcal{F}$ sur $\Sh(\mathcal{C}_{X/Y}, \tau_{X \to Y})$ par les flèches
horizontale inférieure et verticale droite. Le faisceau
$\underline{\mathcal{O}}_Y$ en est un exemple.
En revanche, il existe un morphisme
$\underline{\mathcal{F}} \to \pi^{-1}f^{-1}\mathcal{F}$
qui, comme nous le verrons plus loin, devient un isomorphisme après application
de $\pi_!$.
\end{remark}





\section{D\'eformations et obstructions des modules plats}
\label{section-modules}

\noindent
Dans cette section, nous esquissons la construction d'une th\'eorie des
d\'eformations du champ des faisceaux coh\'erents sur un espace alg\'ebrique
$X$ quelconque au-dessus d'un anneau $\Lambda$. Ce d\'eveloppement est
obsol\`ete en raison de l'expos\'e am\'elior\'e de
Quot, Section \ref{quot-section-not-flat}.

\medskip\noindent
Nous nous placerons dans le cadre suivant. On se donne
\begin{enumerate}
\item un anneau $\Lambda$,
\item un espace alg\'ebrique $X$ sur $\Lambda$,
\item une $\Lambda$-alg\`ebre $A$; posons
$X_A = X \times_{\Spec(\Lambda)} \Spec(A)$, et
\item un $\mathcal{O}_{X_A}$-module $\mathcal{F}$ de pr\'esentation finie,
plat sur $A$.
\end{enumerate}
Dans cette situation, nous consid\'ererons toutes les surjections possibles
$$
0 \to I \to A' \to A \to 0
$$
o\`u $A'$ est une $\Lambda$-alg\`ebre dont le noyau $I$ est un id\'eal de carr\'e
nul dans $A'$. Pour tout $A'$, on obtient un \'epaississement du premier
ordre $X_A \to X_{A'}$ d'espaces alg\'ebriques au-dessus de
$\Spec(\Lambda)$. Pour chacun de ces \'epaississements, nous consid\'erons le
probl\`eme du rel\`evement de $\mathcal{F}$ en un module $\mathcal{F}'$ de
pr\'esentation finie sur $X_{A'}$, plat sur $A'$. Nous voudrions retrouver
dans ce cadre les r\'esultats de Th\'eorie des d\'eformations, Lemme
\ref{defos-lemma-flat-ringed-topoi}.

\medskip\noindent
Plus pr\'ecis\'ement, notons $\textit{Lift}(\mathcal{F}, A')$ la cat\'egorie
des couples $(\mathcal{F}', \alpha)$, o\`u $\mathcal{F}'$ est un module de
pr\'esentation finie sur $X_{A'}$, plat sur $A'$, et o\`u
$\alpha : \mathcal{F}'|_{X_A} \to \mathcal{F}$ est un isomorphisme.
Les morphismes $(\mathcal{F}'_1, \alpha_1) \to (\mathcal{F}'_2, \alpha_2)$
sont les isomorphismes $\mathcal{F}'_1 \to \mathcal{F}'_2$ compatibles avec
$\alpha_1$ et $\alpha_2$.
L'ensemble des classes d'isomorphisme de $\textit{Lift}(\mathcal{F}, A')$
est not\'e $\text{Lift}(\mathcal{F}, A')$.

\medskip\noindent
Soit $\mathcal{G}$ un faisceau de $\mathcal{O}_X \otimes_\Lambda A$-modules
sur $X_\etale$, plat sur $A$. Nous introduisons la cat\'egorie
$\textit{Lift}(\mathcal{G}, A')$ des couples
$(\mathcal{G}', \beta)$, o\`u $\mathcal{G}'$ est un faisceau de
$\mathcal{O}_X \otimes_\Lambda A'$-modules plat sur $A'$ et $\beta$
est un isomorphisme $\mathcal{G}' \otimes_{A'} A \to \mathcal{G}$.

\begin{lemma}
\label{lemma-equivalence}
Notations et hypoth\`eses comme ci-dessus. Soit $p : X_A \to X$ la projection.
Pour tout $A'$, notons $p' : X_{A'} \to X$ la projection. Le foncteur $p'_*$
induit une \'equivalence de cat\'egories entre
\begin{enumerate}
\item la cat\'egorie $\textit{Lift}(\mathcal{F}, A')$, et
\item la cat\'egorie $\textit{Lift}(p_*\mathcal{F}, A')$.
\end{enumerate}
\end{lemma}

\begin{proof}
FIXME.
\end{proof}

\noindent
Soit $\mathcal{H}$ un faisceau de $\mathcal{O} \otimes_\Lambda A$-modules
sur $\mathcal{C}_{X/\Lambda}$, plat sur $A$. Nous introduisons la cat\'egorie
$\textit{Lift}_\mathcal{O}(\mathcal{H}, A')$
dont les objets sont les couples $(\mathcal{H}', \gamma)$, o\`u $\mathcal{H}'$
est un faisceau de $\mathcal{O} \otimes_\Lambda A'$-modules plat sur $A'$,
et o\`u $\gamma : \mathcal{H}' \otimes_{A'} A \to \mathcal{H}$
est un isomorphisme de $\mathcal{O} \otimes_\Lambda A$-modules.

\medskip\noindent
Soit $\mathcal{G}$ un faisceau de $\mathcal{O}_X \otimes_\Lambda A$-modules
sur $X_\etale$, plat sur $A$.
Consid\'erons les morphismes $i$ et $\pi$ de
Complexe cotangent, \'Equation (\ref{cotangent-equation-pi-spaces}).
Notons $\underline{\mathcal{G}} = \pi^{-1}(\mathcal{G})$. Il est
simplement donn\'e par la r\`egle $(U \to \mathbf{A}) \mapsto \mathcal{G}(U)$;
c'est donc un faisceau de
$\underline{\mathcal{O}}_X \otimes_\Lambda A$-modules.
Notons $i_*\underline{\mathcal{G}}$ le m\^eme faisceau, consid\'er\'e comme un
faisceau de $\mathcal{O} \otimes_\Lambda A$-modules.

\begin{lemma}
\label{lemma-second-equivalence}
Notations et hypoth\`eses comme ci-dessus.
Le foncteur $\pi_!$ induit une \'equivalence de cat\'egories entre
\begin{enumerate}
\item la cat\'egorie
$\textit{Lift}_\mathcal{O}(i_*\underline{\mathcal{G}}, A')$, et
\item la cat\'egorie $\textit{Lift}(\mathcal{G}, A')$.
\end{enumerate}
\end{lemma}

\begin{proof}
FIXME.
\end{proof}

\begin{lemma}
\label{lemma-second-equivalence-obs}
Notations et hypoth\`eses du Lemme \ref{lemma-second-equivalence}.
Consid\'erons l'objet
$$
L = L(\Lambda, X, A, \mathcal{G}) = L\pi_!(Li^*(i_*(\underline{\mathcal{G}})))
$$
de $D(\mathcal{O}_X \otimes_\Lambda A)$. Pour toute surjection $A' \to A$
de $\Lambda$-alg\`ebres dont le noyau $I$ est de carr\'e nul, on a
\begin{enumerate}
\item La cat\'egorie $\textit{Lift}(\mathcal{G}, A')$ est non vide si et
seulement si une certaine classe
$\xi \in \Ext^2_{\mathcal{O}_X \otimes A}(L, \mathcal{G} \otimes_A I)$
est nulle.
\item Si $\textit{Lift}(\mathcal{G}, A')$ est non vide, alors
$\text{Lift}(\mathcal{G}, A')$ est un ensemble principal homog\`ene sous
$\Ext^1_{\mathcal{O}_X \otimes A}(L, \mathcal{G} \otimes_A I)$.
\item \'Etant donn\'e un rel\`evement $\mathcal{G}'$, l'ensemble des
automorphismes de $\mathcal{G}'$ dont l'image inverse est
$\text{id}_\mathcal{G}$ est
canoniquement isomorphe \`a
$\Ext^0_{\mathcal{O}_X \otimes A}(L, \mathcal{G} \otimes_A I)$.
\end{enumerate}
\end{lemma}

\begin{proof}
FIXME.
\end{proof}

\noindent
Enfin, rassemblons ces r\'esultats comme suit.

\begin{proposition}
\label{proposition-conclusion}
Pour $\Lambda$, $X$, $A$ et $\mathcal{F}$ comme ci-dessus, il existe un objet
canonique $L = L(\Lambda, X, A, \mathcal{F})$ de $D(X_A)$ tel que, pour toute
surjection $A' \to A$ de $\Lambda$-alg\`ebres dont le noyau $I$ est de carr\'e
nul, on ait
\begin{enumerate}
\item La cat\'egorie $\textit{Lift}(\mathcal{F}, A')$ est non vide si et
seulement si une certaine classe $\xi \in \Ext^2_{X_A}(L,
\mathcal{F} \otimes_A I)$ est nulle.
\item Si $\textit{Lift}(\mathcal{F}, A')$ est non vide, alors
$\text{Lift}(\mathcal{F}, A')$ est un ensemble principal homog\`ene sous
$\Ext^1_{X_A}(L, \mathcal{F} \otimes_A I)$.
\item \'Etant donn\'e un rel\`evement $\mathcal{F}'$, l'ensemble des
automorphismes de $\mathcal{F}'$ dont l'image inverse est
$\text{id}_\mathcal{F}$ est
canoniquement isomorphe \`a
$\Ext^0_{X_A}(L, \mathcal{F} \otimes_A I)$.
\end{enumerate}
\end{proposition}

\begin{proof}
FIXME.
\end{proof}

\begin{lemma}
\label{lemma-pseudo-coherent}
Dans la situation de la Proposition \ref{proposition-conclusion}, si
$X \to \Spec(\Lambda)$ est localement de type fini et si $\Lambda$ est
noeth\'erien, alors $L$ est pseudo-coh\'erent.
\end{lemma}

\begin{proof}
FIXME.
\end{proof}


\section{Le champ des faisceaux cohérents dans le cas non plat}
\label{section-not-flat}

\noindent
Dans Quot, théorème \ref{quot-theorem-coherent-algebraic},
l'hypothèse selon laquelle $f : X \to B$ est plat n'est pas nécessaire.
Dans cette section, nous modifions la méthode de démonstration, en nous
appuyant sur des idées issues de la géométrie algébrique dérivée, afin de nous
affranchir de l'hypothèse de platitude. Une méthode entièrement différente est
employée dans Quot, section \ref{quot-section-not-flat} pour obtenir exactement
le même résultat ; c'est pourquoi la méthode de la présente section est
obsolète.

\medskip\noindent
La seule étape de la démonstration de
Quot, théorème \ref{quot-theorem-coherent-algebraic}
qui utilise la platitude est l'application de
Quot, lemme \ref{quot-lemma-coherent-defo-thy}.
Ce lemme sert à construire une théorie de l'obstruction comme dans
Artin's Axioms, section \ref{artin-section-dual}.
Sa démonstration repose sur
Deformation Theory, lemmes \ref{defos-lemma-flat-ringed-topoi} et
\ref{defos-lemma-verify-iv-ringed-topoi} de
Deformation Theory, section \ref{defos-section-flat-ringed-topoi}.
C'est ainsi qu'apparaît l'hypothèse de platitude de $f$.
Avant de poursuivre, remarquons que les résultats (2) et (3) de
Deformation Theory, lemmes \ref{defos-lemma-flat-ringed-topoi}
restent valables sans supposer $f$ plat, puisqu'ils reposent sur
Deformation Theory, lemmes \ref{defos-lemma-inf-ext-rel-ringed-topoi}
et \ref{defos-lemma-inf-map-rel-ringed-topoi},
qui ne comportent aucune hypothèse de platitude.

\medskip\noindent
Avant d'entrer dans les détails, motivons la construction provenant de la
géométrie algébrique dérivée, car elle éclairera, nous semble-t-il, la suite.
Soit $A$ une algèbre de type fini sur la base localement noethérienne $S$.
Notons $X \otimes^\mathbf{L} A$ un « changement de base dérivé » de $X$ à
$A$, et notons $i : X_A \to X \otimes^\mathbf{L} A$ le morphisme d'inclusion
canonique. L'objet $X \otimes^\mathbf{L} A$ n'est pas (encore) défini dans le
Projet Stacks ; on peut le concevoir comme l'espace algébrique $X_A$ muni d'un
faisceau simplicial d'anneaux
$\mathcal{O}_{X \otimes^\mathbf{L} A}$ dont les faisceaux d'homologie sont
$$
H_i(\mathcal{O}_{X \otimes^\mathbf{L} A}) =
\text{Tor}^{\mathcal{O}_S}_i(\mathcal{O}_X, A).
$$
Le morphisme $X \otimes^\mathbf{L} A \to \Spec(A)$ est plat
(les termes du faisceau simplicial d'anneaux étant plats sur $A$),
de sorte que les résultats usuels sur les déformations de modules plats lui
sont applicables. On obtient ainsi une théorie de l'obstruction faisant
intervenir les groupes
$$
\Ext^i_{X \otimes^\mathbf{L} A}(i_*\mathcal{F},
i_*\mathcal{F} \otimes_A M)
$$
où $i = 0, 1, 2$ correspondent respectivement aux automorphismes
infinitésimaux, aux déformations infinitésimales et aux obstructions.
Remarquons que toute déformation plate de $i_*\mathcal{F}$ sur
$X \otimes^\mathbf{L} A'$ est nécessairement de la forme
$i'_*\mathcal{F}'$, où $\mathcal{F}'$ est une déformation plate de
$\mathcal{F}$. Par adjonction entre les foncteurs $Li^*$ et
$i_* = Ri_*$, ces groupes d'Ext sont égaux à
$$
\Ext^i_{X_A}(Li^*(i_*\mathcal{F}), \mathcal{F} \otimes_A M)
$$
Nous obtenons donc des groupes d'obstruction exactement de la même forme que
dans la démonstration de Quot, lemme \ref{quot-lemma-coherent-defo-thy},
à ceci près que la première occurrence de $\mathcal{F}$ est remplacée par le
complexe $Li^*(i_*\mathcal{F})$.

\medskip\noindent
Nous démontrons ci-dessous la version non plate du lemme en construisant
« directement » $E(\mathcal{F}) = Li^*(i_*\mathcal{F})$, puis en établissant
directement son lien avec la théorie des déformations de $\mathcal{F}$.
En fait, il suffit de construire $\tau_{\geq -2}E(\mathcal{F})$, puisque seuls
nous intéressent les groupes d'Ext
$\Ext^i_{X_A}(Li^*(i_*\mathcal{F}), \mathcal{F} \otimes_A M)$
pour $i = 0, 1, 2$. On peut même déterminer les faisceaux de cohomologie
$$
H^i(E(\mathcal{F})) =
\left\{
\begin{matrix}
0 & \text{si }i > 0 \\
\mathcal{F} & \text{si } i = 0 \\
0 & \text{si } i = -1 \\
\text{Tor}_1^{\mathcal{O}_S}(\mathcal{O}_X, A)
\otimes_{\mathcal{O}_X} \mathcal{F} &
\text{si } i = -2
\end{matrix}
\right.
$$
Cette observation guidera notre construction de $E(\mathcal{F})$ dans les
remarques qui suivent.

\begin{remark}[Construction directe]
\label{remark-construction-E}
Soit $S$ un schéma. Soit $f : X \to B$ un morphisme d'espaces algébriques
sur $S$. Soit $U$ un autre espace algébrique sur $B$. Notons
$q : X \times_B U \to U$ la seconde projection. Considérons le triangle
distingué
$$
Lq^*L_{U/B} \to L_{X \times_B U/B} \to E \to Lq^*L_{U/B}[1]
$$
de Cotangent, section \ref{cotangent-section-fibre-product}.
Pour tout faisceau $\mathcal{F}$ de
$\mathcal{O}_{X \times_B U}$-modules, on dispose de la classe d'Atiyah
$$
\mathcal{F} \to
L_{X \times_B U/B}
\otimes_{\mathcal{O}_{X \times_B U}}^\mathbf{L} \mathcal{F}[1]
$$
voir Cotangent, section \ref{cotangent-section-atiyah-general}.
Composons-la avec le morphisme vers $E$, puis choisissons un triangle distingué
$$
E(\mathcal{F}) \to \mathcal{F} \to
\mathcal{F} \otimes_{\mathcal{O}_{X \times_B U}}^\mathbf{L} E[1] \to
E(\mathcal{F})[1]
$$
dans $D(\mathcal{O}_{X \times_B U})$.
Par construction, la classe d'Atiyah se relève en un morphisme
$$
e_\mathcal{F} :
E(\mathcal{F})
\longrightarrow
Lq^*L_{U/B} \otimes_{\mathcal{O}_{X \times_B U}}^\mathbf{L} \mathcal{F}[1]
$$
qui s'inscrit dans un morphisme de triangles distingués
$$
\xymatrix{
\mathcal{F} \otimes^\mathbf{L} Lq^*L_{U/B}[1] \ar[r] &
\mathcal{F} \otimes^\mathbf{L} L_{X \times_B U/B}[1] \ar[r] &
\mathcal{F} \otimes^\mathbf{L} E[1] \\
E(\mathcal{F}) \ar[r] \ar[u]^{e_\mathcal{F}} &
\mathcal{F} \ar[r] \ar[u]^{Atiyah} &
\mathcal{F} \otimes^\mathbf{L} E[1] \ar[u]^{=}
}
$$
Une fois donnés $S, B, X, f, U, \mathcal{F}$, nous fixons un choix de
$E(\mathcal{F})$ et de $e_\mathcal{F}$.
\end{remark}

\begin{remark}[Construction de la classe d'obstruction]
\label{remark-construction-ob}
Avec les notations de la remarque \ref{remark-construction-E}, soit
$i : U \to U'$ un épaississement du premier ordre de $U$ sur $B$. Soit
$\mathcal{I} \subset \mathcal{O}_{U'}$ le faisceau quasi-cohérent d'idéaux
définissant $U$ dans $U'$. Le triangle fondamental
$$
Li^*L_{U'/B} \to L_{U/B} \to L_{U/U'} \to Li^*L_{U'/B}[1]
$$
et le morphisme $L_{U/U'} \to \mathcal{I}[1]$ définissent ensemble un
morphisme $e_{U'} : L_{U/B} \to \mathcal{I}[1]$. En le combinant au morphisme
$e_\mathcal{F}$ de la remarque précédente, on obtient
$$
(\text{id}_\mathcal{F} \otimes Lq^*e_{U'}) \cup e_\mathcal{F} :
E(\mathcal{F})
\longrightarrow
\mathcal{F} \otimes_{\mathcal{O}_{X \times_B U}} q^*\mathcal{I}[2]
$$
(on a en outre composé avec le morphisme du produit tensoriel dérivé vers le
produit tensoriel usuel). Autrement dit, on obtient un élément
$$
\xi_{U'} \in
\Ext^2_{\mathcal{O}_{X \times_B U}}(
E(\mathcal{F}),
\mathcal{F} \otimes_{\mathcal{O}_{X \times_B U}} q^*\mathcal{I})
$$
\end{remark}

\begin{lemma}
\label{lemma-ob-is-obstruction}
Dans la situation de la remarque \ref{remark-construction-ob}, supposons que
$\mathcal{F}$ soit plat sur $U$. Alors l'annulation de la classe $\xi_{U'}$
est une condition nécessaire et suffisante pour qu'il existe un
$\mathcal{O}_{X \times_B U'}$-module $\mathcal{F}'$, plat sur $U'$, tel que
$i^*\mathcal{F}' \cong \mathcal{F}$.
\end{lemma}

\begin{proof}[Démonstration (esquisse)]
Nous utiliserons le critère de
Deformation Theory, lemme \ref{defos-lemma-inf-obs-ext-rel-ringed-topoi}.
Pour abréger, nous écrirons
$\mathcal{O} = \mathcal{O}_{X \times_B U}$ et
$\mathcal{O}' = \mathcal{O}_{X \times_B U'}$.
Considérons la suite exacte courte
$$
0 \to \mathcal{I} \to \mathcal{O}_{U'} \to \mathcal{O}_U \to 0.
$$
Soit $\mathcal{J} \subset \mathcal{O}'$ le faisceau quasi-cohérent d'idéaux
définissant $X \times_B U$. D'après ce qui précède, on obtient une suite exacte
$$
\text{Tor}_1^{\mathcal{O}_B}(\mathcal{O}_X, \mathcal{O}_U) \to
q^*\mathcal{I} \to \mathcal{J} \to 0
$$
où $\text{Tor}_1^{\mathcal{O}_B}(\mathcal{O}_X, \mathcal{O}_U)$ est une
abréviation de
$$
\text{Tor}_1^{h^{-1}\mathcal{O}_B}(p^{-1}\mathcal{O}_X, q^{-1}\mathcal{O}_U)
\otimes_{(p^{-1}\mathcal{O}_X\otimes_{h^{-1}\mathcal{O}_B}q^{-1}\mathcal{O}_U)}
\mathcal{O}.
$$
En tensorisant par $\mathcal{F}$, on obtient la suite exacte
$$
\mathcal{F} \otimes_\mathcal{O}
\text{Tor}_1^{\mathcal{O}_B}(\mathcal{O}_X, \mathcal{O}_U) \to
\mathcal{F} \otimes_\mathcal{O}
q^*\mathcal{I} \to
\mathcal{F} \otimes_\mathcal{O} \mathcal{J} \to 0
$$
(Remarquons que les rôles des lettres $\mathcal{I}$ et $\mathcal{J}$ sont
intervertis par rapport aux notations de
Deformation Theory, lemme \ref{defos-lemma-inf-obs-ext-rel-ringed-topoi}.)
La condition (1) du lemme affirme que le dernier morphisme ci-dessus est un
isomorphisme, autrement dit que le premier morphisme est nul.
L'annulation de ce morphisme peut se vérifier sur les fibres en des points géométriques
$\overline{z} = (\overline{x}, \overline{u}) : \Spec(k) \to X \times_B U$.
Posons $R = \mathcal{O}_{B, \overline{b}}$,
$A = \mathcal{O}_{X, \overline{x}}$,
$B = \mathcal{O}_{U, \overline{u}}$ et
$C = \mathcal{O}_{\overline{z}}$.
D'après Cotangent, lemme \ref{cotangent-lemma-fibre-product}
et le triangle qui définit $E(\mathcal{F})$, on a
$$
H^{-2}(E(\mathcal{F}))_{\overline{z}} =
\mathcal{F}_{\overline{z}} \otimes \text{Tor}_1^R(A, B)
$$
Le morphisme $\xi_{U'}$ induit donc un morphisme
$$
\mathcal{F}_{\overline{z}} \otimes \text{Tor}_1^R(A, B)
\longrightarrow
\mathcal{F}_{\overline{z}} \otimes_B \mathcal{I}_{\overline{u}}
$$
Nous affirmons que ce morphisme coïncide avec la fibre en ce point du
morphisme décrit ci-dessus (démonstration omise : il s'agit d'un énoncé
purement algébrique). La condition (1) de
Deformation Theory, lemme \ref{defos-lemma-inf-obs-ext-rel-ringed-topoi}
équivaut donc à l'annulation de
$H^{-2}(\xi_{U'}) :
H^{-2}(E(\mathcal{F})) \to \mathcal{F} \otimes \mathcal{I}$.

\medskip\noindent
Pour achever la démonstration, montrons que, sous l'hypothèse que la condition
(1) est satisfaite, la condition (2) équivaut à l'annulation de $\xi_{U'}$.
Dans la suite de la démonstration, nous écrivons
$\mathcal{F} \otimes \mathcal{I}$ pour désigner
$\mathcal{F} \otimes_\mathcal{O} q^*\mathcal{I} =
\mathcal{F} \otimes_\mathcal{O} \mathcal{J}$. L'examen de la suite spectrale
$$
\Ext^i(H^{-j}(E(\mathcal{F})), \mathcal{F} \otimes \mathcal{I})
\Rightarrow
\Ext^{i + j}(E(\mathcal{F}), \mathcal{F} \otimes \mathcal{I})
$$
et les égalités $H^0(E(\mathcal{F})) = \mathcal{F}$ et
$H^{-1}(E(\mathcal{F})) = 0$ montrent que l'on a une suite exacte
$$
0 \to
\Ext^2(\mathcal{F}, \mathcal{F} \otimes \mathcal{I}) \to
\Ext^2(E(\mathcal{F}), \mathcal{F} \otimes \mathcal{I}) \to
\Hom(H^{-2}(E(\mathcal{F})), \mathcal{F} \otimes \mathcal{I})
$$
Ainsi, notre élément $\xi_{U'}$ appartient à
$\Ext^2(\mathcal{F}, \mathcal{F} \otimes \mathcal{I})$.
Il reste à montrer qu'il coïncide avec l'élément de
Deformation Theory, lemme \ref{defos-lemma-inf-obs-ext-rel-ringed-topoi} ;
nous omettons cette vérification.
\end{proof}

\begin{lemma}
\label{lemma-coherent-defo-thy-general}
Dans Quot, situation \ref{quot-situation-coherent}, supposons que $S$ soit un
schéma localement noethérien et que $S = B$. Posons
$\mathcal{X} = \textit{Coh}_{X/B}$. Alors la versalité est ouverte pour
$\mathcal{X}$ (voir Artin's Axioms, définition
\ref{artin-definition-openness-versality}).
\end{lemma}

\begin{proof}[Démonstration (esquisse)]
Soit $U \to S$ un morphisme de schémas de type fini, soit $x$ un objet de
$\mathcal{X}$ au-dessus de $U$, et soit $u_0 \in U$ un point de type fini tel
que $x$ soit versel en $u_0$. En rétrécissant $U$, on peut supposer que $u_0$
est un point fermé (Morphisms, lemme
\ref{morphisms-lemma-point-finite-type}) et que $U = \Spec(A)$, le morphisme
$U \to S$ se factorisant par un ouvert affine $\Spec(\Lambda)$ de $S$.
Nous appliquerons Artin's Axioms, lemme \ref{artin-lemma-dual-openness} pour
démontrer le lemme. Soit $\mathcal{F}$ le module cohérent sur
$X_A = \Spec(A) \times_S X$, plat sur $A$, qui correspond à l'objet $x$.

\medskip\noindent
Choisissons $E(\mathcal{F})$ et $e_\mathcal{F}$ comme dans
la remarque \ref{remark-construction-E}.
La description des faisceaux de cohomologie de $E(\mathcal{F})$ montre que
$$
\Ext^1(E(\mathcal{F}), \mathcal{F} \otimes_A M) =
\Ext^1(\mathcal{F}, \mathcal{F} \otimes_A M)
$$
pour tout $A$-module $M$. D'après cette égalité et
Deformation Theory, lemme \ref{defos-lemma-inf-ext-rel-ringed-topoi},
on dispose d'un isomorphisme de foncteurs
$$
T_x(M) = \Ext^1_{X_A}(E(\mathcal{F}), \mathcal{F} \otimes_A M)
$$
D'après le lemme \ref{lemma-ob-is-obstruction}, toute surjection $A' \to A$
de $\Lambda$-algèbres dont le noyau $I$ est de carré nul fournit une classe
d'obstruction
$$
\xi_{A'} \in \Ext^2_{X_A}(E(\mathcal{F}), \mathcal{F} \otimes_A I)
$$
Appliquons Derived Categories of Spaces, lemme
\ref{spaces-perfect-lemma-compute-ext} au calcul des groupes d'Ext
$\Ext^i_{X_A}(E(\mathcal{F}), \mathcal{F} \otimes_A M)$
pour $i \leq m$, où $m = 2$. Nous omettons la vérification que
$E(\mathcal{F})$ appartient à $D^-_{\textit{Coh}}$ ; indication : appliquer
Cotangent, lemme \ref{cotangent-lemma-cotangent-finite}.
On obtient un objet parfait $K \in D(A)$ et des isomorphismes fonctoriels
$$
H^i(K \otimes_A^\mathbf{L} M)
\longrightarrow
\Ext^i_{X_A}(E(\mathcal{F}), \mathcal{F} \otimes_A M)
$$
pour $i \leq m$, compatibles aux morphismes de bord. Cet objet $K$, muni des
identifications affichées ci-dessus, fournit une donnée comme dans
Artin's Axioms, situation \ref{artin-situation-dual}.
Enfin, la condition (iv) de
Artin's Axioms, lemme \ref{artin-lemma-dual-obstruction}
résulte d'une variante de
Deformation Theory, lemme \ref{defos-lemma-verify-iv-ringed-topoi},
dont nous omettons l'énoncé et la démonstration.
Ainsi, Artin's Axioms, lemme \ref{artin-lemma-dual-openness}
s'applique, ce qui démontre le lemme.
\end{proof}

\begin{theorem}
\label{theorem-coherent-algebraic-general}
Soit $S$ un schéma. Soit $f : X \to B$ un morphisme d'espaces algébriques
sur $S$. Supposons que $f$ soit de présentation finie et séparé. Alors
$\textit{Coh}_{X/B}$ est un champ algébrique sur $S$.
\end{theorem}

\begin{proof}
Ce théorème reproduit Quot, théorème
\ref{quot-theorem-coherent-algebraic-general}.
Nous le reprenons ici parce que le contenu de la présente section en fournit
une seconde démonstration (comme nous l'avons expliqué au début de la section).
En effet, nous raisonnons exactement comme dans la démonstration de
Quot, théorème \ref{quot-theorem-coherent-algebraic},
si ce n'est que nous y substituons
le lemme \ref{lemma-coherent-defo-thy-general} à
Quot, lemme \ref{quot-lemma-coherent-defo-thy}.
\end{proof}




\section{Modifications}
\label{section-modifications}

\noindent
Voici un résultat désormais obsolète concernant la catégorie de
Algébrisation des espaces formels, équation
(\ref{restricted-equation-modification}).
On trouvera le traitement actuel dans Algébrisation des espaces formels,
section \ref{restricted-section-modifications}.

\begin{lemma}
\label{lemma-henselian}
Soit $(A, \mathfrak m, \kappa)$ un anneau local noethérien.
La catégorie de
Algébrisation des espaces formels, équation
(\ref{restricted-equation-modification})
pour $A$ est équivalente à la catégorie donnée par
Algébrisation des espaces formels, équation
(\ref{restricted-equation-modification})
pour l'hensélisé $A^h$ de $A$.
\end{lemma}

\begin{proof}
C'est un cas particulier d'Algébrisation des espaces formels, lemme
\ref{restricted-lemma-equivalence-to-completion}.
\end{proof}

\noindent
Le lemme suivant sur les singularités rationnelles n'est plus nécessaire
dans le chapitre consacré à la résolution des singularités de surfaces.

\begin{lemma}
\label{lemma-double-dual-rational}
Plaçons-nous dans Résolution des surfaces, situation
\ref{resolve-situation-rational}.
Soit $M$ un $A$-module réflexif fini. Notons
$M \otimes_A \mathcal{O}_X$ l'image réciproque du
$\mathcal{O}_S$-module associé. Alors le morphisme canonique de
$M \otimes_A \mathcal{O}_X$ vers son bidual est surjectif.
\end{lemma}

\begin{proof}
Soit $\mathcal{F} = (M \otimes_A \mathcal{O}_X)^{**}$ le bidual, et soit
$\mathcal{F}' \subset \mathcal{F}$ l'image du morphisme d'évaluation
$M \otimes_A \mathcal{O}_X \to \mathcal{F}$. On dispose alors d'une suite
exacte
$$
0 \to \mathcal{F}' \to \mathcal{F} \to \mathcal{Q} \to 0
$$
Puisque $X$ est normal, les anneaux locaux $\mathcal{O}_{X, x}$ sont des
anneaux de valuation discrète aux points de codimension $1$ (voir
Propriétés, lemme \ref{properties-lemma-criterion-normal}).
Par conséquent, $\mathcal{Q}_x = 0$ en ces points, d'après
Compléments d'algèbre, lemme
\ref{more-algebra-lemma-cokernel-map-double-dual-dvr}.
Ainsi, $\mathcal{Q}$ est supporté par un nombre fini de points fermés et est
engendré par ses sections globales, d'après Cohomologie des schémas, lemme
\ref{coherent-lemma-coherent-support-dimension-0}.
On obtient la suite exacte
$$
0 \to H^0(X, \mathcal{F}') \to H^0(X, \mathcal{F}) \to H^0(X, \mathcal{Q}) \to 0
$$
parce que $\mathcal{F}'$ est engendré par ses sections globales
(Résolution des surfaces, lemme
\ref{resolve-lemma-globally-generated}).
Comme $X \to \Spec(A)$ est un isomorphisme au-dessus du complémentaire du
point fermé et que $M$ est réflexif, les morphismes
$$
M \to H^0(X, \mathcal{F}') \to H^0(X, \mathcal{F})
$$
induisent des isomorphismes après localisation en tout idéal premier non
maximal de $A$. Ces morphismes sont donc des isomorphismes, d'après
Compléments d'algèbre, lemme
\ref{more-algebra-lemma-check-isomorphism-via-depth-and-ass},
et le fait que les modules réflexifs sur les anneaux normaux vérifient la
condition $(S_2)$ (Compléments d'algèbre, lemme
\ref{more-algebra-lemma-reflexive-over-normal}).
On conclut donc, comme voulu, que $\mathcal{Q} = 0$.
\end{proof}


\section{Théorie de l'intersection}
\label{section-intersection-theory}

\begin{lemma}
\label{lemma-good-blowing-up}
Soit $b : X' \to X$ l'éclatement d'un schéma projectif lisse
sur un corps $k$ le long d'un sous-schéma fermé lisse $Z \subset X$.
Considérons le diagramme
$$
\xymatrix{
E \ar[r]_j \ar[d]_\pi & X' \ar[d]^b \\
Z \ar[r]^i & X
}
$$
Supposons qu'il existe un élément de $K_0(X)$ dont la restriction à
$Z$ est égale à la classe de $\mathcal{C}_{Z/X}$ dans $K_0(Z)$.
Alors $[Lb^*\mathcal{O}_Z] = [\mathcal{O}_E] \cdot \alpha''$
dans $K_0(X')$ pour un certain $\alpha'' \in K_0(X')$.
\end{lemma}

\begin{proof}
Les schémas $X$, $X'$, $E$, $Z$ sont projectifs et lisses sur
$k$ et nous avons donc $K'_0(X) = K_0(X) = K_0(\textit{Vect}(X)) =
K_0(D^b_{\textit{Coh}}(X)))$
et de même pour les $3$ autres. Voir Catégories dérivées des schémas, Lemmes
\ref{perfect-lemma-Noetherian-Kprime},
\ref{perfect-lemma-Kprime-K} et \ref{perfect-lemma-K-is-old-K}.
Nous passerons librement d'une de ces versions à l'autre dans cette démonstration.
Considérons la suite exacte
$$
0 \to \mathcal{F} \to \pi^*\mathcal{C}_{Z/X} \to \mathcal{C}_{E/X'} \to 0
$$
de $\mathcal{O}_E$-modules localement libres de type fini qui définit $\mathcal{F}$.
Remarquons que $\mathcal{C}_{E/X'} = \mathcal{O}_{X'}(-E)|_E$
est la restriction du $\mathcal{O}_X$-module inversible
$\mathcal{O}_{X'}(-E)$.
Soit $\alpha \in K_0(X)$ un élément tel que
$i^*\alpha = [\mathcal{C}_{Z/X}]$ dans $K_0(Z)$.
Posons $\alpha' = b^*\alpha - [\mathcal{O}_{X'}(-E)]$.
Alors $j^*\alpha' = [\mathcal{F}]$. On en déduit que
$j^*\lambda^i(\alpha') = [\wedge^i(\mathcal{F})]$ d'après
Théories cohomologiques de Weil, Lemme \ref{weil-lemma-lambda-operations}.
Cela signifie que $[\mathcal{O}_E] \cdot \alpha' = [\wedge^i\mathcal{F}]$
dans $K_0(X)$; voir
Catégories dérivées des schémas, Lemme \ref{perfect-lemma-projection-formula}.
Soit $r$ la codimension maximale d'une composante irréductible de $Z$
dans $X$. Un calcul que nous omettons montre que
$H^{-i}(Lb^*\mathcal{O}_Z) = \wedge^i\mathcal{F}$
pour $i \geq 0, 1, \ldots, r - 1$ et est nul dans les autres degrés.
Il s'ensuit que, dans $K_0(X)$, on a
\begin{align*}
[Lb^*\mathcal{O}_Z] & =
\sum\nolimits_{i = 0, \ldots, r - 1} (-1)^i[\wedge^i\mathcal{F}] \\
& =
\sum\nolimits_{i = 0, \ldots, r - 1} (-1)^i[\mathcal{O}_E] \lambda^i(\alpha') \\
& =
[\mathcal{O}_E] \left(\sum\nolimits_{i = 0, \ldots, r - 1}
(-1)^i \lambda^i(\alpha')\right)
\end{align*}
Cela démontre le lemme en prenant
$\alpha'' = \sum_{i = 0, \ldots, r - 1} (-1)^i \lambda^i(\alpha')$.
\end{proof}



\begin{lemma}
\label{lemma-gysin-factors-principal}
Soit $(S, \delta)$ comme dans Homologie de Chow, Situation \ref{chow-situation-setup}.
Soit $X$ localement de type fini sur $S$.
Supposons $X$ intègre et $n = \dim_\delta(X)$.
Soit $a \in \Gamma(X, \mathcal{O}_X)$ une fonction non nulle.
Soit $i : D = Z(a) \to X$ l'immersion fermée du sous-schéma des zéros de $a$.
Soit $f \in R(X)^*$.
Dans ce cas, $i^*\text{div}_X(f) = 0$ dans $A_{n - 2}(D)$.
\end{lemma}

\begin{proof}
Cas particulier d'Homologie de Chow, Lemme \ref{chow-lemma-gysin-factors-general}.
\end{proof}

\begin{remark}
\label{remark-not-true-not-quasi-compact}
Cette remarque disait autrefois que l'on ne savait pas si les flèches de
Homologie de Chow, Lemme \ref{chow-lemma-cycles-k-group}, étaient des isomorphismes
en général. Nous avons cependant trouvé depuis une démonstration de ce fait.
\end{remark}







\subsection{Lemmes d'éclatement}
\label{section-blowing-up-lemmas}

\noindent
Dans cette section, nous démontrons quelques lemmes permettant de représenter
des diviseurs de Cartier par des diviseurs de Cartier effectifs convenables
sur des éclatements. Ces lemmes se trouvent dans \cite[Section 2.4]{F}.
Nous avons adapté leur énoncé afin qu'ils restent valables
sans hypothèse de type fini. Il se peut que le morphisme $b$
du Lemme \ref{lemma-blowing-up-intersections} soit un composé
d'une infinité d'éclatements, mais, au-dessus de tout ouvert quasi-compact
$W \subset X$ donné, un nombre fini d'éclatements suffit
(et c'est le résultat de loc.\ cit.).

\begin{lemma}
\label{lemma-push-pull-effective-Cartier}
Soit $(S, \delta)$ comme dans Homologie de Chow, Situation \ref{chow-situation-setup}.
Soient $X$, $Y$ localement de type fini sur $S$.
Soit $f : X \to Y$ un morphisme propre.
Soit $D \subset Y$ un diviseur de Cartier effectif.
Supposons $X$, $Y$ intègres, $n = \dim_\delta(X) = \dim_\delta(Y)$ et
que $f$ soit dominant. Alors
$$
f_*[f^{-1}(D)]_{n - 1} = [R(X) : R(Y)] [D]_{n - 1}.
$$
En particulier, si $f$ est birationnel, alors $f_*[f^{-1}(D)]_{n - 1} = [D]_{n - 1}$.
\end{lemma}

\begin{proof}
Cela résulte immédiatement d'Homologie de Chow, Lemme \ref{chow-lemma-equal-c1-as-cycles},
et du fait que $D$ est le sous-schéma des zéros
de la section canonique $1_D$ de $\mathcal{O}_X(D)$.
\end{proof}

\begin{lemma}
\label{lemma-blowing-up-denominators}
Soit $(S, \delta)$ comme dans Homologie de Chow, Situation \ref{chow-situation-setup}.
Soit $X$ localement de type fini sur $S$.
Supposons $X$ intègre avec $\dim_\delta(X) = n$.
Soit $\mathcal{L}$ un $\mathcal{O}_X$-module inversible.
Soit $s$ une section méromorphe non nulle de $\mathcal{L}$.
Soit $U \subset X$ le plus grand sous-schéma ouvert tel que
$s$ corresponde à une section de $\mathcal{L}$ sur $U$.
Il existe un morphisme projectif
$$
\pi : X' \longrightarrow X
$$
tel que
\begin{enumerate}
\item $X'$ est intègre,
\item $\pi|_{\pi^{-1}(U)} : \pi^{-1}(U) \to U$ est un isomorphisme,
\item il existe des diviseurs de Cartier effectifs $D, E \subset X'$
tels que
$$
\pi^*\mathcal{L} = \mathcal{O}_{X'}(D - E),
$$
\item la section méromorphe $s$ correspond, par l'isomorphisme ci-dessus,
à la section méromorphe $1_D \otimes (1_E)^{-1}$ (voir
Diviseurs, Définition
\ref{divisors-definition-invertible-sheaf-effective-Cartier-divisor}),
\item on a
$$
\pi_*([D]_{n - 1} - [E]_{n - 1}) = \text{div}_\mathcal{L}(s)
$$
dans $Z_{n - 1}(X)$.
\end{enumerate}
\end{lemma}

\begin{proof}
Soit $\mathcal{I} \subset \mathcal{O}_X$ le faisceau quasi-cohérent d'idéaux
des dénominateurs de $s$; voir Diviseurs, Définition
\ref{divisors-definition-regular-meromorphic-ideal-denominators}.
D'après Diviseurs, Lemme \ref{divisors-lemma-blowing-up-denominators},
on obtient (2), (3) et (4).
D'après Diviseurs, Lemme \ref{divisors-lemma-blow-up-integral-scheme},
on obtient (1). D'après Diviseurs, Lemme \ref{divisors-lemma-blowing-up-projective},
le morphisme $\pi$ est projectif.
Il reste à démontrer (5).
D'après Homologie de Chow, Lemme \ref{chow-lemma-equal-c1-as-cycles}, on a
$$
\pi_*(\text{div}_{\mathcal{L}'}(s')) = \text{div}_\mathcal{L}(s).
$$
Il suffit donc de montrer que
$\text{div}_{\mathcal{L}'}(s') = [D]_{n - 1} - [E]_{n - 1}$.
Cela résulte de l'égalité
$s' = 1_D \otimes 1_E^{-1}$ et de l'additivité; voir
Diviseurs, Lemme \ref{divisors-lemma-c1-additive}.
\end{proof}

\begin{definition}
\label{definition-epsilon}
Soit $(S, \delta)$ comme dans Homologie de Chow, Situation \ref{chow-situation-setup}.
Soit $X$ localement de type fini sur $S$.
Supposons $X$ intègre et $\dim_\delta(X) = n$.
Soient $D_1, D_2$ deux diviseurs de Cartier effectifs sur $X$.
Soit $Z \subset X$ un sous-schéma fermé intègre
avec $\dim_\delta(Z) = n - 1$. L'{\it invariant $\epsilon$}
de cette situation est
$$
\epsilon_Z(D_1, D_2) = n_Z \cdot m_Z
$$
où $n_Z$, resp.\ $m_Z$, est le coefficient de
$Z$ dans le $(n - 1)$-cycle $[D_1]_{n - 1}$, resp.\ $[D_2]_{n - 1}$.
\end{definition}

\begin{lemma}
\label{lemma-two-divisors}
Soit $(S, \delta)$ comme dans Homologie de Chow, Situation \ref{chow-situation-setup}.
Soit $X$ localement de type fini sur $S$.
Supposons $X$ intègre et $\dim_\delta(X) = n$.
Soient $D_1, D_2$ deux diviseurs de Cartier effectifs sur $X$.
Soit $Z$ un sous-schéma ouvert et fermé du schéma $D_1 \cap D_2$.
Supposons $\dim_\delta(D_1 \cap D_2 \setminus Z) \leq n - 2$.
Alors il existe un morphisme
$b : X' \to X$ et des diviseurs de Cartier
$D_1', D_2', E$ sur $X'$ ayant les propriétés suivantes :
\begin{enumerate}
\item $X'$ est intègre,
\item $b$ est projectif,
\item $b$ est l'éclatement de $X$ le long du sous-schéma fermé $Z$,
\item $E = b^{-1}(Z)$,
\item $b^{-1}(D_1) = D'_1 + E$, et $b^{-1}D_2 = D_2' + E$,
\item $\dim_\delta(D'_1 \cap D'_2) \leq n - 2$, et si
$Z = D_1 \cap D_2$ alors $D'_1 \cap D'_2 = \emptyset$,
\item pour tout sous-schéma fermé intègre $W'$
avec $\dim_\delta(W') = n - 1$, on a
\begin{enumerate}
\item si $\epsilon_{W'}(D'_1, E) > 0$, alors, en posant
$W = b(W')$, on a
$\dim_\delta(W) = n - 1$ et
$$
\epsilon_{W'}(D'_1, E) < \epsilon_W(D_1, D_2),
$$
\item si $\epsilon_{W'}(D'_2, E) > 0$, alors, en posant
$W = b(W')$, on a
$\dim_\delta(W) = n - 1$ et
$$
\epsilon_{W'}(D'_2, E) < \epsilon_W(D_1, D_2),
$$
\end{enumerate}
\end{enumerate}
\end{lemma}

\begin{proof}
Remarquons que le faisceau quasi-cohérent d'idéaux
$\mathcal{I} = \mathcal{I}_{D_1} + \mathcal{I}_{D_2}$
définit l'intersection schématique $D_1 \cap D_2 \subset X$.
Comme $Z$ est une réunion de composantes connexes de $D_1 \cap D_2$,
on voit que, pour tout $z \in Z$, le noyau de
$\mathcal{O}_{X, z} \to \mathcal{O}_{Z, z}$ est égal à $\mathcal{I}_z$.
Soit $b : X' \to X$ l'éclatement de $X$ le long de $Z$. (Ainsi, localement pour la topologie de Zariski
au voisinage de $Z$, c'est l'éclatement de $X$ le long de $\mathcal{I}$.)
Notons $E = b^{-1}(Z)$ le diviseur de Cartier effectif correspondant; voir
Diviseurs,
Lemme \ref{divisors-lemma-blowing-up-gives-effective-Cartier-divisor}.
Comme $Z \subset D_1$, on a $E \subset f^{-1}(D_1)$ et donc
$D_1 = D_1' + E$ pour un certain diviseur de Cartier effectif $D'_1 \subset X'$;
voir Diviseurs, Lemme \ref{divisors-lemma-difference-effective-Cartier-divisors}.
De même, $D_2 = D_2' + E$. Cela établit les assertions (1) -- (5).

\medskip\noindent
Remarquons que, si $W'$ est comme en (7) (a) ou (7) (b), alors l'image $W$
de $W'$ est contenue dans $D_1 \cap D_2$. Si $W$ n'est pas contenue dans
$Z$, alors $b$ est un isomorphisme au point générique de $W$ et
on voit que $\dim_\delta(W) = \dim_\delta(W') = n - 1$, ce qui
contredit l'hypothèse
$\dim_\delta(D_1 \cap D_2 \setminus Z) \leq n - 2$.
Ainsi, $W \subset Z$. Cela signifie que,
pour démontrer (6) et (7), on peut travailler localement au voisinage de $Z$ sur $X$.

\medskip\noindent
On peut donc supposer que $X = \Spec(A)$, où
$A$ est un anneau intègre noethérien, et que $D_1 = \Spec(A/a)$,
$D_2 = \Spec(A/b)$ et $Z = D_1 \cap D_2$.
Posons $I = (a, b)$. Comme $A$ est intègre et $a, b \not = 0$, on peut
recouvrir l'éclatement par deux cartes affines, à savoir
$U = \Spec(A[s]/(as - b))$ et $V = \Spec(A[t]/(bt -a))$.
Ces cartes sont recollées au moyen de l'isomorphisme
$A[s, s^{-1}]/(as - b) \cong A[t, t^{-1}]/(bt - a)$
qui envoie $s$ sur $t^{-1}$.
Le diviseur de Cartier effectif $E$ est décrit par
$\Spec(A[s]/(as - b, a)) \subset U$ et
$\Spec(A[t]/(bt - a, b)) \subset V$.
Le sous-schéma fermé $D'_1$ correspond à
$\Spec(A[t]/(bt - a, t)) \subset U$.
Le sous-schéma fermé $D'_2$ correspond à
$\Spec(A[s]/(as -b, s)) \subset V$.
Puisque ``$ts = 1$'', on a $D'_1 \cap D'_2 = \emptyset$.

\medskip\noindent
Supposons donné un idéal premier $\mathfrak q \subset A[s]/(as - b)$
de hauteur un tel que $s, a \in \mathfrak q$.
Soit $\mathfrak p \subset A$ l'idéal premier correspondant de $A$.
Remarquons que $a, b \in \mathfrak p$.
D'après la formule des dimensions, on a aussi $\dim(A_{\mathfrak p}) = 1$.
La dernière assertion à démontrer est
$$
\text{ord}_{A_{\mathfrak p}}(a)
\text{ord}_{A_{\mathfrak p}}(b)
>
\text{ord}_{B_{\mathfrak q}}(a)
\text{ord}_{B_{\mathfrak q}}(s)
$$
où $B = A[s]/(as - b)$. D'après
Algèbre, Lemme \ref{algebra-lemma-quasi-finite-extension-dim-1},
on a $\text{ord}_{A_{\mathfrak p}}(x) \geq \text{ord}_{B_{\mathfrak q}}(x)$
pour $x = a, b$. Puisque $\text{ord}_{B_{\mathfrak q}}(s) > 0$, on conclut
par additivité de la fonction $\text{ord}$ et par le fait que
$as = b$.
\end{proof}

\begin{definition}
\label{definition-locally-finite-sum-effective-Cartier-divisors}
Soit $X$ un schéma.
Soit $\{D_i\}_{i \in I}$ une famille localement finie
de diviseurs de Cartier effectifs sur $X$.
Supposons donnée une application
$I \to \mathbf{Z}_{\geq 0}$, $i \mapsto n_i$.
La {\it somme des diviseurs de Cartier effectifs}
$D = \sum n_i D_i$ est l'unique diviseur de Cartier effectif
$D \subset X$ tel que, pour tout ouvert quasi-compact $U \subset X$,
on ait $D|_U = \sum_{D_i \cap U \not = \emptyset} n_iD_i|_U$,
la somme du membre de droite étant celle de Diviseurs,
Définition \ref{divisors-definition-sum-effective-Cartier-divisors}.
\end{definition}

\begin{lemma}
\label{lemma-sum-divisors-associated-Weil}
Soit $(S, \delta)$ comme dans Homologie de Chow, situation \ref{chow-situation-setup}.
Soit $X$ localement de type fini sur $S$.
Supposons $X$ intègre et $\dim_\delta(X) = n$.
Soit $\{D_i\}_{i \in I}$ une famille localement finie
de diviseurs de Cartier effectifs sur $X$.
Donnons-nous des $n_i \geq 0$ pour $i \in I$.
Alors
$$
[D]_{n - 1} = \sum\nolimits_i n_i[D_i]_{n - 1}
$$
dans $Z_{n - 1}(X)$.
\end{lemma}

\begin{proof}
Puisque nous démontrons une égalité de cycles, nous pouvons raisonner localement sur $X$.
On se ramène ainsi à une somme finie, puis, par récurrence, à une somme de
deux diviseurs de Cartier effectifs $D = D_1 + D_2$.
D'après Homologie de Chow, lemme \ref{chow-lemma-compute-c1}, on a
$D_1 = \text{div}_{\mathcal{O}_X(D_1)}(1_{D_1})$, où
$1_{D_1}$ désigne la section canonique de $\mathcal{O}_X(D_1)$.
La même assertion vaut bien entendu pour $D_2$ et $D$.
Comme $1_D = 1_{D_1} \otimes 1_{D_2}$ sous l'identification
$\mathcal{O}_X(D) = \mathcal{O}_X(D_1) \otimes \mathcal{O}_X(D_2)$,
le résultat découle de Diviseurs, lemme \ref{divisors-lemma-c1-additive}.
\end{proof}

\begin{lemma}
\label{lemma-blowing-up-intersections}
Soit $(S, \delta)$ comme dans Homologie de Chow, situation \ref{chow-situation-setup}.
Soit $X$ localement de type fini sur $S$.
Supposons $X$ intègre et $\dim_\delta(X) = d$.
Soit $\{D_i\}_{i \in I}$ une famille localement finie de
diviseurs de Cartier effectifs sur $X$.
Supposons que, pour tout $\{i, j, k\} \subset I$ tel que $\#\{i, j, k\} = 3$,
on ait $D_i \cap D_j \cap D_k = \emptyset$.
Il existe alors
\begin{enumerate}
\item un sous-schéma ouvert $U \subset X$ tel que
$\dim_\delta(X \setminus U) \leq d - 3$,
\item un morphisme $b : U' \to U$, et
\item des diviseurs de Cartier effectifs $\{D'_j\}_{j \in J}$ sur $U'$,
\end{enumerate}
vérifiant les propriétés suivantes :
\begin{enumerate}
\item $b$ est un morphisme propre $b : U' \to U$,
\item $U'$ est intègre,
\item $b$ est un isomorphisme au-dessus du complémentaire de la réunion des
intersections deux à deux des $D_i|_U$,
\item $\{D'_j\}_{j \in J}$ est une famille localement finie de diviseurs de
Cartier effectifs sur $U'$,
\item $\dim_\delta(D'_j \cap D'_{j'}) \leq d - 2$ si $j \not = j'$, et
\item $b^{-1}(D_i|_U) = \sum n_{ij} D'_j$ pour certains $n_{ij} \geq 0$.
\end{enumerate}
De plus, si $X$ est quasi-compact, on peut supposer $U = X$ ci-dessus.
\end{lemma}

\begin{proof}
Commençons par le cas quasi-compact, qui est peut-être
le plus intéressant. Dans ce cas, nous construisons par récurrence une suite
d'éclatements
$$
X = X_0 \xleftarrow{b_0} X_1 \xleftarrow{b_1} X_2 \leftarrow \ldots
$$
et des ensembles finis de diviseurs de Cartier effectifs $\{D_{n, i}\}_{i \in I_n}$.
À chaque étape, ceux-ci auront la propriété que toute intersection triple
$D_{n, i} \cap D_{n, j} \cap D_{n, k}$ est vide.
En outre, pour tout $n \geq 0$, nous aurons
$I_{n + 1} = I_n \amalg P(I_n)$, où $P(I_n)$ désigne
l'ensemble des paires d'éléments de $I_n$. Enfin, nous aurons
$$
b_n^{-1}(D_{n, i}) = D_{n + 1, i} +
\sum\nolimits_{i' \in I_n, i' \not = i} D_{n + 1, \{i, i'\}}
$$
On en déduit que, pour tout $n \geq 0$,
$(b_0 \circ \ldots \circ b_n)^{-1}(D_i)$ est une combinaison linéaire
à coefficients entiers positifs ou nuls des diviseurs $D_{n + 1, j}$, $j \in I_{n + 1}$.

\medskip\noindent
Pour amorcer la récurrence, posons $X_0 = X$,
$I_0 = I$ et $D_{0, i} = D_i$.

\medskip\noindent
Étant donné $(X_n, \{D_{n, i}\}_{i \in I_n})$, soit $X_{n + 1}$
l'éclatement de $X_n$ le long du sous-schéma fermé
$Z_n = \bigcup_{\{i, i'\} \in P(I_n)} D_{n, i} \cap D_{n, i'}$.
Remarquons que les sous-schémas fermés $D_{n, i} \cap D_{n, i'}$ sont deux à deux
disjoints en raison de notre hypothèse sur les intersections triples.
Autrement dit, on peut écrire
$Z_n = \coprod_{\{i, i'\} \in P(I_n)} D_{n, i} \cap D_{n, i'}$.
De plus, dans un voisinage de Zariski de $D_{n, i} \cap D_{n, i'}$, le
morphisme $b_n$ est l'éclatement du schéma $X_n$
le long du sous-schéma fermé $D_{n, i} \cap D_{n, i'}$, et les conclusions
du lemme \ref{lemma-two-divisors} s'appliquent.
En posant donc $D_{n + 1, \{i, i'\}} = b_n^{-1}(D_{n, i} \cap D_{n, i'})$,
on obtient un diviseur de Cartier effectif.
Les diviseurs de Cartier $D_{n + 1, \{i, i'\}}$ sont deux à deux disjoints.
Manifestement, on a
$b_n^{-1}(D_{n, i}) \supset D_{n + 1, \{i, i'\}}$ pour
tout $i' \in I_n$, $i' \not = i$. Par conséquent, en appliquant
Diviseurs, lemme \ref{divisors-lemma-difference-effective-Cartier-divisors},
on voit que l'on a bien $b_n^{-1}(D_{n, i}) = D_{n + 1, i} +
\sum\nolimits_{i' \in I_n, i' \not = i} D_{n + 1, \{i, i'\}}$
pour un certain diviseur de Cartier effectif $D_{n + 1, i}$ sur $X_{n + 1}$.
Au voisinage de $D_{n + 1, \{i, i'\}}$, ces diviseurs
$D_{n + 1, i}$ jouent le rôle des diviseurs primés du
lemme \ref{lemma-two-divisors}. En particulier, on en déduit que
$D_{n + 1, i} \cap D_{n + 1, i'} = \emptyset$ si $i \not = i'$,
$i, i' \in I_n$, d'après le point (6) du lemme \ref{lemma-two-divisors}.
Cela implique déjà que les intersections triples
des diviseurs $D_{n + 1, i}$ sont vides.

\medskip\noindent
À ce stade, nous pouvons utiliser la quasi-compacité de $X$
pour conclure que l'invariant
\begin{equation}
\label{equation-invariant}
\epsilon(X, \{D_i\}_{i \in I}) =
\max\{\epsilon_Z(D_i, D_{i'}) \mid
Z \subset X,
\dim_\delta(Z) = d - 1,
\{i, i'\} \in P(I)\}
\end{equation}
est fini, car chaque $D_i$ n'a en définitive qu'un nombre fini de composantes
irréductibles. Nous affirmons que, pour un certain $n$, l'invariant
$\epsilon(X_n, \{D_{n, i}\}_{i \in I_n})$ est nul. En effet, sinon,
le lemme \ref{lemma-two-divisors} donnerait une suite strictement décroissante
$$
\epsilon(X, \{D_i\}_{i \in I})
=
\epsilon(X_0, \{D_{0, i}\}_{i \in I_0})
>
\epsilon(X_1, \{D_{1, i}\}_{i \in I_1})
>
\ldots
$$
d'entiers strictement positifs, ce qui est contradictoire. Choisissons $n$ tel que
l'invariant $\epsilon(X_n, \{D_{n, i}\}_{i \in I_n})$ soit nul.
Cela signifie qu'il n'existe ni sous-schéma fermé intègre $Z \subset X_n$,
ni paire d'indices $i, i' \in I_n$
tels que $\epsilon_Z(D_{n, i}, D_{n, i'}) > 0$.
Autrement dit, $\dim_\delta(D_{n, i} \cap D_{n, i'}) \leq d - 2$ pour
toute paire $\{i, i'\} \in P(I_n)$, comme souhaité.

\medskip\noindent
Passons maintenant au cas général où l'on ne suppose plus
le schéma $X$ quasi-compact. Le problème que pose la méthode de
la première partie de la démonstration est que l'on peut obtenir une suite infinie
d'éclatements dont les centres dominent un point fixé de $X$. Pour
éviter cela, nous retirons à chaque étape des sous-ensembles fermés convenables de codimension $\geq 3$.
Plus précisément, nous construirons par récurrence
une suite de morphismes de la forme suivante :
$$
\xymatrix{
X = X_0 \\
U_0 \ar[u]^{j_0} & X_1 \ar[l]_{b_0} \\
 & U_1 \ar[u]^{j_1} & X_2 \ar[l]_{b_1} \\
 & & U_2 \ar[u]^{j_2} & X_3 \ar[l]_{b_2}
}
$$
Chacun des morphismes $j_n : U_n \to X_n$ sera une immersion ouverte.
Chacun des morphismes $b_n : X_{n + 1} \to U_n$ sera un morphisme propre
birationnel entre schémas intègres. Comme dans le cas quasi-compact, nous aurons
des diviseurs de Cartier effectifs $\{D_{n, i}\}_{i \in I_n}$ sur $X_n$.
À chaque étape, ceux-ci auront la propriété que toute intersection triple
$D_{n, i} \cap D_{n, j} \cap D_{n, k}$ est vide.
En outre, pour tout $n \geq 0$, nous aurons
$I_{n + 1} = I_n \amalg P(I_n)$, où $P(I_n)$ désigne
l'ensemble des paires d'éléments de $I_n$.
Enfin, nous ferons en sorte que
$$
b_n^{-1}(D_{n, i}|_{U_n}) = D_{n + 1, i} +
\sum\nolimits_{i' \in I_n, i' \not = i} D_{n + 1, \{i, i'\}}
$$

\medskip\noindent
Amorçons la récurrence en posant $X_0 = X$,
$I_0 = I$ et $D_{0, i} = D_i$.

\medskip\noindent
Étant donné $(X_n, \{D_{n, i}\})$, construisons le sous-schéma ouvert
$U_n$ comme suit. Pour toute paire $\{i, i'\} \in P(I_n)$, considérons
le sous-schéma fermé $D_{n, i} \cap D_{n, i'}$. Ses composantes irréductibles
« bonnes » sont celles de $\delta$-dimension $d - 2$, et ses composantes irréductibles
« mauvaises » celles de $\delta$-dimension $d - 1$.
Posons
$$
\text{Bad}(i, i')
=
\bigcup\nolimits_{W \subset D_{n, i} \cap D_{n, i'}
\text{ comp.\ irr.\ avec }\dim_\delta(W) = d - 1} W
$$
et, de même,
$$
\text{Good}(i, i')
=
\bigcup\nolimits_{W \subset D_{n, i} \cap D_{n, i'}
\text{ comp.\ irr.\ avec }\dim_\delta(W) = d - 2} W.
$$
Alors $D_{n, i} \cap D_{n, i'} = \text{Bad}(i, i') \cup \text{Good}(i, i')$
et, de plus, on a
$\dim_\delta(\text{Bad}(i, i') \cap \text{Good}(i, i')) \leq d - 3$.
Voici le choix de $U_n$ :
$$
U_n
=
X_n
\setminus
\bigcup\nolimits_{\{i, i'\} \in P(I_n)}
\text{Bad}(i, i') \cap \text{Good}(i, i').
$$
Notre condition sur les intersections triples des diviseurs $D_{n, i}$
montre que la réunion est en fait disjointe. De plus,
on voit que, comme schéma,
$$
D_{n, i}|_{U_n} \cap D_{n, i'}|_{U_n}
=
Z_{n, i, i'} \amalg G_{n, i, i'}
$$
où $Z_{n, i, i'}$ est $\delta$-équidimensionnel de dimension $d - 1$
et $G_{n, i, i'}$ est $\delta$-équidimensionnel de dimension $d - 2$.
(Ainsi, topologiquement, $Z_{n, i, i'}$ est la réunion des composantes mauvaises,
privée de ses intersections avec les composantes bonnes.) Enfin, posons
$$
Z_n =
\bigcup\nolimits_{\{i, i'\} \in P(I_n)} Z_{n, i, i'} =
\coprod\nolimits_{\{i, i'\} \in P(I_n)} Z_{n, i, i'},
$$
et soit $b_n : X_{n + 1} \to U_n$ l'éclatement de centre $Z_n$.
Remarquons que le lemme \ref{lemma-two-divisors}
s'applique au morphisme $b_n : X_{n + 1} \to U_n$, localement au voisinage
de chacun des lieux $D_{n, i}|_{U_n} \cap D_{n, i'}|_{U_n}$. Ainsi,
exactement comme dans la première partie de la démonstration, nous obtenons des diviseurs de
Cartier effectifs $D_{n + 1, \{i, i'\}}$ pour $\{i, i'\} \in P(I_n)$
et des diviseurs de Cartier effectifs $D_{n + 1, i}$ pour $i \in I_n$
tels que
$b_n^{-1}(D_{n, i}|_{U_n}) = D_{n + 1, i} +
\sum\nolimits_{i' \in I_n, i' \not = i} D_{n + 1, \{i, i'\}}$.
Pour tout $n$, notons $\pi_n : X_n \to X$ le morphisme obtenu
comme le composé $j_0 \circ b_0 \circ \ldots \circ j_{n - 1} \circ b_{n - 1}$.

\medskip\noindent
{\bf Assertion :} étant donné un ouvert quasi-compact $V \subset X$,
pour tout $n$ assez grand, les morphismes
$$
\pi_n^{-1}(V) \leftarrow \pi_{n + 1}^{-1}(V) \leftarrow \ldots
$$
sont tous des isomorphismes. En effet, si le morphisme
$\pi_n^{-1}(V) \leftarrow \pi_{n + 1}^{-1}(V)$ n'est pas un isomorphisme,
alors $Z_{n, i, i'} \cap \pi_n^{-1}(V) \not = \emptyset$ pour une certaine
paire $\{i, i'\} \in P(I_n)$. Il existe donc une composante irréductible
$W \subset D_{n, i} \cap D_{n, i'}$ telle que $\dim_\delta(W) = d - 1$.
En particulier, on a $\epsilon_W(D_{n, i}, D_{n, i'}) > 0$.
En appliquant à plusieurs reprises le lemme \ref{lemma-two-divisors}, on obtient
$$
\epsilon_W(D_{n, i}, D_{n, i'})
<
\epsilon(V, \{D_i|_V\}) - n
$$
où $\epsilon(V, \{D_i|_V\})$ est défini comme dans (\ref{equation-invariant}).
Puisque $V$ est quasi-compact, on a
$\epsilon(V, \{D_i|_V\}) < \infty$ et, en choisissant $n > \epsilon(V, \{D_i|_V\})$,
on obtient le résultat.

\medskip\noindent
Notons que, par construction, la différence $X_n \setminus U_n$
vérifie $\dim_\delta(X_n \setminus U_n) \leq d - 3$.
Soit $T_n = \pi_n(X_n \setminus U_n)$ son image dans $X$.
En remontant le diagramme de morphismes ci-dessus et en utilisant le fait que chaque $b_n$ est fermé,
on obtient que $T_0 \cup \ldots \cup T_n$ est un sous-ensemble fermé de $X$
pour tout $n$. Par construction, tout $t \in T_n$ vérifie $\delta(t) \leq d - 3$.
Ainsi, $\overline{T_n} \subset X$ est un sous-ensemble fermé
tel que $\dim_\delta(T_n) \leq d - 3$. D'après l'assertion ci-dessus,
pour tout ouvert quasi-compact $V \subset X$, on a
$T_n \cap V \not = \emptyset$ pour au plus un nombre fini de $n$.
Ainsi, $\{\overline{T_n}\}_{n \geq 0}$ est une famille localement finie
de sous-ensembles fermés, et l'on peut poser
$U = X \setminus \bigcup \overline{T_n}$. Ce sera le
$U$ du lemme.

\medskip\noindent
Par construction de $U$, on a $U_n \cap \pi_n^{-1}(U) = \pi_n^{-1}(U)$.
Tous les morphismes
$$
b_n : \pi_{n + 1}^{-1}(U) \longrightarrow \pi_n^{-1}(U)
$$
sont donc propres. De plus, d'après l'assertion, ils deviennent des isomorphismes
à partir d'un certain rang au-dessus de tout ouvert quasi-compact de $X$. On peut donc définir
$$
U' = \lim_n \pi_n^{-1}(U).
$$
Le morphisme induit $b : U' \to U$ est propre, car la propreté est locale
sur $U$ et, au-dessus de tout ouvert quasi-compact, la limite est stationnaire. De même,
posons $J = \bigcup_{n \geq 0} I_n$ au moyen des inclusions
$I_n \to I_{n + 1}$ issues de la construction. Pour $j \in J$, choisissons
un $n_0$ tel que $j$ corresponde à $i \in I_{n_0}$ et définissons
$D'_j = \lim_{n \geq n_0} D_{n, i}$. Cette définition a encore un sens,
car, localement sur $X$, les morphismes se stabilisent.
Les autres assertions du lemme se vérifient comme dans le cas où
$X$ est quasi-compact.
\end{proof}









\section{Commutativité de l'intersection des diviseurs}
\label{section-commutativity}

\noindent
Les résultats de cette section ont été initialement utilisés pour donner une autre
démonstration des lemmes de Homologie de Chow, Section \ref{chow-section-commutativity},
ainsi qu'une version faible de
Homologie de Chow, Lemme \ref{chow-lemma-gysin-commutes-gysin}.

\begin{lemma}
\label{lemma-improved-additivity}
Soit $(S, \delta)$ comme dans Homologie de Chow, Situation \ref{chow-situation-setup}.
Soit $X$ localement de type fini sur $S$.
Soit $\{i_j : D_j \to X \}_{j \in J}$ une famille localement finie
de diviseurs de Cartier effectifs sur $X$. Soient $n_j > 0$, $j\in J$.
Posons $D = \sum_{j \in J} n_j D_j$, et notons $i : D \to X$ le
morphisme d'inclusion. Soit $\alpha \in Z_{k + 1}(X)$. Alors
$$
p : \coprod\nolimits_{j \in J} D_j \longrightarrow D
$$
est propre et
$$
i^*\alpha = p_*\left(\sum n_j i_j^*\alpha\right)
$$
dans $\CH_k(D)$.
\end{lemma}

\begin{proof}
La démonstration de ce lemme est un peu plus longue que prévu
en raison d'une subtilité concernant les sommes infinies d'équivalences rationnelles.
Dans le cas quasi-compact, la famille $D_j$ est finie et le résultat
est tout à fait élémentaire : il découle directement de
Homologie de Chow, Lemme \ref{chow-lemma-compute-c1}, de
Diviseurs, Lemme \ref{divisors-lemma-c1-additive}, et des définitions.

\medskip\noindent
Le morphisme $p$ est propre puisque la famille $\{D_j\}_{j \in J}$
est localement finie. Écrivons $\alpha = \sum_{a \in A} m_a [W_a]$
avec $W_a \subset X$ un sous-schéma intègre fermé de
$\delta$-dimension $k + 1$.
Notons $i_a : W_a \to X$ l'immersion fermée.
Nous supposons que $m_a \not = 0$ pour tout $a \in A$ et que
$\{W_a\}_{a \in A}$ est localement finie sur $X$.

\medskip\noindent
Observons que, d'après
Homologie de Chow, Définition \ref{chow-definition-gysin-homomorphism},
la classe $i^*\alpha$ est celle d'un cycle
$\sum m_a\beta_a$ pour certains $\beta_a \in Z_k(W_a \cap D)$.
Plus précisément, si $W_a \not \subset D$, alors $\beta_a = [D \cap W_a]_k$,
et si $W_a \subset D$, alors $\beta_a$ est un cycle
représentant $c_1(\mathcal{O}_X(D)) \cap [W_a]$.

\medskip\noindent
Pour chaque $a \in A$, écrivons $J = J_{a, 1} \amalg J_{a, 2} \amalg J_{a, 3}$,
où
\begin{enumerate}
\item $j \in J_{a, 1}$ si et seulement si $W_a \cap D_j = \emptyset$,
\item $j \in J_{a, 2}$ si et seulement si
$W_a \not = W_a \cap D_1 \not = \emptyset$, et
\item $j \in J_{a, 3}$ si et seulement si $W_a \subset D_j$.
\end{enumerate}
Puisque la famille $\{D_j\}$ est localement finie, nous voyons que
$J_{a, 3}$ est un ensemble fini. Pour tous $a \in A$ et $j \in J$,
nous choisissons un cycle $\beta_{a, j} \in Z_k(W_a \cap D_j)$ comme suit :
\begin{enumerate}
\item si $j \in J_{a, 1}$, nous posons $\beta_{a, j} = 0$,
\item si $j \in J_{a, 2}$, nous posons $\beta_{a, j} = [D_j \cap W_a]_k$, et
\item si $j \in J_{a, 3}$, nous choisissons $\beta_{a, j} \in Z_k(W_a)$
représentant $c_1(i_a^*\mathcal{O}_X(D_j)) \cap [W_j]$.
\end{enumerate}
Nous affirmons que
$$
\beta_a \sim_{rat}
\sum\nolimits_{j \in J} n_j \beta_{a, j}
$$
dans $\CH_k(W_a \cap D)$.

\medskip\noindent
Cas I : $W_a \not \subset D$. Dans ce cas, $J_{a, 3} = \emptyset$.
Il suffit donc de montrer que
$[D \cap W_a]_k = \sum n_j [D_j \cap W_a]_k$ comme cycles.
C'est le Lemme \ref{lemma-sum-divisors-associated-Weil}.

\medskip\noindent
Cas II : $W_a \subset D$. Dans ce cas, $\beta_a$ est un cycle représentant
$c_1(i_a^*\mathcal{O}_X(D)) \cap [W_a]$.
Écrivons $D = D_{a, 1} + D_{a, 2} + D_{a, 3}$ avec
$D_{a, s} = \sum_{j \in J_{a, s}} n_jD_j$.
D'après Diviseurs, Lemme \ref{divisors-lemma-c1-additive}, nous avons
\begin{eqnarray*}
c_1(i_a^*\mathcal{O}_X(D)) \cap [W_a] & = &
c_1(i_a^*\mathcal{O}_X(D_{a, 1})) \cap [W_a] +
c_1(i_a^*\mathcal{O}_X(D_{a, 2})) \cap [W_a] \\
& &
 + c_1(i_a^*\mathcal{O}_X(D_{a, 3})) \cap [W_a].
\end{eqnarray*}
Il est clair que le premier terme de la somme est nul.
Comme $J_{a, 3}$ est fini, nous voyons que le dernier terme coïncide
avec $\sum\nolimits_{j \in J_{a, 3}} n_jc_1(i_a^*\mathcal{L}_j) \cap [W_a]$ ;
voir Diviseurs, Lemme \ref{divisors-lemma-c1-additive}.
Ce cycle est représenté par $\sum_{j \in J_{a, 3}} n_j \beta_{a, j}$.
Enfin, le Cas I montre que le terme intermédiaire est représenté par le cycle
$\sum\nolimits_{j \in J_{a, 2}} n_j[D_j \cap W_a]_k =
\sum_{j \in J_{a, 2}} n_j\beta_{a, j}$.
D'où l'affirmation dans ce cas.

\medskip\noindent
Nous sommes maintenant en mesure d'achever la démonstration du lemme.
En effet, nous avons $i^*D \sim_{rat} \sum m_a\beta_a$ par notre
choix de $\beta_a$. Pour chaque $a$, nous avons
$\beta_a \sim_{rat} \sum_j \beta_{a, j}$, l'équivalence rationnelle
ayant lieu sur $D \cap W_a$.
Comme la famille de sous-schémas fermés $D \cap W_a$
est localement finie sur $D$, nous voyons aussi que
$\sum m_a \beta_a \sim_{rat} \sum_{a, j} m_a\beta_{a, j}$
sur $D$ (voir
Homologie de Chow, Remarque \ref{chow-remark-infinite-sums-rational-equivalences}).
Enfin, il est clair que $\sum_a m_a\beta_{a, j}$ (considéré
comme un cycle sur $D_j$) représente $i_j^*\alpha$, et donc
$\sum_{a, j} m_a\beta_{a, j}$ représente $p_* \sum_j i_j^*\alpha$,
ce qui achève la démonstration.
\end{proof}

\begin{lemma}
\label{lemma-commutativity-effective-Cartier-proper-intersection}
Soit $(S, \delta)$ comme dans Homologie de Chow, Situation \ref{chow-situation-setup}.
Soit $X$ localement de type fini sur $S$.
Supposons $X$ intègre et $\dim_\delta(X) = n$.
Soient $D$ et $D'$ des diviseurs de Cartier effectifs sur $X$.
Supposons $\dim_\delta(D \cap D') = n - 2$. Soient $i : D \to X$,
resp.\ $i' : D' \to X$, les immersions fermées correspondantes.
Alors
\begin{enumerate}
\item il existe un cycle $\alpha \in Z_{n - 2}(D \cap D')$
dont l'image directe sur $D$ représente
$i^*[D']_{n - 1} \in \CH_{n - 2}(D)$
et dont l'image directe sur $D'$ représente
$(i')^*[D]_{n - 1} \in \CH_{n - 2}(D')$, et
\item nous avons
$$
D \cdot [D']_{n - 1}
=
D' \cdot [D]_{n - 1}
$$
dans $\CH_{n - 2}(X)$.
\end{enumerate}
\end{lemma}

\begin{proof}
La partie (2) est une conséquence immédiate de la partie (1).
Écrivons $[D]_{n - 1} = \sum n_a[Z_a]$ et
$[D']_{n - 1} = \sum m_b[Z_b]$, où les $Z_a$ sont les composantes irréductibles
de $D$ et les $[Z_b]$ les composantes irréductibles
de $D'$. D'après
Homologie de Chow, Définition \ref{chow-definition-gysin-homomorphism},
nous avons $i^*D' = \sum m_b i^*[Z_b]$ et $(i')^*D = \sum n_a(i')^*[Z_a]$.
Par hypothèse, aucune des composantes irréductibles $Z_b$
n'est contenue dans $D$, et donc $i^*[Z_b] = [Z_b\cap D]_{n - 2}$
par définition. De même, $(i')^*[Z_a] = [Z_a \cap D']_{n - 2}$.
Nous cherchons donc à démontrer l'égalité de cycles
$$
\sum n_a[Z_a \cap D']_{n - 2} = \sum m_b[Z_b \cap D]_{n - 2}
$$
qui sont bien supportés sur $D \cap D'$.
Soit $W \subset X$ un sous-schéma intègre fermé
tel que $\dim_\delta(W) = n - 2$. Soit $\xi \in W$ son point générique.
Posons $R = \mathcal{O}_{X, \xi}$. C'est un anneau local noethérien intègre.
Notons que $\dim(R) = 2$. Soit $f \in R$, resp.\ $f' \in R$,
un élément définissant l'idéal de $D$, resp.\ de $D'$.
Par hypothèse, $\dim(R/(f, f')) = 0$. Soient
$\mathfrak q'_1, \ldots, \mathfrak q'_t \subset R$ les idéaux premiers minimaux
contenant $(f')$, et $\mathfrak q_1, \ldots, \mathfrak q_s \subset R$ les idéaux premiers minimaux
contenant $(f)$.
L'égalité ci-dessus se ramène à l'égalité
$$
\sum_{i = 1, \ldots, s}
\text{length}_{R_{\mathfrak q_i}}(R_{\mathfrak q_i}/(f))
\text{ord}_{R/\mathfrak q_i}(f')
=
\sum_{j = 1, \ldots, t}
\text{length}_{R_{\mathfrak q'_j}}(R_{\mathfrak q'_j}/(f'))
\text{ord}_{R/\mathfrak q'_j}(f).
$$
D'après Homologie de Chow, Lemme \ref{chow-lemma-length-multiplication}, 
appliqué avec $M = R/(f)$, le membre de gauche de
cette égalité est égal à
$$
\text{length}_R(R/(f, f'))
-
\text{length}_R(\Ker(f' : R/(f) \to R/(f)))
$$
Or nous observons que
$\Ker(f' : R/(f) \to R/(f))$ est canoniquement isomorphe
à $((f) \cap (f'))/(ff')$ par l'application $x \bmod (f) \mapsto
f'x \bmod (ff')$. Le membre de gauche vaut donc
$$
\text{length}_R(R/(f, f'))
-
\text{length}_R((f) \cap (f')/(ff'))
$$
Comme cette expression est symétrique en $f$ et $f'$, la démonstration est achevée.
\end{proof}

\begin{lemma}
\label{lemma-commutativity-effective-Cartier-proper-intersection-infinite}
Soit $(S, \delta)$ comme dans Homologie de Chow, Situation \ref{chow-situation-setup}.
Soit $X$ localement de type fini sur $S$.
Supposons $X$ intègre et $\dim_\delta(X) = n$.
Soit $\{D_j\}_{j \in J}$ une famille localement finie de
diviseurs de Cartier effectifs sur $X$. Soient $n_j, m_j \geq 0$ des
familles d'entiers positifs ou nuls. Posons
$D = \sum n_j D_j$ et $D' = \sum m_j D_j$.
Supposons que $\dim_\delta(D_j \cap D_{j'}) = n - 2$ pour tous
$j \not = j'$. Alors $D \cdot [D']_{n - 1} = D' \cdot [D]_{n - 1}$ dans
$\CH_{n - 2}(X)$.
\end{lemma}

\begin{proof}
Ce lemme découle immédiatement des
Lemmes \ref{lemma-sum-divisors-associated-Weil} et
\ref{lemma-commutativity-effective-Cartier-proper-intersection}
lorsque les sommes sont finies, par exemple si $X$ est quasi-compact.
Le lecteur pourra donc omettre la démonstration.

\medskip\noindent
Voici la démonstration dans le cas général.
Soient $i_j : D_j \to X$ les immersions fermées
Soit $p : \coprod D_j \to X$ le coproduit des morphismes $i_j$.
Soit $\{Z_a\}_{a \in A}$ la famille des composantes irréductibles de
$\bigcup D_j$. Pour chaque $j$, écrivons
$$
[D_j]_{n - 1} = \sum d_{j, a}[Z_a].
$$
D'après le Lemme \ref{lemma-sum-divisors-associated-Weil}, nous avons
$$
[D]_{n - 1} = \sum n_j d_{j, a} [Z_a],
\quad
[D']_{n - 1} = \sum m_j d_{j, a} [Z_a].
$$
D'après le Lemme \ref{lemma-improved-additivity},
nous avons
$$
D \cdot [D']_{n - 1} = p_*\left(\sum n_j i_j^*[D']_{n - 1} \right),
\quad
D' \cdot [D]_{n - 1} = p_*\left(\sum m_{j'} i_{j'}^*[D]_{n - 1} \right).
$$
Comme dans la définition des homomorphismes de Gysin (voir
Homologie de Chow, Définition \ref{chow-definition-gysin-homomorphism}),
nous choisissons des cycles $\beta_{a, j}$ sur $D_j \cap Z_a$ représentant
$i_j^*[Z_a]$. (Notons qu'en fait $\beta_{a, j} = [D_j \cap Z_a]_{n - 2}$
si $Z_a$ n'est pas contenu dans $D_j$, c'est-à-dire qu'il n'y a aucun choix dans ce cas.)
Puisque la restriction de $p$ à chacun des $D_j$ est une immersion fermée,
nous pouvons (et nous allons) considérer $\beta_{a, j}$ comme un cycle sur $X$.
En substituant les formules pour $[D]_{n - 1}$ et $[D']_{n - 1}$ obtenues
ci-dessus, on voit que
$$
D \cdot [D']_{n - 1} =
\sum\nolimits_{j, j', a} n_j m_{j'} d_{j', a} \beta_{a, j},
\quad
D' \cdot [D]_{n - 1} =
\sum\nolimits_{j, j', a} m_{j'} n_j d_{j, a} \beta_{a, j'}.
$$
De plus, avec les mêmes conventions, on a aussi
$$
D_j \cdot [D_{j'}]_{n - 1} = \sum d_{j', a} \beta_{a, j}.
$$
Avec ces notations, le
lemme \ref{lemma-commutativity-effective-Cartier-proper-intersection}
(voir aussi sa démonstration)
affirme que, pour $j \not = j'$, les cycles
$\sum d_{j', a} \beta_{a, j}$ et $\sum d_{j, a} \beta_{a, j'}$
sont égaux en tant que cycles ! On en déduit que
\begin{eqnarray*}
D \cdot [D']_{n - 1}
& = &
\sum\nolimits_{j, j', a} n_j m_{j'} d_{j', a} \beta_{a, j} \\
& = &
\sum\nolimits_{j \not = j'} n_j m_{j'}
\left(\sum\nolimits_a d_{j', a} \beta_{a, j}\right) +
\sum\nolimits_{j, a} n_j m_j d_{j, a} \beta_{a, j} \\
& = &
\sum\nolimits_{j \not = j'} n_j m_{j'}
\left(\sum\nolimits_a d_{j, a} \beta_{a, j'}\right) +
\sum\nolimits_{j, a} n_j m_j d_{j, a} \beta_{a, j} \\
& = &
\sum\nolimits_{j, j', a} m_{j'} n_j d_{j, a} \beta_{a, j'} \\
& = &
D' \cdot [D]_{n - 1}
\end{eqnarray*}
ce qui achève la démonstration.
\end{proof}

\begin{lemma}
\label{lemma-commutativity-effective-Cartier}
Soit $(S, \delta)$ comme dans Homologie de Chow, situation \ref{chow-situation-setup}.
Soit $X$ localement de type fini sur $S$.
Supposons que $X$ soit intègre et que $\dim_\delta(X) = n$.
Soient $D$ et $D'$ des diviseurs de Cartier effectifs sur $X$.
Alors
$$
D \cdot [D']_{n - 1} = D' \cdot [D]_{n - 1}
$$
dans $\CH_{n - 2}(X)$.
\end{lemma}

\begin{proof}[Première démonstration du lemme \ref{lemma-commutativity-effective-Cartier}]
Montrons d'abord l'assertion lorsque $X$ est quasi-compact.
Dans ce cas, appliquons
le lemme \ref{lemma-blowing-up-intersections} à $X$ et à
l'ensemble à deux éléments $\{D, D'\}$ de diviseurs de Cartier effectifs.
On obtient ainsi un morphisme propre $b : X' \to X$,
une famille finie de diviseurs de Cartier
effectifs $D'_j \subset X'$ dont les intersections deux à deux sont de codimension $\geq 2$,
avec $b^{-1}(D) = \sum n_j D'_j$, et $b^{-1}(D') = \sum m_j D'_j$.
Notons que $b_*[b^{-1}(D)]_{n - 1} = [D]_{n - 1}$ dans $Z_{n - 1}(X)$
et de même pour $D'$,
voir le lemme \ref{lemma-push-pull-effective-Cartier}.
Par conséquent, d'après Homologie de Chow, lemme \ref{chow-lemma-pushforward-cap-c1}, on a
$$
D \cdot [D']_{n - 1} = b_*\left(b^{-1}(D) \cdot [b^{-1}(D')]_{n - 1}\right)
$$
dans $\CH_{n - 2}(X)$ et de même pour l'autre terme. Le
lemme résulte donc de l'égalité
$b^{-1}(D) \cdot [b^{-1}(D')]_{n - 1} = b^{-1}(D') \cdot [b^{-1}(D)]_{n - 1}$
dans $\CH_{n - 2}(X')$ donnée par le lemme
\ref{lemma-commutativity-effective-Cartier-proper-intersection-infinite}.

\medskip\noindent
Notons que, dans la démonstration ci-dessus, chacun des lemmes cités s'applique aussi
au cas général (où $X$ n'est pas supposé quasi-compact). La
seule modification mineure dans le cas général est que le morphisme
$b : U' \to U$ obtenu en appliquant
le lemme \ref{lemma-blowing-up-intersections}
a pour but
un ouvert $U \subset X$ dont le complémentaire est de codimension $\geq 3$.
D'après Homologie de Chow, lemme \ref{chow-lemma-restrict-to-open}, on voit donc que
$\CH_{n - 2}(U) = \CH_{n - 2}(X)$
et, après avoir remplacé $X$ par $U$, le reste de la démonstration s'applique
sans changement.
\end{proof}

\begin{proof}[Seconde démonstration du lemme \ref{lemma-commutativity-effective-Cartier}]
Soient $\mathcal{I} = \mathcal{O}_X(-D)$ et
$\mathcal{I}' = \mathcal{O}_X(-D')$ les faisceaux
d'idéaux inversibles de $D$ et de $D'$.
Posons
$\mathcal{I}_{D'} = \mathcal{I} \otimes_{\mathcal{O}_X} \mathcal{O}_{D'}$
et
$\mathcal{I}'_D = \mathcal{I}' \otimes_{\mathcal{O}_X} \mathcal{O}_D$.
On peut restreindre le morphisme d'inclusion $\mathcal{I} \to \mathcal{O}_X$
à $D'$ pour obtenir un morphisme
$$
\varphi : \mathcal{I}_{D'} \longrightarrow \mathcal{O}_{D'}
$$
et, de même,
$$
\psi : \mathcal{I}'_D \longrightarrow \mathcal{O}_D
$$
Il est clair que
$$
\Coker(\varphi)
\cong
\mathcal{O}_{D \cap D'}
\cong
\Coker(\psi)
$$
et que
$$
\Ker(\varphi)
\cong
\frac{\mathcal{I} \cap \mathcal{I}'}{\mathcal{I}\mathcal{I}'}
\cong
\Ker(\psi).
$$
On voit donc que
$$
\gamma =
[\mathcal{I}_{D'}] - [\mathcal{O}_{D'}]
=
[\mathcal{I}'_D] - [\mathcal{O}_D]
$$
dans $K_0(\textit{Coh}_{\leq n - 1}(X))$. D'autre part, il est clair que
$$
[\mathcal{I}'_D]_{n - 1} = [D]_{n - 1}, \quad
[\mathcal{I}_{D'}]_{n - 1} = [D']_{n - 1}.
$$
On a également
$$
\mathcal{O}_X(D') \otimes \mathcal{I}'_D = \mathcal{O}_D, \quad
\mathcal{O}_X(D) \otimes \mathcal{I}_{D'} = \mathcal{O}_{D'}.
$$
D'après Homologie de Chow, lemme \ref{chow-lemma-coherent-sheaf-cap-c1}
(appliqué deux fois),
cela signifie que l'élément $\gamma$ appartient à $B_{n - 2}(X)$ et
a pour image à la fois $c_1(\mathcal{O}_X(D')) \cap [D]_{n - 1}$ et
$c_1(\mathcal{O}_X(D)) \cap [D']_{n - 1}$, ce qui achève la démonstration (puisque
l'application $B_{n - 2}(X) \to \CH_{n - 2}(X)$ est bien définie -- ce qui constitue
le point clef de cette démonstration).
\end{proof}
      
 
    
     
  
   
 




\section{Modules dualisants sur les modèles propres réguliers}
\label{section-dualizing}

\noindent
Dans Réduction semi-stable, situation
\ref{models-situation-regular-model}, on pose
$\omega_{X/R}^\bullet = f^!\mathcal{O}_{\Spec(R)}$,
le complexe dualisant relatif de $f : X \to \Spec(R)$ introduit dans
Dualité pour les schémas, remarque
\ref{duality-remark-relative-dualizing-complex}.
Comme $f$ est de Gorenstein de dimension relative $1$, d'après
Réduction semi-stable, lemme \ref{models-lemma-gorenstein}, les lemmes
\ref{duality-lemma-affine-flat-Noetherian-gorenstein},
\ref{duality-lemma-CM-shriek} et
\ref{duality-lemma-gorenstein-CM-morphism}
de Dualité pour les schémas donnent
$$
\omega_{X/R}^\bullet = \omega_X[1]
$$
pour un certain $\mathcal{O}_X$-module inversible $\omega_X$.
Ce module inversible est souvent appelé le
{\it module dualisant relatif de $X$ sur $R$}.
Comme $R$ est régulier (donc de Gorenstein) de dimension $1$, le complexe
$\omega_R^\bullet = R[1]$ est un complexe dualisant normalisé pour $R$.
Ainsi, $\omega_X = H^{-2}(f^!\omega_R^\bullet)$ et $\omega_X$ n'est pas
seulement un module dualisant relatif, mais aussi un module dualisant; voir
Dualité pour les schémas, exemple
\ref{duality-example-proper-over-local}.
Le module $\omega_X$ représente donc le foncteur
$$
\textit{Coh}(\mathcal{O}_X) \to \textit{Ens},\quad
\mathcal{F} \mapsto \Hom_R(H^1(X, \mathcal{F}), R)
$$
d'après Dualité pour les schémas, lemme
\ref{duality-lemma-dualizing-module-proper-over-A}.
On obtient ainsi une autre définition du module dualisant relatif dans
Réduction semi-stable, situation \ref{models-situation-regular-model}.
La formation de $\omega_X$ commute à tout changement de base
(pour tout morphisme propre de Gorenstein de dimension relative donnée);
cela résulte du fait correspondant pour le complexe dualisant relatif,
exposé dans Dualité pour les schémas, remarque
\ref{duality-remark-relative-dualizing-complex},
qui remonte à Dualité pour les schémas, lemme
\ref{duality-lemma-proper-flat-base-change}.
L'image réciproque de $\omega_X$ est donc le module dualisant $\omega_C$ de
$C$ sur $K$ étudié dans Courbes algébriques, lemme
\ref{curves-lemma-duality-dim-1-CM}.
Notons que $\omega_C$ est isomorphe à $\Omega_{C/K}$ d'après
Courbes algébriques, lemme \ref{curves-lemma-duality-dim-1}.
De même, $\omega_X|_{X_k}$ est le module dualisant $\omega_{X_k}$ de $X_k$
sur $k$.

\begin{lemma}
\label{lemma-dualizing-components}
Dans Réduction semi-stable, situation
\ref{models-situation-regular-model}, le module dualisant de $C_i$ sur $k$
est
$$
\omega_{C_i} = \omega_X(C_i)|_{C_i}
$$
où $\omega_X$ est défini ci-dessus.
\end{lemma}

\begin{proof}
Soit $t : C_i \to X$ l'immersion fermée. Comme $t$ est l'inclusion d'un
diviseur de Cartier effectif, les lemmes
\ref{duality-lemma-twisted-inverse-image-closed} et
\ref{duality-lemma-sheaf-with-exact-support-effective-Cartier}
de Dualité pour les schémas donnent
$t^!(\mathcal{L}) = \mathcal{L}(C_i)|_{C_i}$ pour tout
$\mathcal{O}_X$-module inversible $\mathcal{L}$.
Considérons le diagramme commutatif
$$
\xymatrix{
C_i \ar[r]_t \ar[d]_g & X \ar[d]^f \\
\Spec(k) \ar[r]^s & \Spec(R)
}
$$
La courbe $C_i$ est de Gorenstein (Réduction semi-stable, lemme
\ref{models-lemma-gorenstein}) et son module dualisant inversible
$\omega_{C_i}$ est caractérisé par la propriété
$\omega_{C_i}[0] = g^!\mathcal{O}_{\Spec(k)}$. Voir Courbes algébriques,
lemme \ref{curves-lemma-duality-dim-1}, sa démonstration, ainsi que les lemmes
\ref{curves-lemma-duality-dim-1-CM} et \ref{curves-lemma-rr}
de Courbes algébriques.
D'autre part, $s^!(R[1]) = k$, donc
$$
\omega_{C_i}[0] =
g^! s^!(R[1]) = t^!f^!(R[1]) = t^!\omega_X
$$
En combinant ces observations, on obtient l'énoncé du lemme.
\end{proof}






\section{Renvois en double ou scindés}
\label{section-duplicates}

\noindent
Cette section rassemble les doublons ainsi que les renvois qui énonçaient
autrefois plusieurs faits simultanément, mais qui sont désormais scindés en
leurs différentes composantes.

\begin{lemma}
\label{lemma-characterize-epi-mono}
Soit $X$ un espace topologique.
\begin{enumerate}
\item Les épimorphismes (resp.\ les monomorphismes) dans la catégorie des
préfaisceaux sont exactement les morphismes surjectifs (resp.\ injectifs) de
préfaisceaux.
\item Les épimorphismes (resp.\ les monomorphismes) dans la catégorie des
faisceaux sont exactement les morphismes surjectifs (resp.\ injectifs) de
faisceaux, et exactement ceux qui sont surjectifs (resp.\ injectifs) sur
toutes les fibres.
\item Le morphisme de faisceaux associé à un morphisme surjectif
(resp.\ injectif) de préfaisceaux d'ensembles est surjectif
(resp.\ injectif).
\end{enumerate}
\end{lemma}

\begin{proof}
On combine les lemmes \ref{sheaves-lemma-characterize-epi-mono} et
\ref{sheaves-lemma-sheafify-epi-mono} de Faisceaux.
\end{proof}

\begin{lemma}
\label{lemma-directed-colimit-finite-type}
Soit $X$ un schéma quasi-compact et quasi-séparé.
Soit $\mathcal{F}$ un $\mathcal{O}_X$-module quasi-cohérent.
Alors $\mathcal{F}$ est la limite inductive filtrante de ses sous-modules
quasi-cohérents de type fini.
\end{lemma}

\begin{proof}
C'est un doublon de Propriétés, lemme
\ref{properties-lemma-quasi-coherent-colimit-finite-type}.
\end{proof}

\begin{lemma}
\label{lemma-points-monomorphism}
Soit $S$ un schéma et soit $X$ un espace algébrique sur $S$.
L'application $\{\Spec(k) \to X \text{ monomorphisme}\} \to |X|$ est injective.
\end{lemma}

\begin{proof}
C'est un doublon de Propriétés des espaces, lemme
\ref{spaces-properties-lemma-points-monomorphism}.
\end{proof}

\begin{theorem}
\label{theorem-equivalence-sheaves-point}
Soit $S = \Spec(K)$, où $K$ est un corps.
Soit $\overline{s}$ un point géométrique de $S$.
Notons $G = \text{Gal}_{\kappa(s)}$ le groupe de Galois absolu.
Il existe alors une équivalence de catégories
$\Sh(S_\etale) \to G\textit{-Ens}$,
$\mathcal{F} \mapsto \mathcal{F}_{\overline{s}}$.
\end{theorem}

\begin{proof}
C'est un doublon de Cohomologie \'etale, théorème
\ref{etale-cohomology-theorem-equivalence-sheaves-point}.
\end{proof}

\begin{remark}
\label{remark-tangent-spaces}
Ce renvoi mène ici en raison d'une étiquette dupliquée. Pour le contenu
proprement dit, voir Théorie formelle des déformations, section
\ref{formal-defos-section-tangent-spaces}.
\end{remark}

\begin{lemma}
\label{lemma-locally-ringed-space-direct-summand-free}
Soit $X$ un espace localement annelé. Tout facteur direct d'un
$\mathcal{O}_X$-module libre de type fini est localement libre de type fini.
\end{lemma}

\begin{proof}
C'est un doublon de Modules, lemme
\ref{modules-lemma-direct-summand-of-locally-free-is-locally-free}.
\end{proof}

\begin{lemma}
\label{lemma-characterize-injective}
Soit $R$ un anneau et soit $E$ un $R$-module. Les conditions suivantes sont
équivalentes :
\begin{enumerate}
\item $E$ est un $R$-module injectif;
\item pour tout idéal $I \subset R$ et tout morphisme de modules
$\varphi : I \to E$, il existe un prolongement de $\varphi$ en un morphisme
de $R$-modules $R \to E$.
\end{enumerate}
\end{lemma}

\begin{proof}
C'est le critère de Baer; voir Injectifs, lemme
\ref{injectives-lemma-criterion-baer}.
\end{proof}

\begin{lemma}
\label{lemma-periodic-length}
Soit $R$ un anneau local.
\begin{enumerate}
\item Si $(M, N, \varphi, \psi)$ est un complexe $2$-périodique tel que
$M$ et $N$ soient de longueur finie, alors
$e_R(M, N, \varphi, \psi) = \text{longueur}_R(M) - \text{longueur}_R(N)$.
\item Si $(M, \varphi, \psi)$ est un complexe $(2, 1)$-périodique tel que
$M$ soit de longueur finie, alors
$e_R(M, \varphi, \psi) = 0$.
\item Supposons donnée une suite exacte de complexes $2$-périodiques
$$
0 \to (M_1, N_1, \varphi_1, \psi_1)
\to (M_2, N_2, \varphi_2, \psi_2)
\to (M_3, N_3, \varphi_3, \psi_3)
\to 0
$$
Si deux de ces trois complexes ont des modules de cohomologie de longueur
finie, il en va de même du troisième, et l'on a
$$
e_R(M_2, N_2, \varphi_2, \psi_2) =
e_R(M_1, N_1, \varphi_1, \psi_1) +
e_R(M_3, N_3, \varphi_3, \psi_3).
$$
\end{enumerate}
\end{lemma}

\begin{proof}
Cela résulte des lemmes \ref{chow-lemma-additivity-periodic-length} et
\ref{chow-lemma-finite-periodic-length} d'Homologie de Chow.
\end{proof}




\begin{lemma}
\label{lemma-extensions-of-rings}
Soient $A$ un anneau et $I$ un $A$-module.
\begin{enumerate}
\item L'ensemble des extensions d'anneaux $0 \to I \to A' \to A \to 0$
où $I$ est un idéal de carré nul est en bijection canonique avec
$\Ext^1_A(\NL_{A/\mathbf{Z}}, I)$.
\item Étant donnés un homomorphisme d'anneaux $A \to B$, un $B$-module $N$,
un homomorphisme de $A$-modules $c : I \to N$, ainsi que les extensions
d'anneaux suivantes, à noyaux de carré nul :
\begin{enumerate}
\item[(a)] $0 \to I \to A' \to A \to 0$ correspondant à
$\alpha \in \Ext^1_A(\NL_{A/\mathbf{Z}}, I)$, et
\item[(b)] $0 \to N \to B' \to B \to 0$ correspondant à
$\beta \in \Ext^1_B(\NL_{B/\mathbf{Z}}, N)$,
\end{enumerate}
alors il existe un homomorphisme $A' \to B'$ s'inscrivant dans
le diagramme de Théorie des déformations, équation (\ref{defos-equation-to-solve}),
si et seulement si $\beta$ et $\alpha$
ont la même image dans
$\Ext^1_A(\NL_{A/\mathbf{Z}}, N)$.
\end{enumerate}
\end{lemma}

\begin{proof}
Cela résulte de Théorie des déformations, lemmes
\ref{defos-lemma-extensions-of-algebras} et
\ref{defos-lemma-extensions-of-algebras-functorial}.
\end{proof}

\begin{lemma}
\label{lemma-extensions-of-ringed-spaces}
Soient $(S, \mathcal{O}_S)$ un espace annelé et $\mathcal{J}$
un $\mathcal{O}_S$-module.
\begin{enumerate}
\item L'ensemble des extensions de faisceaux d'anneaux
$0 \to \mathcal{J} \to \mathcal{O}_{S'} \to \mathcal{O}_S \to 0$
où $\mathcal{J}$ est un idéal de carré nul est en bijection canonique avec
$\Ext^1_{\mathcal{O}_S}(\NL_{S/\mathbf{Z}}, \mathcal{J})$.
\item Étant donnés un morphisme d'espaces annelés
$f : (X, \mathcal{O}_X) \to (S, \mathcal{O}_S)$, un $\mathcal{O}_X$-module
$\mathcal{G}$, un $f$-morphisme $c : \mathcal{J} \to \mathcal{G}$, ainsi que
les extensions suivantes de faisceaux d'anneaux, à noyaux de carré nul :
\begin{enumerate}
\item[(a)] $0 \to \mathcal{J} \to \mathcal{O}_{S'} \to \mathcal{O}_S \to 0$
correspondant à
$\alpha \in \Ext^1_{\mathcal{O}_S}(\NL_{S/\mathbf{Z}}, \mathcal{J})$,
\item[(b)] $0 \to \mathcal{G} \to \mathcal{O}_{X'} \to \mathcal{O}_X \to 0$
correspondant à
$\beta \in \Ext^1_{\mathcal{O}_X}(\NL_{X/\mathbf{Z}}, \mathcal{G})$,
\end{enumerate}
alors il existe un morphisme $X' \to S'$ s'inscrivant dans
le diagramme de Théorie des déformations, équation
(\ref{defos-equation-to-solve-ringed-spaces}),
si et seulement si $\beta$ et $\alpha$
ont la même image dans
$\Ext^1_{\mathcal{O}_X}(Lf^*\NL_{S/\mathbf{Z}}, \mathcal{G})$.
\end{enumerate}
\end{lemma}

\begin{proof}
Cela résulte de Théorie des déformations, lemmes
\ref{defos-lemma-extensions-of-relative-ringed-spaces} et
\ref{defos-lemma-extensions-of-relative-ringed-spaces-functorial}.
\end{proof}

\begin{lemma}
\label{lemma-extensions-of-ringed-topoi}
Soient $(\Sh(\mathcal{B}), \mathcal{O}_\mathcal{B})$ un topos annelé
et $\mathcal{J}$ un $\mathcal{O}_\mathcal{B}$-module.
\begin{enumerate}
\item L'ensemble des extensions de faisceaux d'anneaux
$0 \to \mathcal{J} \to \mathcal{O}_{\mathcal{B}'} \to
\mathcal{O}_\mathcal{B} \to 0$
où $\mathcal{J}$ est un idéal de carré nul est en bijection canonique avec
$\Ext^1_{\mathcal{O}_\mathcal{B}}(
\NL_{\mathcal{O}_\mathcal{B}/\mathbf{Z}}, \mathcal{J})$.
\item Étant donnés un morphisme de topos annelés
$f : (\Sh(\mathcal{C}), \mathcal{O}) \to
(\Sh(\mathcal{B}), \mathcal{O}_\mathcal{B})$, un $\mathcal{O}$-module
$\mathcal{G}$, un homomorphisme de $f^{-1}\mathcal{O}_\mathcal{B}$-modules
$c : f^{-1}\mathcal{J} \to \mathcal{G}$, ainsi que les extensions suivantes
de faisceaux d'anneaux, à noyaux de carré nul :
\begin{enumerate}
\item[(a)] $0 \to \mathcal{J} \to \mathcal{O}_{\mathcal{B}'} \to
\mathcal{O}_\mathcal{B} \to 0$ correspondant à
$\alpha \in \Ext^1_{\mathcal{O}_\mathcal{B}}(
\NL_{\mathcal{O}_\mathcal{B}/\mathbf{Z}}, \mathcal{J})$,
\item[(b)] $0 \to \mathcal{G} \to \mathcal{O}' \to \mathcal{O} \to 0$
correspondant à
$\beta \in \Ext^1_\mathcal{O}(\NL_{\mathcal{O}/\mathbf{Z}}, \mathcal{G})$,
\end{enumerate}
alors il existe un morphisme $(\Sh(\mathcal{C}), \mathcal{O}') \to
(\Sh(\mathcal{B}), \mathcal{O}_{\mathcal{B}'})$ s'inscrivant dans
le diagramme de Théorie des déformations, équation (\ref{defos-equation-to-solve-ringed-topoi}),
si et seulement si $\beta$ et $\alpha$ ont la même image dans
$\Ext^1_\mathcal{O}(
Lf^*\NL_{\mathcal{O}_\mathcal{B}/\mathbf{Z}}, \mathcal{G})$.
\end{enumerate}
\end{lemma}

\begin{proof}
Cela résulte de Théorie des déformations, lemmes
\ref{defos-lemma-extensions-of-relative-ringed-topoi} et
\ref{defos-lemma-extensions-of-relative-ringed-topoi-functorial}.
\end{proof}

\begin{remark}
\label{remark-examples-formal-defos}
Cette étiquette renvoyait autrefois à une section décrivant plusieurs exemples
de problèmes de déformation. Ces exemples ont désormais chacun leur propre section.
Voir Problèmes de déformation, sections
\ref{examples-defos-section-finite-projective-modules},
\ref{examples-defos-section-representations},
\ref{examples-defos-section-continuous-representations}, et
\ref{examples-defos-section-graded-algebras}.
\end{remark}

\begin{lemma}
\label{lemma-examples-have-RS}
Les exemples de Problèmes de déformation
\ref{examples-defos-example-finite-projective-modules},
\ref{examples-defos-example-representations},
\ref{examples-defos-example-continuous-representations}, et
\ref{examples-defos-example-graded-algebras}
satisfont à la condition de Rim--Schlessinger (RS).
\end{lemma}

\begin{proof}
Cela résulte de Problèmes de déformation, lemmes
\ref{examples-defos-lemma-finite-projective-modules-RS},
\ref{examples-defos-lemma-representations-RS},
\ref{examples-defos-lemma-continuous-representations-RS}, et
\ref{examples-defos-lemma-graded-algebras-RS}.
\end{proof}

\begin{lemma}
\label{lemma-tangent-and-inf}
On a les identifications canoniques suivantes d'espaces vectoriels sur $k$ :
\begin{enumerate}
\item Dans Problèmes de déformation, exemple
\ref{examples-defos-example-finite-projective-modules},
si $x_0 = (k, V)$, alors $T_{x_0}\mathcal{F} = (0)$
et $\text{Inf}_{x_0}(\mathcal{F}) = \text{End}_k(V)$
sont de dimension finie.
\item Dans Problèmes de déformation, exemple
\ref{examples-defos-example-representations},
si $x_0 = (k, V, \rho_0)$, alors
$T_{x_0}\mathcal{F} = \Ext^1_{k[\Gamma]}(V, V) = H^1(\Gamma, \text{End}_k(V))$
et $\text{Inf}_{x_0}(\mathcal{F}) = H^0(\Gamma, \text{End}_k(V))$
sont de dimension finie si $\Gamma$ est engendré par un nombre fini d'éléments.
\item Dans Problèmes de déformation, exemple
\ref{examples-defos-example-continuous-representations},
si $x_0 = (k, V, \rho_0)$, alors
$T_{x_0}\mathcal{F} = H^1_{cont}(\Gamma, \text{End}_k(V))$
et
$\text{Inf}_{x_0}(\mathcal{F}) = H^0_{cont}(\Gamma, \text{End}_k(V))$
sont de dimension finie si $\Gamma$ est topologiquement de type fini.
\item Dans Problèmes de déformation, exemple
\ref{examples-defos-example-graded-algebras},
si $x_0 = (k, P)$, alors
$T_{x_0}\mathcal{F}$ et $\text{Inf}_{x_0}(\mathcal{F}) = \text{Der}_k(P, P)$
sont de dimension finie si $P$ est une $k$-algèbre de type fini.
\end{enumerate}
\end{lemma}

\begin{proof}
Cela résulte de Problèmes de déformation, lemmes
\ref{examples-defos-lemma-finite-projective-modules-TI},
\ref{examples-defos-lemma-representations-TI},
\ref{examples-defos-lemma-continuous-representations-TI}, et
\ref{examples-defos-lemma-graded-algebras-TI}.
\end{proof}






\begin{multicols}{2}[\section{Autres chapitres}]
\noindent
Préliminaires
\begin{enumerate}
\item \hyperref[introduction-section-phantom]{Introduction}
\item \hyperref[conventions-section-phantom]{Conventions}
\item \hyperref[sets-section-phantom]{Théorie des ensembles}
\item \hyperref[categories-section-phantom]{Catégories}
\item \hyperref[topology-section-phantom]{Topologie}
\item \hyperref[sheaves-section-phantom]{Faisceaux sur les espaces}
\item \hyperref[sites-section-phantom]{Sites et faisceaux}
\item \hyperref[stacks-section-phantom]{Champs}
\item \hyperref[fields-section-phantom]{Corps}
\item \hyperref[algebra-section-phantom]{Algèbre commutative}
\item \hyperref[brauer-section-phantom]{Groupes de Brauer}
\item \hyperref[homology-section-phantom]{Algèbre homologique}
\item \hyperref[derived-section-phantom]{Catégories dérivées}
\item \hyperref[simplicial-section-phantom]{Méthodes simpliciales}
\item \hyperref[more-algebra-section-phantom]{Compléments d'algèbre}
\item \hyperref[smoothing-section-phantom]{Lissage des morphismes d'anneaux}
\item \hyperref[modules-section-phantom]{Faisceaux de modules}
\item \hyperref[sites-modules-section-phantom]{Modules sur les sites}
\item \hyperref[injectives-section-phantom]{Objets injectifs}
\item \hyperref[cohomology-section-phantom]{Cohomologie des faisceaux}
\item \hyperref[sites-cohomology-section-phantom]{Cohomologie sur les sites}
\item \hyperref[dga-section-phantom]{Algèbre différentielle graduée}
\item \hyperref[dpa-section-phantom]{Algèbre à puissances divisées}
\item \hyperref[sdga-section-phantom]{Faisceaux différentiels gradués}
\item \hyperref[hypercovering-section-phantom]{Hyperrecouvrements}
\end{enumerate}
Schémas
\begin{enumerate}
\setcounter{enumi}{25}
\item \hyperref[schemes-section-phantom]{Schémas}
\item \hyperref[constructions-section-phantom]{Constructions de schémas}
\item \hyperref[properties-section-phantom]{Propriétés des schémas}
\item \hyperref[morphisms-section-phantom]{Morphismes de schémas}
\item \hyperref[coherent-section-phantom]{Cohomologie des schémas}
\item \hyperref[divisors-section-phantom]{Diviseurs}
\item \hyperref[limits-section-phantom]{Limites de schémas}
\item \hyperref[varieties-section-phantom]{Variétés}
\item \hyperref[topologies-section-phantom]{Topologies sur les schémas}
\item \hyperref[descent-section-phantom]{Descente}
\item \hyperref[perfect-section-phantom]{Catégories dérivées de schémas}
\item \hyperref[more-morphisms-section-phantom]{Compléments sur les morphismes}
\item \hyperref[flat-section-phantom]{Compléments sur la platitude}
\item \hyperref[groupoids-section-phantom]{Schémas en groupoïdes}
\item \hyperref[more-groupoids-section-phantom]{Compléments sur les schémas en groupoïdes}
\item \hyperref[etale-section-phantom]{Morphismes étales de schémas}
\end{enumerate}
Thèmes de théorie des schémas
\begin{enumerate}
\setcounter{enumi}{41}
\item \hyperref[chow-section-phantom]{Homologie de Chow}
\item \hyperref[intersection-section-phantom]{Théorie de l'intersection}
\item \hyperref[pic-section-phantom]{Schémas de Picard des courbes}
\item \hyperref[weil-section-phantom]{Théories cohomologiques de Weil}
\item \hyperref[adequate-section-phantom]{Modules adéquats}
\item \hyperref[dualizing-section-phantom]{Complexes dualisants}
\item \hyperref[duality-section-phantom]{Dualité pour les schémas}
\item \hyperref[discriminant-section-phantom]{Discriminants et différentes}
\item \hyperref[derham-section-phantom]{Cohomologie de de Rham}
\item \hyperref[local-cohomology-section-phantom]{Cohomologie locale}
\item \hyperref[algebraization-section-phantom]{Géométrie algébrique et formelle}
\item \hyperref[curves-section-phantom]{Courbes algébriques}
\item \hyperref[resolve-section-phantom]{Résolution des surfaces}
\item \hyperref[models-section-phantom]{Réduction semi-stable}
\item \hyperref[functors-section-phantom]{Foncteurs et morphismes}
\item \hyperref[equiv-section-phantom]{Catégories dérivées de variétés}
\item \hyperref[pione-section-phantom]{Groupes fondamentaux de schémas}
\item \hyperref[etale-cohomology-section-phantom]{Cohomologie étale}
\item \hyperref[crystalline-section-phantom]{Cohomologie cristalline}
\item \hyperref[proetale-section-phantom]{Cohomologie pro-étale}
\item \hyperref[relative-cycles-section-phantom]{Cycles relatifs}
\item \hyperref[more-etale-section-phantom]{Compléments de cohomologie étale}
\item \hyperref[trace-section-phantom]{La formule des traces}
\end{enumerate}
Espaces algébriques
\begin{enumerate}
\setcounter{enumi}{64}
\item \hyperref[spaces-section-phantom]{Espaces algébriques}
\item \hyperref[spaces-properties-section-phantom]{Propriétés des espaces algébriques}
\item \hyperref[spaces-morphisms-section-phantom]{Morphismes d'espaces algébriques}
\item \hyperref[decent-spaces-section-phantom]{Espaces algébriques décents}
\item \hyperref[spaces-cohomology-section-phantom]{Cohomologie des espaces algébriques}
\item \hyperref[spaces-limits-section-phantom]{Limites d'espaces algébriques}
\item \hyperref[spaces-divisors-section-phantom]{Diviseurs sur les espaces algébriques}
\item \hyperref[spaces-over-fields-section-phantom]{Espaces algébriques sur des corps}
\item \hyperref[spaces-topologies-section-phantom]{Topologies sur les espaces algébriques}
\item \hyperref[spaces-descent-section-phantom]{Descente et espaces algébriques}
\item \hyperref[spaces-perfect-section-phantom]{Catégories dérivées d'espaces}
\item \hyperref[spaces-more-morphisms-section-phantom]{Compléments sur les morphismes d'espaces}
\item \hyperref[spaces-flat-section-phantom]{Platitude sur les espaces algébriques}
\item \hyperref[spaces-groupoids-section-phantom]{Groupoïdes dans les espaces algébriques}
\item \hyperref[spaces-more-groupoids-section-phantom]{Compléments sur les groupoïdes dans les espaces}
\item \hyperref[bootstrap-section-phantom]{Amorçage}
\item \hyperref[spaces-pushouts-section-phantom]{Sommes amalgamées d'espaces algébriques}
\end{enumerate}
Thèmes de géométrie
\begin{enumerate}
\setcounter{enumi}{81}
\item \hyperref[spaces-chow-section-phantom]{Groupes de Chow des espaces}
\item \hyperref[groupoids-quotients-section-phantom]{Quotients de groupoïdes}
\item \hyperref[spaces-more-cohomology-section-phantom]{Compléments sur la cohomologie des espaces}
\item \hyperref[spaces-simplicial-section-phantom]{Espaces simpliciaux}
\item \hyperref[spaces-duality-section-phantom]{Dualité pour les espaces}
\item \hyperref[formal-spaces-section-phantom]{Espaces algébriques formels}
\item \hyperref[restricted-section-phantom]{Algébrisation des espaces formels}
\item \hyperref[spaces-resolve-section-phantom]{Résolution des surfaces revisitée}
\end{enumerate}
Théorie des déformations
\begin{enumerate}
\setcounter{enumi}{89}
\item \hyperref[formal-defos-section-phantom]{Théorie formelle des déformations}
\item \hyperref[defos-section-phantom]{Théorie des déformations}
\item \hyperref[cotangent-section-phantom]{Le complexe cotangent}
\item \hyperref[examples-defos-section-phantom]{Problèmes de déformation}
\end{enumerate}
Champs algébriques
\begin{enumerate}
\setcounter{enumi}{93}
\item \hyperref[algebraic-section-phantom]{Champs algébriques}
\item \hyperref[examples-stacks-section-phantom]{Exemples de champs}
\item \hyperref[stacks-sheaves-section-phantom]{Faisceaux sur les champs algébriques}
\item \hyperref[criteria-section-phantom]{Critères de représentabilité}
\item \hyperref[artin-section-phantom]{Axiomes d'Artin}
\item \hyperref[quot-section-phantom]{Espaces de Quot et de Hilbert}
\item \hyperref[stacks-properties-section-phantom]{Propriétés des champs algébriques}
\item \hyperref[stacks-morphisms-section-phantom]{Morphismes de champs algébriques}
\item \hyperref[stacks-limits-section-phantom]{Limites de champs algébriques}
\item \hyperref[stacks-cohomology-section-phantom]{Cohomologie des champs algébriques}
\item \hyperref[stacks-perfect-section-phantom]{Catégories dérivées de champs}
\item \hyperref[stacks-introduction-section-phantom]{Introduction aux champs algébriques}
\item \hyperref[stacks-more-morphisms-section-phantom]{Compléments sur les morphismes de champs}
\item \hyperref[stacks-geometry-section-phantom]{La géométrie des champs}
\end{enumerate}
Thèmes de théorie des espaces de modules
\begin{enumerate}
\setcounter{enumi}{107}
\item \hyperref[moduli-section-phantom]{Champs de modules}
\item \hyperref[moduli-curves-section-phantom]{Espaces de modules de courbes}
\end{enumerate}
Divers
\begin{enumerate}
\setcounter{enumi}{109}
\item \hyperref[examples-section-phantom]{Exemples}
\item \hyperref[exercises-section-phantom]{Exercices}
\item \hyperref[guide-section-phantom]{Guide bibliographique}
\item \hyperref[desirables-section-phantom]{Desiderata}
\item \hyperref[coding-section-phantom]{Style de programmation}
\item \hyperref[obsolete-section-phantom]{Obsolète}
\item \hyperref[fdl-section-phantom]{Licence de documentation libre GNU}
\item \hyperref[index-section-phantom]{Index généré automatiquement}
\end{enumerate}
\end{multicols}


\bibliography{my}
\bibliographystyle{amsalpha}

\end{document}
